\documentclass[11pt,a4paper]{report}
\usepackage{fontspec,amsmath,amssymb,amsthm,microtype,xcolor}
\setmainfont{lmroman10-regular.otf}[BoldFont=lmroman10-bold.otf,ItalicFont=lmroman10-italic.otf,BoldItalicFont=lmroman10-bolditalic.otf]
\setmonofont{lmmono10-regular.otf}[Scale=MatchLowercase]
\usepackage[margin=25mm]{geometry}
\usepackage[unicode,hidelinks]{hyperref}
\setlength{\emergencystretch}{3em}
\setlength{\parskip}{0.4em}
\setlength{\parindent}{0pt}
\setcounter{tocdepth}{1}
\makeatletter
\renewcommand*\l@section{\@dottedtocline{1}{1.5em}{3.8em}}
\renewcommand*\@pnumwidth{2.5em}
\renewcommand*\@tocrmarg{3.5em}
\makeatother
\setcounter{secnumdepth}{2}
\allowdisplaybreaks[2]
\raggedbottom
\title{\texorpdfstring{Unadjusted Langevin under Strong Convexity}{Unadjusted Langevin under Strong Convexity}\\[0.6em]\large Mathematical proofs}
\author{LangevinLean}
\date{}
\begin{document}
\pagenumbering{roman}
\maketitle
\chapter*{About this document}
\addcontentsline{toc}{chapter}{About this document}
This volume collects the theorem statements, mathematical arguments and supporting proofs for this development. Every locally authored proof module required by the displayed results, proof roadmap and linked extensions is included once, in dependency order. Foundational results from Mathlib are used as standard mathematical results.



The strongly convex convergence theorems construct their Brownian reference and prove Gibbs stationarity. The LSI extension derives its entropy bound from smoothness, LSI, and finite initial KL.



\textbf{Paper references}



The two step-size regimes and the coefficient 1.65 follow Theorem 1(a)–(b) of Dalalyan and Karagulyan (2019). Dalalyan (2017), Theorem 1 and Section 6, gives the earlier synchronous-coupling argument.


\begin{itemize}

\item Arnak S. Dalalyan and Avetik G. Karagulyan (2019). User-friendly guarantees for the Langevin Monte Carlo with inaccurate gradient. Stochastic Processes and their Applications 129(12), 5278–5311. arXiv
\item Arnak Dalalyan (2017). Further and stronger analogy between sampling and optimization: Langevin Monte Carlo and gradient descent. Proceedings of the 2017 Conference on Learning Theory, Proceedings of Machine Learning Research 65, 678–689.
\end{itemize}
\tableofcontents
\clearpage
\pagenumbering{arabic}
\chapter{Assumptions and main results}
\section{\texorpdfstring{Small-\allowbreak{}step Gaussian ULA}{Small-step Gaussian ULA}}


Let \(d\ge1\), and let \(U:\mathbb R^d\to\mathbb R\) be differentiable and \(m\)-strongly convex, with \(m>0\) and \(L\)-Lipschitz gradient. These assumptions imply \(m\le L\). Define


\[\pi(dx)=Z^{-1}e^{-U(x)}dx,\qquad Z=\int e^{-U(x)}dx\in(0,\infty).\]


Let \(\mu_0\) be a probability law with \(W_2(\mu_0,\pi)<\infty\). A finite second moment suffices. For \(h>0\), take independent standard Gaussian vectors \(\xi_k\), independent of initialization, and set


\[X_{k+1}=X_k-h\nabla U(X_k)+\sqrt{2h}\,\xi_k,\qquad\mu_k=\operatorname{Law}(X_k).\]


In the Lean statement, \(\mathbb R^d\) is \texttt{Euclidean\allowbreak{}Space \ensuremath{\mathbb R} \ensuremath{\iota}} with a finite nonempty index type, and \(d=|\iota|\). \texttt{ula\allowbreak{}Law\allowbreak{}Step} is the exact Gaussian transition on probability laws; its identification with this recursion is proved.



If \(0<h\le2/(m+L)\), then for every \(k\in\mathbb N\),

\[W_2(\mu_k,\pi)\le(1-mh)^kW_2(\mu_0,\pi)+\frac{33}{20}\frac Lm\sqrt{hd}.\]

The rational coefficient \(33/20\) is exactly \(1.65\). The initial term uses the true Wasserstein infimum over couplings.


{\small\color{gray}Supporting proof: Section~\ref{proof-module-UnadjustedLangevin-CanonicalULA-lean}.\par}
\section{\texorpdfstring{The other stable step-\allowbreak{}size regime}{The other stable step-size regime}}


Let \(d\ge1\), and let \(U:\mathbb R^d\to\mathbb R\) be differentiable and \(m\)-strongly convex, with \(m>0\) and \(L\)-Lipschitz gradient. These assumptions imply \(m\le L\). Define


\[\pi(dx)=Z^{-1}e^{-U(x)}dx,\qquad Z=\int e^{-U(x)}dx\in(0,\infty).\]


Let \(\mu_0\) be a probability law with \(W_2(\mu_0,\pi)<\infty\). A finite second moment suffices. For \(h>0\), take independent standard Gaussian vectors \(\xi_k\), independent of initialization, and set


\[X_{k+1}=X_k-h\nabla U(X_k)+\sqrt{2h}\,\xi_k,\qquad\mu_k=\operatorname{Law}(X_k).\]


In the Lean statement, \(\mathbb R^d\) is \texttt{Euclidean\allowbreak{}Space \ensuremath{\mathbb R} \ensuremath{\iota}} with a finite nonempty index type, and \(d=|\iota|\). \texttt{ula\allowbreak{}Law\allowbreak{}Step} is the exact Gaussian transition on probability laws; its identification with this recursion is proved.



If \(2/(m+L)\le h\) and \(Lh<2\), then for every \(k\in\mathbb N\),

\[W_2(\mu_k,\pi)\le(Lh-1)^kW_2(\mu_0,\pi)+\frac{33}{20}\frac{Lh}{2-Lh}\sqrt{hd}.\]

The strict upper endpoint ensures the denominator is positive. Both regimes follow from the same scalar recurrence, using their respective contraction factors.


{\small\color{gray}Supporting proof: Section~\ref{proof-module-UnadjustedLangevin-CanonicalULA-lean}.\par}
\section{Proof organization}
The supporting proofs below establish the following steps.
\begin{enumerate}
\item \textbf{Contract the gradient step}. \(q=1-mh\) or \(Lh-1\) See Section~\ref{proof-module-UnadjustedLangevin-SharpContraction-lean}.
\item \textbf{Construct the Gibbs target}. Positive normalization and \(L^2\) moments See Section~\ref{proof-module-UnadjustedLangevin-Gibbs-lean}.
\item \textbf{Construct Brownian noise}. Independent Gaussian coordinates See Section~\ref{proof-module-UnadjustedLangevin-CanonicalBrownian-lean}.
\item \textbf{Bound the local Euler defect}. Integrate the Lipschitz drift error See Section~\ref{proof-module-UnadjustedLangevin-LocalError-lean}.
\item \textbf{Prove Gibbs stationarity}. Metropolis invariance and Euler limit See Section~\ref{proof-module-UnadjustedLangevin-GibbsStationarity-lean}.
\item \textbf{Solve the integral equation}. \(L^2\) Picard construction See Section~\ref{proof-module-UnadjustedLangevin-LangevinIntegral-lean}.
\item \textbf{Solve the scalar recurrence}. \(a_{k+1}\le q a_k+b\) See Section~\ref{proof-module-UnadjustedLangevin-Recurrence-lean}.
\item \textbf{Compare synchronous paths}. Noise cancellation and initial infimum See Section~\ref{proof-module-UnadjustedLangevin-HilbertDalalyan-lean}.
\item \textbf{Identify the Gaussian laws}. Independent transitions give the chain See Section~\ref{proof-module-UnadjustedLangevin-ChainLaw-lean}.
\item \textbf{Canonical ULA theorem}. Both stable step-size regimes See Section~\ref{proof-module-UnadjustedLangevin-CanonicalULA-lean}.
\end{enumerate}
\chapter{Supporting proofs}
The following mathematical arguments are reproduced in full from the authored exposition of the local proof dependencies. Each module appears after its local imports.

\section{\texorpdfstring{A first Euler-\allowbreak{}step estimate}{A first Euler-step estimate}}
\label{proof-module-UnadjustedLangevin-GradientStep-lean}
{\small\color{gray}Source module: \texttt{Unadjusted\allowbreak{}Langevin/\allowbreak{}Gradient\allowbreak{}Step.\allowbreak{}lean}.\par}


Before using the full lower and upper models together, it is useful to record the estimate obtained directly from strong monotonicity and a Lipschitz bound. This estimate applies to arbitrary vectors \(g_x,g_y\); they need not arise as gradients.



The formal inputs \texttt{gx} and \texttt{gy} are arbitrary vectors satisfying the displayed monotonicity and Lipschitz inequalities. Thus \texttt{gradient\_\allowbreak{}step\_\allowbreak{}sq\_\allowbreak{}bound} is an algebraic Hilbert-space result; the potential's actual gradients are substituted later. \texttt{gradient\_\allowbreak{}step\_\allowbreak{}contract} takes a nonnegative square root and does not require a separate sign assumption on the scalar coefficient.

\subsection{\texorpdfstring{The squared Euler-\allowbreak{}step bound}{The squared Euler-step bound}}


Let \(x,y,g_x,g_y\) lie in a real inner-product space, let \(m\in\mathbb R\), \(L\ge0\), and \(h\ge0\). Assume

\[m\|x-y\|^2\le\langle g_x-g_y,x-y\rangle,\qquad\|g_x-g_y\|\le L\|x-y\|.\]

Then

\[\|(x-hg_x)-(y-hg_y)\|^2\le(1-2hm+h^2L^2)\|x-y\|^2.\]

Both sides of the Lipschitz estimate are nonnegative, so squaring preserves the inequality:

\[\|g_x-g_y\|^2\le L^2\|x-y\|^2.\]

Multiply strong monotonicity by \(2h\ge0\), and multiply the squared Lipschitz estimate by \(h^2\ge0\):

\[2hm\|x-y\|^2\le2h\langle g_x-g_y,x-y\rangle,\qquad h^2\|g_x-g_y\|^2\le h^2L^2\|x-y\|^2.\]

Write \(d=x-y\), \(v=g_x-g_y\). Expanding the squared norm and applying the preceding estimates gives

\[\begin{aligned}\|(x-hg_x)-(y-hg_y)\|^2&=\|d-hv\|^2\\&=\|d\|^2-2h\langle v,d\rangle+h^2\|v\|^2\\&\le(1-2hm+h^2L^2)\|d\|^2.\end{aligned}\]\subsection{\texorpdfstring{Taking the square root}{Taking the square root}}


Let \(H\) be a real inner-product space, \(x,y,g_x,g_y\in H\), \(m\in\mathbb R\), \(L\ge0\), and \(h\ge0\). Assume

\[m\|x-y\|^2\le\langle g_x-g_y,x-y\rangle,\qquad\|g_x-g_y\|\le L\|x-y\|.\]

Then

\[\|(x-hg_x)-(y-hg_y)\|\le\sqrt{1-2hm+h^2L^2}\,\|x-y\|.\]

No separate sign assumption on \(1-2hm+h^2L^2\) is required; Lean uses the real square root, which is zero on nonpositive inputs.



Apply the preceding squared estimate. With \(A=\|(x-hg_x)-(y-hg_y)\|\), \(r=\|x-y\|\), and \(c=1-2hm+h^2L^2\), it says \(A^2\le cr^2\).



Since \(A\ge0\), take square roots and use \(r^2\ge0\):

\[A\le\sqrt{cr^2}=\sqrt c\,\sqrt{r^2}=\sqrt c\,r.\]

The square-root product identity is valid here even without a sign hypothesis on \(c\), because the second factor \(r^2\) is nonnegative.

\section{\texorpdfstring{Supporting planes and strong monotonicity}{Supporting planes and strong monotonicity}}
\label{proof-module-UnadjustedLangevin-StrongConvexity-lean}
{\small\color{gray}Source module: \texttt{Unadjusted\allowbreak{}Langevin/\allowbreak{}Strong\allowbreak{}Convexity.\allowbreak{}lean}.\par}


We begin with the deterministic geometry behind synchronous Langevin coupling. Throughout this part, \(\langle\cdot,\cdot\rangle\) and \(\|\cdot\|\) are the real inner product and its norm. Strong convexity uses the normalization \(m/2\): the function \(v\mapsto f(v)-\frac m2\|v\|^2\) is convex.



The formal normalization is \texttt{Strong\allowbreak{}Convex\allowbreak{}On Set.\allowbreak{}univ m f}, meaning convexity after subtracting \(m\|x\|^2/2\). \texttt{fderiv \ensuremath{\mathbb R} f x} is the continuous linear derivative; \texttt{gradient f x} is its Hilbert-space representative, so \(Df(x)[v]=\langle\nabla f(x),v\rangle\). The proofs explicitly move between these two representations.

\subsection{\texorpdfstring{The supporting-\allowbreak{}plane inequality}{The supporting-plane inequality}}


Let \(H\) be a real inner-product space and let \(f:H\to\mathbb R\) be differentiable and convex on all of \(H\). For any \(x,y\in H\),

\[f(x)+Df(x)[y-x]\le f(y).\]

Restrict \(f\) to the affine line \(\ell(t)=x+t(y-x)\). Since affine composition preserves convexity, \(\varphi=f\circ\ell\) is convex on \(\mathbb R\), with \(\varphi(0)=f(x)\) and \(\varphi(1)=f(y)\).



The derivative of the line is its constant direction:

\[\ell'(0)=y-x.\]

By the chain rule, the derivative of the restricted function is

\[\varphi'(0)=Df(x)[y-x].\]

For a differentiable convex function, the derivative at the left endpoint is bounded by the secant slope to a later point. Applying this at \(0<1\) gives

\[Df(x)[y-x]\le\frac{\varphi(1)-\varphi(0)}{1-0}.\]

Evaluate the two endpoints and rearrange:

\[Df(x)[y-x]\le f(y)-f(x),\qquad f(x)+Df(x)[y-x]\le f(y).\]\subsection{\texorpdfstring{Strong monotonicity of the gradient}{Strong monotonicity of the gradient}}


Let \(H\) be a real Hilbert space. If \(f:H\to\mathbb R\) is differentiable and \(m\)-strongly convex, where \(m\in\mathbb R\), then for every \(x,y\in H\),

\[m\|x-y\|^2\le\langle\nabla f(x)-\nabla f(y),x-y\rangle.\]

Remove the strong-convexity quadratic. Set

\[g(v)=f(v)-\frac m2\|v\|^2.\]

The definition of strong convexity says precisely that \(g\) is convex.



The squared norm has derivative \(D\|v\|^2[u]=2\langle v,u\rangle\). Consequently \(g\) is differentiable and

\[Dg(v)[u]=Df(v)[u]-m\langle v,u\rangle.\]

Apply the supporting-plane inequality in both directions, and substitute the derivative just computed:

\[g(x)+Dg(x)[y-x]\le g(y),\qquad g(y)+Dg(y)[x-y]\le g(x).\]

Express the derivative of \(f\) using its gradient. The two inequalities become

\[\begin{aligned}f(x)-\tfrac m2\|x\|^2+\langle\nabla f(x),y-x\rangle-m\langle x,y-x\rangle&\le f(y)-\tfrac m2\|y\|^2,\\ f(y)-\tfrac m2\|y\|^2+\langle\nabla f(y),x-y\rangle-m\langle y,x-y\rangle&\le f(x)-\tfrac m2\|x\|^2.\end{aligned}\]

Add these inequalities. The function values and quadratic values cancel. Using \(\langle x-y,x-y\rangle=\|x-y\|^2\), what remains is

\[0\le\langle\nabla f(x)-\nabla f(y),x-y\rangle-m\|x-y\|^2.\]

This is the required monotonicity inequality.

\subsection{\texorpdfstring{The quadratic lower model}{The quadratic lower model}}


Let \(H\) be a real Hilbert space, \(f:H\to\mathbb R\) differentiable, and \(m\in\mathbb R\). If \(f\) is \(m\)-strongly convex on all of \(H\), then for every \(x,y\in H\),

\[f(x)+\langle\nabla f(x),y-x\rangle+\frac m2\|y-x\|^2\le f(y).\]

Set \(g(v)=f(v)-\frac m2\|v\|^2\). Strong convexity makes \(g\) convex.



Differentiate this shifted potential:

\[Dg(v)[u]=Df(v)[u]-m\langle v,u\rangle.\]

In particular, \(g\) is differentiable on \(H\).



Its supporting-plane inequality at \(x\), together with \(Df(x)[u]=\langle\nabla f(x),u\rangle\), yields

\[f(x)-\frac m2\|x\|^2+\langle\nabla f(x),y-x\rangle-m\langle x,y-x\rangle\le f(y)-\frac m2\|y\|^2.\]

Collect the quadratic correction by the identity

\[\|y\|^2-\|x\|^2-2\langle x,y-x\rangle=\|y-x\|^2.\]

Moving terms gives the desired lower model.

\section{\texorpdfstring{The Langevin update and cancellation of shared noise}{The Langevin update and cancellation of shared noise}}
\label{proof-module-UnadjustedLangevin-ULAStep-lean}
{\small\color{gray}Source module: \texttt{Unadjusted\allowbreak{}Langevin/\allowbreak{}ULAStep.\allowbreak{}lean}.\par}


We now work in finite-dimensional Euclidean space. A ULA step adds the same scaled noise vector to a gradient step. When two starting points use this same noise, it cancels exactly from their difference. The deterministic estimates can therefore be applied pointwise, before introducing any probability law.



\texttt{ula\allowbreak{}Update} is the pointwise Euler map. Its definition accepts arbitrary vector noise; Gaussian laws enter later in realization theorems. \texttt{ula\allowbreak{}Update\_\allowbreak{}sub} uses the same innovation on both sides, so that contribution cancels exactly. The monotonicity of the actual gradient is supplied by the strong-convexity theorem when applying the deterministic step bound.

\subsection{\texorpdfstring{The Euclidean state space}{The Euclidean state space}}


For a natural number \(d\), write \(E_d=\mathbb R^d\) with its Euclidean inner product and norm. Lean represents its vectors as real-valued functions on \(\operatorname{Fin}(d)\), the set of indices \(0,\ldots,d-1\); when \(d=0\), this index set is empty.



Introduce the abbreviation

\[E_d:=\operatorname{EuclideanSpace}_{\mathbb R}(\operatorname{Fin}(d)).\]

This equips the coordinate functions with the Euclidean norm \(\|x\|=(\sum_{i=0}^{d-1}x_i^2)^{1/2}\), so the subsequent gradient and distance estimates use the usual geometry of \(\mathbb R^d\).

\subsection{\texorpdfstring{One unadjusted Langevin step}{One unadjusted Langevin step}}


For a potential \(U:E_d\to\mathbb R\), a real step size \(h\), a state \(x\in E_d\), and a noise vector \(z\in E_d\), define

\[T_h^z(x)=x-h\nabla U(x)+\sqrt{2h}\,z.\]

The definition itself imposes no sign condition on \(h\); later contraction results assume \(h\ge0\).



The update consists of the current state, the drift increment \(-h\nabla U(x)\), and the noise increment \(\sqrt{2h}\,z\).

\subsection{\texorpdfstring{Shared noise cancels}{Shared noise cancels}}


For every potential \(U:E_d\to\mathbb R\), \(h\in\mathbb R\), and vectors \(x,y,z\in E_d\),

\[T_h^z(x)-T_h^z(y)=(x-h\nabla U(x))-(y-h\nabla U(y)).\]

Expand the two updates. The two copies of \(\sqrt{2h}\,z\) subtract to zero:

\[(x-h\nabla U(x)+\sqrt{2h}\,z)-(y-h\nabla U(y)+\sqrt{2h}\,z)=(x-h\nabla U(x))-(y-h\nabla U(y)).\]\subsection{\texorpdfstring{The initial shared-\allowbreak{}noise estimate}{The initial shared-noise estimate}}


Suppose \(U:E_d\to\mathbb R\) is differentiable and \(m\)-strongly convex, \(\nabla U\) is \(L\)-Lipschitz with \(L\ge0\), and \(h\ge0\). For every \(x,y,z\in E_d\),

\[\|T_h^z(x)-T_h^z(y)\|\le\sqrt{1-2hm+h^2L^2}\,\|x-y\|.\]

For \(h>0\) this factor is strictly below one whenever \(hL^2<2m\).



Cancel the common noise. The left-hand side becomes the distance between the two deterministic gradient steps.



Strong convexity gives

\[m\|x-y\|^2\le\langle\nabla U(x)-\nabla U(y),x-y\rangle,\]

and the Lipschitz assumption gives \(\|\nabla U(x)-\nabla U(y)\|\le L\|x-y\|\). Insert these two inequalities into the Euler-step square-root bound.

\section{\texorpdfstring{Gaussian and Brownian noise moments}{Gaussian and Brownian noise moments}}
\label{proof-module-UnadjustedLangevin-GaussianMoments-lean}
{\small\color{gray}Source module: \texttt{Unadjusted\allowbreak{}Langevin/\allowbreak{}Gaussian\allowbreak{}Moments.\allowbreak{}lean}.\par}


Fix a natural-number dimension \(d\in\mathbb N\), write \(E_d=\mathbb R^d\) with its Euclidean norm, and let \(P\) be a measure on \(\Omega\). A law statement \(X_i\sim\mathcal N(0,v)\) below means that the coordinate is almost everywhere measurable and has the stated pushforward measure. We prove the required \(L^2\) norms from these laws. Coordinate independence is unnecessary for second moments because \(\|X\|^2=\sum_iX_i^2\) has no cross terms.



For Brownian processes the input is mathlib’s actual scalar pre-Brownian law predicate, whose marginal and increment laws supply the Gaussian hypotheses. These facts establish the noise moments; existence and stationarity of a coupled Langevin reference require separate inputs.



The \(L^2\) norms and expectations below concern the actual measure \(P\), with \texttt{Mem\allowbreak{}Lp} giving square integrability before passing to a quotient class. Laws are pushforward identities; independence is used only for cross-term cancellation, while marginal Gaussian second moments supply the dimension factor.

\subsection{\texorpdfstring{Second moment of a centered Gaussian}{Second moment of a centered Gaussian}}


For nonnegative variance \(v\), the centered Gaussian law has second moment

\[\int_{\mathbb R}x^2\,d\mathcal N(0,v)(x)=v.\]

The known Gaussian variance is \(v\). Expand variance as \(\int(x-\mathbb E[x])^2\), and insert the Gaussian mean \(\mathbb E[x]=0\).



Simplify \(x-0=x\). The variance integral is now exactly the required second moment.

\subsection{\texorpdfstring{Gaussian coordinates give a square-\allowbreak{}integrable vector}{Gaussian coordinates give a square-integrable vector}}


Let \(X:\Omega\to E_d\), and suppose every coordinate satisfies \(X_i\sim\mathcal N(0,v)\) under \(P\), with \(v\ge0\). Then \(X\in L^2(P;E_d)\).



In finite-dimensional Euclidean space, square integrability is equivalent to square integrability of each coordinate. Each coordinate belongs to \(L^2\) because its Gaussian law has a finite second moment; transfer that membership along the equality of laws.

\subsection{\texorpdfstring{Sum the coordinate moments}{Sum the coordinate moments}}


Under the same coordinate Gaussian laws,

\[\int_\Omega\|X(\omega)\|^2\,dP(\omega)=dv.\]

For each coordinate \(i\), retain the proved square integrability \(X_i\in L^2(P)\). It will justify integrating the finite sum of squared coordinates.



Change variables using the law of \(X_i\), then apply the scalar Gaussian moment:

\[\int_\Omega X_i^2\,dP=\int_{\mathbb R}x^2\,d\mathcal N(0,v)=v.\]

Expand the squared Euclidean norm and interchange the finite sum and integral:

\[\int_\Omega\|X\|^2\,dP=\sum_{i=0}^{d-1}\int_\Omega X_i^2\,dP.\]

Every summand is integrable by its \(L^2\) membership.



Each of the \(d\) summands equals \(v\); hence their sum is \(dv\).

\subsection{\texorpdfstring{The Gaussian vector as an L\ensuremath{{}^{2}} class}{The Gaussian vector as an L² class}}


Let \([X]\) be the \(L^2\) equivalence class constructed from \(X\), whose coordinates have laws \(\mathcal N(0,v)\). Then

\[\|[X]\|_{L^2}^2=dv.\]

Express the squared \(L^2\) norm as the self-inner-product, then use the definition of the \(L^2\) inner product as an integral.



The chosen representative of \([X]\) equals \(X\) almost everywhere. Therefore the integrand \(\langle[X](\omega),[X](\omega)\rangle\) agrees almost everywhere with \(\|X(\omega)\|^2\), so the two integrals agree.



Substitute the integral identity and the vector second-moment calculation \(\int\|X\|^2\,dP=dv\).

\subsection{\texorpdfstring{The square-\allowbreak{}root Gaussian norm}{The square-root Gaussian norm}}


For the same Gaussian coordinate laws,

\[\|[X]\|_{L^2}=\sqrt{dv}.\]

This includes \(d=0\) and \(v=0\).



Replace \(dv\) by \(\|[X]\|^2\). Since norms are nonnegative, \(\sqrt{\|[X]\|^2}=\|[X]\|\).

\subsection{\texorpdfstring{Include the Langevin diffusion factor}{Include the Langevin diffusion factor}}


If every coordinate of \(X\) has law \(\mathcal N(0,v)\), then the Langevin-scaled vector satisfies

\[\|\sqrt2[X]\|_{L^2}=\sqrt{2d}\sqrt v.\]

Norm homogeneity gives \(\|\sqrt2[X]\|=\sqrt2\|[X]\|\), because \(\sqrt2\ge0\). Insert \(\|[X]\|=\sqrt{dv}\), then split the square roots of the nonnegative products.



Rearrange \(\sqrt2\sqrt d\sqrt v\) into \(\sqrt{2d}\sqrt v\).

\subsection{\texorpdfstring{Assemble Brownian coordinates}{Assemble Brownian coordinates}}


For scalar processes \(W_i:[0,\infty)\times\Omega\to\mathbb R\), define the Euclidean vector

\[W_t(\omega):=(W_0(t,\omega),\ldots,W_{d-1}(t,\omega))\in E_d.\]

Collect the coordinates at the fixed time \(t\) and outcome \(\omega\), using the finite-dimensional \(\ell^2\) norm.

\subsection{\texorpdfstring{Brownian marginals are Gaussian}{Brownian marginals are Gaussian}}


Assume each scalar \(W_i\) satisfies \(\mathrm{IsPreBrownianReal}\) under \(P\). For each \(t\ge0\) and coordinate \(i\),

\[W_t^{(i)}\sim\mathcal N(0,t).\]

Apply the scalar Brownian marginal-law theorem at time \(t\). The \(i\)-th coordinate of the assembled vector is definitionally the original scalar process.

\subsection{\texorpdfstring{Brownian vectors belong to L\ensuremath{{}^{2}}}{Brownian vectors belong to L²}}


Under the scalar pre-Brownian assumptions, for every \(t\ge0\),

\[W_t\in L^2(P;E_d).\]

Apply Gaussian vector square integrability with variance \(v=t\), using the just-proved marginal laws of all coordinates.

\subsection{\texorpdfstring{The exact Brownian noise norm}{The exact Brownian noise norm}}


For scalar pre-Brownian coordinates and \(t\ge0\),

\[\|\sqrt2[W_t]\|_{L^2}=\sqrt{2d}\sqrt t.\]

Specialize the scaled Gaussian vector norm to the Brownian vector, whose coordinate variance is \(t\). This derives the square-root noise scale from the Brownian laws.

\subsection{\texorpdfstring{Define an increment from any starting time}{Define an increment from any starting time}}


For \(a,t\ge0\), define the vector increment over the interval \([a,a+t]\) by

\[\Delta_aW_t(\omega):=W_{a+t}(\omega)-W_a(\omega).\]

Subtract the two Euclidean vectors at the interval endpoints. Each coordinate is the corresponding scalar Brownian increment.

\subsection{\texorpdfstring{Increment variance is elapsed time}{Increment variance is elapsed time}}


Under the scalar pre-Brownian assumptions, for every \(a,t\ge0\) and every coordinate \(i\),

\[W_i(a+t)-W_i(a)\sim\mathcal N(0,t).\]

Expand the vector increment to its scalar coordinate. Apply the Brownian difference-law theorem at times \(a+t\) and \(a\); its variance is the nonnegative distance between those times.



It remains to identify that variance. As real numbers,

\[|(a+t)-a|=|t|=t,\]

because \(t\ge0\). Thus the increment law is exactly \(\mathcal N(0,t)\).

\subsection{\texorpdfstring{Increments are square integrable}{Increments are square integrable}}


For pre-Brownian coordinates and any \(a,t\ge0\),

\[\Delta_aW_t\in L^2(P;E_d).\]

Each increment coordinate has the centered Gaussian law of variance \(t\), so the Gaussian vector integrability theorem applies.

\subsection{\texorpdfstring{The same noise norm on every reference segment}{The same noise norm on every reference segment}}


For scalar pre-Brownian coordinates and every \(a,t\ge0\),

\[\boxed{\|\sqrt2[W_{a+t}-W_a]\|_{L^2}=\sqrt{2d}\sqrt t.}\]

Apply the scaled Gaussian vector norm to the increment coordinate laws. The answer depends only on elapsed time \(t\), so every reference interval has the same moment bound.

\section{\texorpdfstring{Characteristic functions are continuous in L\ensuremath{{}^{2}}}{Characteristic functions are continuous in L²}}
\label{proof-module-UnadjustedLangevin-LpCharacteristic-lean}
{\small\color{gray}Source module: \texttt{Unadjusted\allowbreak{}Langevin/\allowbreak{}Lp\allowbreak{}Characteristic.\allowbreak{}lean}.\par}


At each fixed frequency, the characteristic function of the law of a real \(L^2\) variable depends continuously on that variable. Consequently centered Gaussian laws are closed under \(L^2\) convergence, with variance given by the squared norm.

\subsection{\texorpdfstring{centered\allowbreak{}Characteristic}{centeredCharacteristic}}


Set \(c_t(x)=e^{itx}-1\), represented by the unit-radius circle map centered at \(-1\). Subtracting one makes \(c_t(0)=0\), which is the condition needed by the Lipschitz composition construction on \(L^2\).

\subsection{\texorpdfstring{centered\allowbreak{}Characteristic lipschitz}{centeredCharacteristic lipschitz}}


Multiplication by \(t\) is \(|t|\)-Lipschitz, and the unit circle parametrization is \(1\)-Lipschitz. Composing gives \(|c_t(x)-c_t(y)|\le|t||x-y|\). The proof identifies the nonnegative Lipschitz constants before applying the composition theorem.

\subsection{\texorpdfstring{centered\allowbreak{}Characteristic zero}{centeredCharacteristic zero}}


At zero the exponential is one, hence \(c_t(0)=e^0-1=0\). Unfolding the circle-map definition makes this a scalar identity; this vanishing proof permits \texttt{comp\allowbreak{}Lp} to lift the map without an integrability assumption on a nonzero constant.

\subsection{\texorpdfstring{characteristic\allowbreak{}Lp}{characteristicLp}}


Lift \(c_t\) to \(L^2\) by Lipschitz composition, then add the constant class one. Thus \(\operatorname{characteristicLp}(t,X)\) represents \(e^{itX}\). The probability assumption makes the constant function square integrable.

\subsection{\texorpdfstring{characteristic\allowbreak{}Lp coe}{characteristicLp coe}}


Combine the almost-everywhere representative identities for lifted composition, \(L^2\) addition, and the constant class. Their intersection has full measure; there the representative is \((e^{itX}-1)+1=e^{itX}\). This connects quotient-space operations with the random function used in the characteristic integral.

\subsection{\texorpdfstring{char\allowbreak{}Fun map eq inner characteristic\allowbreak{}Lp}{charFun map eq inner characteristicLp}}


The pushforward integral theorem gives \(\widehat{P_X}(t)=\int e^{itX}dP\). In complex \(L^2\), inner product with the constant one has exactly this integral because its conjugate is one. The representative identity above therefore yields \(\widehat{P_X}(t)=\langle1,\operatorname{characteristicLp}(t,X)\rangle\).

\subsection{\texorpdfstring{continuous char\allowbreak{}Fun map Lp}{continuous charFun map Lp}}


Lipschitz composition is continuous on \(L^2\); adding the constant one preserves continuity. Pairing with one is a continuous linear functional. The preceding integral identity rewrites their composition as \(X\mapsto\widehat{P_X}(t)\), proving continuity of the actual-law characteristic function at every fixed frequency.

\subsection{\texorpdfstring{continuous gaussian\allowbreak{}Real char\allowbreak{}Fun Lp}{continuous gaussianReal charFun Lp}}


The centered Gaussian characteristic function is \(e^{-t^2\|X\|_2^2/2}\). Norm, squaring, conversion of a nonnegative real to \texttt{\ensuremath{\mathbb R}\ensuremath{\ge}0}, and the exponential are continuous. Lean unfolds the Gaussian formula and composes these continuity facts; the variance is the actual squared \(L^2\) norm.

\subsection{\texorpdfstring{has\allowbreak{}Law gaussian\allowbreak{}Real of tendsto Lp eventually}{hasLaw gaussianReal of tendsto Lp eventually}}


At each frequency, \(L^2\) convergence makes the actual characteristic functions converge. It also makes the centered Gaussian characteristic functions with variances \(\|X_i\|_2^2\) converge to variance \(\|Y\|_2^2\). Eventual Gaussianity identifies the two nets eventually, so uniqueness of limits identifies their limits. \texttt{Measure.\allowbreak{}ext\_\allowbreak{}of\_\allowbreak{}char\allowbreak{}Fun} then gives the claimed Gaussian law of \(Y\); a nontrivial filter is required for uniqueness.

\subsection{\texorpdfstring{has\allowbreak{}Law gaussian\allowbreak{}Real of tendsto Lp}{hasLaw gaussianReal of tendsto Lp}}


If every approximant has the stated Gaussian law, then the law identity holds eventually for any filter. Feed \texttt{Eventually.\allowbreak{}of\_\allowbreak{}forall} to the preceding closure theorem. This version is suited to finite orthogonal sums converging to an infinite Gaussian series.

\section{\texorpdfstring{Construct an isonormal Gaussian family from product coordinates}{Construct an isonormal Gaussian family from product coordinates}}
\label{proof-module-UnadjustedLangevin-GaussianHilbert-lean}
{\small\color{gray}Source module: \texttt{Unadjusted\allowbreak{}Langevin/\allowbreak{}Gaussian\allowbreak{}Hilbert.\allowbreak{}lean}.\par}


An explicit product of standard normal coordinates realizes a linear isometry \(I:H\to L^2\) with \(I(x)\sim\mathcal N(0,\|x\|^2)\) and covariance \(\mathbb E[I(x)I(y)]=\langle x,y\rangle\).

\subsection{\texorpdfstring{source}{source}}


Use the product probability \(Q=\bigotimes_{i\in\iota}\mathcal N(0,1)\) on real coordinate families. The construction uses \texttt{Measure.\allowbreak{}infinite\allowbreak{}Pi}; independence and coordinate marginals are properties of this concrete source rather than extra assumptions on an unspecified Gaussian process.

\subsection{\texorpdfstring{coordinate has\allowbreak{}Law}{coordinate hasLaw}}


The coordinate projection \(\omega\mapsto\omega_i\) pushes the product measure to its \(i\)-th factor \(\mathcal N(0,1)\). The product marginal theorem proves both almost-everywhere measurability and the pushforward identity, packaged as \texttt{Has\allowbreak{}Law}.

\subsection{\texorpdfstring{coordinate\allowbreak{}Lp}{coordinateLp}}


A standard Gaussian has finite second moment, so the coordinate projection belongs to \(L^2(Q)\). Apply its \texttt{Mem\allowbreak{}Lp.\allowbreak{}to\allowbreak{}Lp} construction to obtain the quotient class \(Z_i\); the definition stores the original coordinate function together with its integrability proof.

\subsection{\texorpdfstring{coordinate\allowbreak{}Lp coe}{coordinateLp coe}}


The representative selected for \(Z_i\) agrees with \(\omega_i\) almost everywhere by the defining property of \texttt{to\allowbreak{}Lp}. This equality is the bridge used to transfer coordinate laws, means, and independence calculations to the \(L^2\) classes.

\subsection{\texorpdfstring{coordinates independent}{coordinates independent}}


The coordinate projections of an infinite product measure are jointly independent. Invoke that product theorem for the scalar Gaussian factors. This is independence of every finite set of coordinates, which is precisely what finite Gaussian linear combinations require.

\subsection{\texorpdfstring{coordinate mean}{coordinate mean}}


Change variables through the coordinate law and use the centered Gaussian mean: \(\int\omega_i\,dQ=\int x\,d\mathcal N(0,1)=0\). Gaussian integrability makes this an ordinary expectation identity.

\subsection{\texorpdfstring{coordinate\allowbreak{}Lp inner}{coordinateLp inner}}


For \(i=j\), the \(L^2\) inner product is \(\mathbb EZ_i^2=1\). For \(i\ne j\), independence factors \(\mathbb EZ_iZ_j=(\mathbb EZ_i)(\mathbb EZ_j)=0\). Almost-everywhere representative identities move between \(Z_i\) and the product coordinates. Thus \(\langle Z_i,Z_j\rangle=\mathbf1_{i=j}\).

\subsection{\texorpdfstring{coordinates orthonormal}{coordinates orthonormal}}


The preceding inner-product formula is the orthonormality criterion for a family in a real Hilbert space. Lean converts the diagonal value one and off-diagonal value zero into \texttt{Orthonormal \ensuremath{\mathbb R} coordinate\allowbreak{}Lp}, the input for the orthogonal-series extension.

\subsection{\texorpdfstring{source is\allowbreak{}Probability\allowbreak{}Measure}{source isProbabilityMeasure}}


Each scalar Gaussian factor is a probability measure. The product-probability instance therefore gives \(Q(\Omega)=1\). This instance is what permits constant functions and integrable functions to be treated in the usual probability-space \(L^2\) theory below.

\subsection{\texorpdfstring{synthesis}{synthesis}}


Extend the orthonormal family \(Z_i\) to a linear isometry \(S:\ell^2(\iota)\to L^2(Q)\). The orthogonal-family extension theorem constructs the limit of finite sums and preserves squared norm, \(\|S(a)\|_2^2=\sum_i a_i^2\). It constructs an \(L^2\) limit rather than assuming pointwise convergence of a random series.

\subsection{\texorpdfstring{synthesis has\allowbreak{}Sum}{synthesis hasSum}}


The synthesis extension theorem states \(\sum_i a_iZ_i=S(a)\), with convergence of finite partial sums in \(L^2\). In Lean \texttt{Has\allowbreak{}Sum} is convergence along finite subsets, so the identity does not require choosing an enumeration of the basis index set.

\subsection{\texorpdfstring{isonormal}{isonormal}}


For a Hilbert basis \(b\), compose its coordinate isometry \(b.\mathrm{repr}:H\to\ell^2(\iota)\) with \texttt{synthesis}. The resulting linear isometry \(I_b\) sends \(x\) to the Gaussian series with coefficients \(\langle b_i,x\rangle\), and preserves all inner products.

\subsection{\texorpdfstring{isonormal has\allowbreak{}Sum}{isonormal hasSum}}


Apply \texttt{synthesis\_\allowbreak{}has\allowbreak{}Sum} to the coefficient family \(b.\mathrm{repr}(x)\). Unfolding the composition gives \(I_b(x)=\sum_i\langle b_i,x\rangle Z_i\) in \(L^2\). This is the exact series represented by the bundled isonormal isometry.

\subsection{\texorpdfstring{integral eq inner one}{integral eq inner one}}


On a probability space, the constant one is in \(L^2\). The definition of the real \(L^2\) inner product gives \(\mathbb EX=\langle1,X\rangle\), after replacing the selected constant representative by one almost everywhere. This exhibits expectation as a continuous functional.

\subsection{\texorpdfstring{coordinate\allowbreak{}Lp mean}{coordinateLp mean}}


Use \(Z_i=\omega_i\) almost everywhere to replace the \(L^2\) representative inside the integral. The already established centered coordinate expectation then gives \(\mathbb EZ_i=0\).

\subsection{\texorpdfstring{synthesis mean}{synthesis mean}}


Each finite partial sum \(\sum_{i\in F}a_iZ_i\) has mean zero by linearity and the coordinate means. Pair the \(L^2\)-convergent series with the constant one; continuity of the inner product passes the zero expectations to the limit, giving \(\mathbb ES(a)=0\).

\subsection{\texorpdfstring{isonormal mean}{isonormal mean}}


The isonormal variable is \texttt{synthesis} applied to the Hilbert-basis coefficients of \(x\). Specialize the synthesis mean theorem to these coefficients to get \(\mathbb EI_b(x)=0\), without redoing the infinite-series limit.

\subsection{\texorpdfstring{variance eq norm sq}{variance eq norm sq}}


For an \(L^2\) random variable with mean zero, the variance formula reduces to \(\operatorname{Var}(X)=\int X^2dQ\). The \(L^2\) norm identity identifies this with \(\|X\|_2^2\). The explicit mean premise is used to remove the centering term.

\subsection{\texorpdfstring{partial sum coe}{partial sum coe}}


Induct over the finite set. The zero sum has zero representative almost everywhere; addition and scalar multiplication of \(L^2\) classes agree almost everywhere with the corresponding pointwise operations. Intersect the finitely many full-measure sets to obtain \(\sum_{i\in F}a_iZ_i=\sum_{i\in F}a_i\omega_i\) almost everywhere.

\subsection{\texorpdfstring{partial sum has\allowbreak{}Gaussian\allowbreak{}Law}{partial sum hasGaussianLaw}}


Independent scalar standard Gaussians remain Gaussian after a finite linear combination. Apply that finite-sum Gaussian theorem to \(a_i\omega_i\), and transfer the law across \texttt{partial\_\allowbreak{}sum\_\allowbreak{}coe}. This proves Gaussianity of the actual \(L^2\) partial sum, not merely of a formally similar expression.

\subsection{\texorpdfstring{partial sum mean}{partial sum mean}}


The integral of a finite \(L^2\) sum is the sum of its integrals, because every summand is integrable. Each coordinate mean is zero, hence \(\mathbb E\sum_{i\in F}a_iZ_i=0\). This supplies the centering parameter needed to identify its Gaussian law.

\subsection{\texorpdfstring{partial sum has\allowbreak{}Law}{partial sum hasLaw}}


A scalar Gaussian law is determined by its mean and variance. The finite sum is Gaussian and centered; its variance is its squared \(L^2\) norm. Therefore its law is \(\mathcal N(0,\|\sum_{i\in F}a_iZ_i\|_2^2)\), with the variance converted to a nonnegative real.

\subsection{\texorpdfstring{synthesis has\allowbreak{}Law}{synthesis hasLaw}}


Finite partial sums converge in \(L^2\) to \(S(a)\), and each has the centered Gaussian law with its own squared norm as variance. Apply \texttt{has\allowbreak{}Law\_\allowbreak{}gaussian\allowbreak{}Real\_\allowbreak{}of\_\allowbreak{}tendsto\_\allowbreak{}Lp}; the isometry then rewrites the limiting variance as \(\|a\|_{\ell^2}^2\). This proves Gaussianity of the completed infinite series.

\subsection{\texorpdfstring{isonormal has\allowbreak{}Law}{isonormal hasLaw}}


Apply the synthesis law theorem to \(b.\mathrm{repr}(x)\). The Hilbert-basis coordinate map preserves norm, so the variance becomes \(\|x\|^2\). Hence \(I_b(x)\sim\mathcal N(0,\|x\|^2)\), including the degenerate law for \(x=0\).

\subsection{\texorpdfstring{isonormal has\allowbreak{}Gaussian\allowbreak{}Law}{isonormal hasGaussianLaw}}


The exact law theorem identifies \(I_b(x)\) with a scalar Gaussian probability. Forgetting its explicit mean and variance parameters yields the bundled \texttt{Has\allowbreak{}Gaussian\allowbreak{}Law} predicate used by the finite-dimensional Gaussian-process construction.

\subsection{\texorpdfstring{isonormal covariance}{isonormal covariance}}


Both variables are centered, so their covariance is \(\mathbb E[I_b(x)I_b(y)]\), the \(L^2\) inner product. The linear isometry preserves that inner product, giving \(\operatorname{Cov}(I_b(x),I_b(y))=\langle x,y\rangle_H\). Lean checks integrability through their \(L^2\) membership.

\section{\texorpdfstring{Scalar Gaussian linear maps give jointly Gaussian processes}{Scalar Gaussian linear maps give jointly Gaussian processes}}
\label{proof-module-UnadjustedLangevin-GaussianLinearProcess-lean}
{\small\color{gray}Source module: \texttt{Unadjusted\allowbreak{}Langevin/\allowbreak{}Gaussian\allowbreak{}Linear\allowbreak{}Process.\allowbreak{}lean}.\par}


If a linear map \(T:H\to L^2(P;\mathbb R)\) sends every input to a Gaussian variable, all finite collections of its evaluations are jointly Gaussian.

\subsection{\texorpdfstring{has\allowbreak{}Gaussian\allowbreak{}Law pi of linear\allowbreak{}Map Lp}{hasGaussianLaw pi of linearMap Lp}}


Use the finite-dimensional Gaussian criterion: every scalar linear functional of \((T(f_i))_i\) must be Gaussian. A finite linear combination equals \(T(\sum_i a_if_i)\) as an \(L^2\) class by linearity, hence agrees almost everywhere as a random variable. The assumed Gaussianity of every \(T(x)\) proves the criterion.

\subsection{\texorpdfstring{is\allowbreak{}Gaussian\allowbreak{}Process of linear\allowbreak{}Map Lp}{isGaussianProcess of linearMap Lp}}


For each finite set of times, apply the preceding joint-Gaussian theorem to the corresponding finite family of vectors \(f(t)\). This verifies the finite-dimensional-law definition of \texttt{Is\allowbreak{}Gaussian\allowbreak{}Process}. The index set may be uncountable; this statement alone asserts no sample-path continuity.

\section{\texorpdfstring{An explicit Gaussian construction of pre-\allowbreak{}Brownian coordinates}{An explicit Gaussian construction of pre-Brownian coordinates}}
\label{proof-module-UnadjustedLangevin-CanonicalBrownian-lean}
{\small\color{gray}Source module: \texttt{Unadjusted\allowbreak{}Langevin/\allowbreak{}Canonical\allowbreak{}Brownian.\allowbreak{}lean}.\par}


The isonormal family on the time Hilbert space constructs a Gaussian process with covariance \(\min(s,t)\). Independent product copies supply the Brownian coordinates required by the canonical ULA theorem. The construction proves pre-Brownian laws; continuous sample paths are not asserted or needed, because continuity is established in \(L^2\).

\subsection{\texorpdfstring{Time\allowbreak{}Hilbert}{TimeHilbert}}


Take the deterministic Hilbert space \(L^2(\mathbb R,dt)\). Its vectors will encode time intervals before an isonormal map turns their inner products into covariances. It is distinct from the random-variable \(L^2\) space on the Gaussian source.

\subsection{\texorpdfstring{time\allowbreak{}Indicator}{timeIndicator}}


Let \(e_t=\mathbf1_{(0,t]}\) as a class in \(L^2(dt)\), for \(t\ge0\). Lean constructs it with \texttt{indicator\allowbreak{}Const\allowbreak{}Lp} and the finite Lebesgue measure of \texttt{Ioc 0 t}. Endpoints do not affect this Lebesgue \(L^2\) class, but this half-open convention makes interval intersection identities direct.

\subsection{\texorpdfstring{time\allowbreak{}Indicator inner}{timeIndicator inner}}


The product of the two indicators is the indicator of \([0,\min(s,t)]\) almost everywhere. Integrating it gives \(\langle e_s,e_t\rangle=\min(s,t)\). The proof evaluates the Lebesgue measure of that finite interval, producing the Brownian covariance before introducing randomness.

\subsection{\texorpdfstring{basis\allowbreak{}Index}{basisIndex}}


Choose the index type supplied by a Hilbert basis of the time Hilbert space. It is a set of Hilbert-space vectors used as a type; no explicit countable enumeration is required by the orthogonal-family construction.

\subsection{\texorpdfstring{basis}{basis}}


Select a Hilbert basis for \texttt{Time\allowbreak{}Hilbert} by Hilbert-space basis existence. This basis supplies the coordinate isometry for the explicit isonormal Gaussian family; all subsequent covariance identities follow from isometry and do not depend on which basis was selected.

\subsection{\texorpdfstring{Sample}{Sample}}


A sample is a real family indexed by the chosen time-Hilbert basis. This concrete product space stores the independent standard Gaussian coefficients used to synthesize all Brownian times.

\subsection{\texorpdfstring{probability}{probability}}


Use \texttt{Gaussian\allowbreak{}Hilbert.\allowbreak{}source} on the basis index type, namely the product of standard scalar Gaussians. Its established probability instance provides normalization and independent coordinate laws for the Brownian construction.

\subsection{\texorpdfstring{instance at 53}{instance at 53}}


The scalar source is the Gaussian coordinate product probability. Reuse \texttt{Gaussian\allowbreak{}Hilbert.\allowbreak{}source\_\allowbreak{}is\allowbreak{}Probability\allowbreak{}Measure} to provide normalization for the chosen basis index. This instance enables the \(L^2\) expectation and covariance theorems for the constructed Brownian process.

\subsection{\texorpdfstring{value\allowbreak{}Lp}{valueLp}}


Apply the isonormal isometry to \(e_t\), obtaining \(B_t=I(e_t)\) as an \(L^2\) random variable. The construction immediately gives a centered Gaussian with variance \(\|e_t\|^2=t\); joint Gaussianity follows from linearity.

\subsection{\texorpdfstring{process}{process}}


At each time select the strongly measurable representative of \texttt{value\allowbreak{}Lp}. This defines an actual function \(B_t(\omega)\) for use in \texttt{Has\allowbreak{}Law} and finite-dimensional distributions. Equalities of \(L^2\) classes still translate to almost-everywhere equalities of these representatives.

\subsection{\texorpdfstring{process measurable}{process measurable}}


The representative selected for an \(L^2\) class is measurable in this setting. Apply the representative measurability theorem to \texttt{value\allowbreak{}Lp t}, obtaining measurability of the actual time-\(t\) random variable.

\subsection{\texorpdfstring{process mean}{process mean}}


The isonormal mean theorem gives \(\mathbb EI(e_t)=0\). Since \texttt{process t} is precisely the selected representative of this class, this is the expectation of the displayed random function \(B_t\).

\subsection{\texorpdfstring{process covariance}{process covariance}}


The isonormal covariance formula gives \(\operatorname{Cov}(I(e_s),I(e_t))=\langle e_s,e_t\rangle\). Substitute \texttt{time\allowbreak{}Indicator\_\allowbreak{}inner} to obtain \(\min(s,t)\). This calculation establishes the Brownian covariance for the concrete source and representatives.

\subsection{\texorpdfstring{is\allowbreak{}Gaussian\allowbreak{}Process}{isGaussianProcess}}


The isonormal isometry is linear and every output is scalar Gaussian. Apply \texttt{is\allowbreak{}Gaussian\allowbreak{}Process\_\allowbreak{}of\_\allowbreak{}linear\allowbreak{}Map\_\allowbreak{}Lp} to the family \(t\mapsto e_t\). Every finite collection of the selected random variables is therefore jointly Gaussian.

\subsection{\texorpdfstring{is\allowbreak{}Pre\allowbreak{}Brownian}{isPreBrownian}}


Combine centered Gaussian finite-dimensional laws with covariance \(\min(s,t)\). The Gaussian-process Brownian characterization yields \texttt{Is\allowbreak{}Pre\allowbreak{}Brownian\allowbreak{}Real}. This proves Brownian finite-dimensional laws; the prefix “pre” records that sample-path continuity is not part of this predicate.

\subsection{\texorpdfstring{Vector\allowbreak{}Sample}{VectorSample}}


For a coordinate index type \(I\), take one scalar Brownian sample for each \(i\in I\). This is the product sample type \(I\to\mathrm{Sample}\), allowing independent entire coordinate processes rather than only independent values at one time.

\subsection{\texorpdfstring{vector\allowbreak{}Probability}{vectorProbability}}


Equip the vector sample space with the infinite product of the scalar Brownian source probabilities. The coordinate processes inherit independence from this actual product construction.

\subsection{\texorpdfstring{coordinate}{coordinate}}


Define the \(i\)-th coordinate process by \(B_i(t,\omega)=B(t,\omega_i)\). It first projects to the \(i\)-th scalar sample, then evaluates the already constructed scalar process at \(t\).

\subsection{\texorpdfstring{coordinate is\allowbreak{}Pre\allowbreak{}Brownian}{coordinate isPreBrownian}}


The product coordinate projection has the scalar Brownian source as its law. Pull the scalar pre-Brownian process back along this projection; its finite-dimensional distributions, mean, and covariance are preserved. Thus each vector coordinate has the established scalar Brownian laws.

\subsection{\texorpdfstring{coordinate paths independent}{coordinate paths independent}}


The product coordinate projections are independent as sample-valued maps. Compose each projection with the map sending a scalar sample to its entire time-indexed process. Independence is preserved under these coordinatewise functions, giving independence of the whole coordinate paths.

\subsection{\texorpdfstring{coordinate grid paths independent}{coordinate grid paths independent}}


Restrict each independent coordinate path to the deterministic grid \(nh\). This is a coordinatewise function of the path, so independence is preserved. The result supplies exactly the countable grid-path hypothesis used in the innovation theorems.

\subsection{\texorpdfstring{instance at 89}{instance at 89}}


Each scalar Brownian sample source has mass one. The infinite product over vector coordinates therefore also has mass one. This instance concerns the actual vector sample measure, not an assumed normalization of an abstract process law.

\section{\texorpdfstring{The quadratic upper model}{The quadratic upper model}}
\label{proof-module-UnadjustedLangevin-Analysis-Smooth-lean}
{\small\color{gray}Source module: \texttt{Unadjusted\allowbreak{}Langevin/\allowbreak{}Analysis/\allowbreak{}Smooth.\allowbreak{}lean}.\par}


We next obtain the complementary upper bound from a Lipschitz gradient. This argument uses only first derivatives: restrict the potential to a segment, subtract its affine approximation and the proposed quadratic error, and show that the resulting scalar function decreases.



\texttt{smooth\_\allowbreak{}upper\_\allowbreak{}bound} concerns the gradient of the same differentiable potential occurring in its function values. The argument uses only first derivatives and a Lipschitz gradient; the scalar line restriction and its interval integral provide the quadratic remainder, with no Hessian assumption.

\subsection{\texorpdfstring{The descent lemma}{The descent lemma}}


Let \(H\) be a real Hilbert space and \(L\ge0\). Suppose \(f:H\to\mathbb R\) is differentiable and

\[\|\nabla f(u)-\nabla f(v)\|\le L\|u-v\|\quad(u,v\in H).\]

Then for every \(x,y\),

\[f(y)\le f(x)+\langle\nabla f(x),y-x\rangle+\frac L2\|y-x\|^2.\]

Set \(d=y-x\), \(a=\langle\nabla f(x),d\rangle\), \(b=\|d\|^2\), and introduce the corrected segment function

\[q(t)=f(x+td)-f(x)-ta-\frac L2t^2b.\]

It suffices to show \(q(1)\le q(0)=0\).



The segment has derivative \(d\). The chain rule therefore gives, for every real \(t\),

\[\frac{d}{dt}f(x+td)=\langle\nabla f(x+td),d\rangle.\]

Subtract the derivatives of the affine and quadratic terms:

\[q'(t)=\langle\nabla f(x+td),d\rangle-a-Ltb.\]

For \(t\) in the interior of \([0,1]\), we have \(t\ge0\). Lipschitz continuity gives

\[\|\nabla f(x+td)-\nabla f(x)\|\le L\|td\|=Lt\|d\|.\]

By Cauchy–Schwarz,

\[\langle\nabla f(x+td)-\nabla f(x),d\rangle\le\|\nabla f(x+td)-\nabla f(x)\|\,\|d\|\le Lt\|d\|^2.\]

Substituting \(a\) and \(b\) shows that \(q'(t)\le0\).



The derivative formula also makes \(q\) continuous on \([0,1]\) and differentiable in its interior. The derivative criterion for monotonicity now implies that \(q\) is nonincreasing on this interval.



Compare the two endpoints:

\[q(1)\le q(0),\qquad f(x+d)-f(x)-a-\frac L2b\le0.\]

Finally \(x+d=x+(y-x)=y\). Substitute the definitions of \(a,b,d\) and rearrange to obtain the quadratic upper bound.

\section{\texorpdfstring{The sharp deterministic contraction}{The sharp deterministic contraction}}
\label{proof-module-UnadjustedLangevin-SharpContraction-lean}
{\small\color{gray}Source module: \texttt{Unadjusted\allowbreak{}Langevin/\allowbreak{}Sharp\allowbreak{}Contraction.\allowbreak{}lean}.\par}


The preceding square-root factor uses strong monotonicity and Lipschitz continuity separately. The quadratic lower and upper models contain more information. We combine them into the interpolation inequality, then interpret that inequality as a ball containing the gradient difference. This yields the two sharp factors \(1-mh\) and \(Lh-1\), with the change of regime at \(h=2/(m+L)\).



The proof separates the two regimes for the maximum contraction factor \(\max(|1-mh|,|1-Lh|)\). Each scalar bound follows after the strong monotonicity and cocoercivity estimates have supplied the needed inner-product control. The final norm theorem removes the square by nonnegativity, retaining the precise step-regime hypotheses.

\subsection{\texorpdfstring{Cocoercivity from two quadratic models}{Cocoercivity from two quadratic models}}


Let \(H\) be a real inner-product space, \(f:H\to\mathbb R\), \(g:H\to H\), and \(K\ge0\). Suppose that for all \(a,b\in H\),

\[f(a)+\langle g(a),b-a\rangle\le f(b)\le f(a)+\langle g(a),b-a\rangle+\frac K2\|b-a\|^2.\]

Then for all \(x,y\),

\[\|g(x)-g(y)\|^2\le K\langle g(x)-g(y),x-y\rangle.\]

The proof includes \(K=0\), and does not separately assume that \(g\) is a gradient.



Fix \(a,b\), \(t\ge0\), and set \(v=g(a)-g(b)\), \(z=a-tv\). Compare the lower model based at \(b\) with the upper model based at \(a\), both evaluated at \(z\):

\[f(b)+\langle g(b),z-b\rangle\le f(z)\le f(a)+\langle g(a),z-a\rangle+\frac K2\|z-a\|^2.\]

The displacements are \(z-a=-tv\) and \(z-b=(a-b)-tv\). Substitute them into the two models. Since \(t\ge0\), the quadratic term is \(\|z-a\|^2=t^2\|v\|^2\).



The difference of the linear terms satisfies

\[\langle g(a),v\rangle-\langle g(b),v\rangle=\|v\|^2.\]

Eliminating \(f(z)\) therefore gives

\[\left(t-\frac K2t^2\right)\|g(a)-g(b)\|^2\le f(a)-f(b)-\langle g(b),a-b\rangle.\]

Apply this estimate to \((a,b)=(x,y)\) and \((y,x)\), with the same \(t\ge0\). Adding will cancel the function values.



The two gradient norms are equal and the reversed displacement changes the sign of its inner product. Thus, writing \(A=\|g(x)-g(y)\|^2\) and \(B=\langle g(x)-g(y),x-y\rangle\), we have for every \(t\ge0\),

\[(2t-Kt^2)A\le B.\]

First consider \(K=0\). Taking \(t=1\) gives \(2A\le B\). As \(A\ge0\), this also shows \(B\ge0\).



If the desired conclusion \(A\le0\) were false, then \(A>0\). We could choose

\[t=\frac{B+1}{2A}\ge0.\]

The estimate \(2tA\le B\) would then force \(B+1\le B\).



To justify that substitution, \(A>0\) implies \(\|g(x)-g(y)\|\ne0\), so its square can be cancelled from the denominator. Hence indeed

\[2\frac{B+1}{2A}A=B+1.\]

The contradiction \(B+1\le B\) excludes \(A>0\). Thus \(A=0\), proving the claim when \(K=0\).



Now suppose \(K>0\). Choose \(t=K^{-1}\) in the common estimate:

\[\frac AK\le B.\]

Multiplying by \(K\) gives \(A\le KB\), as required.

\subsection{\texorpdfstring{The strong-\allowbreak{}smooth interpolation inequality}{The strong-smooth interpolation inequality}}


Let \(H\) be a real Hilbert space, \(f:H\to\mathbb R\) differentiable and \(m\)-strongly convex, and let \(\nabla f\) be \(L\)-Lipschitz, where \(L\ge0\) and \(m\le L\). Then for every \(x,y\),

\[mL\|x-y\|^2+\|\nabla f(x)-\nabla f(y)\|^2\le(m+L)\langle\nabla f(x)-\nabla f(y),x-y\rangle.\]

The inequality holds also at the boundary \(m=L\).



Subtract the strong-convexity quadratic and its gradient:

\[F(v)=f(v)-\frac m2\|v\|^2,\qquad G(v)=\nabla f(v)-mv.\]

The quadratic lower model for \(f\) becomes an ordinary supporting plane for \(F\):

\[F(a)+\langle G(a),b-a\rangle\le F(b).\]

This follows by expanding the squared distance and collecting the terms involving \(m\).



Apply the smooth upper model to \(f\), and substitute the same definitions of \(F\) and \(G\).



Expanding the quadratic terms shows that the upper curvature decreases from \(L\) to \(L-m\):

\[F(b)\le F(a)+\langle G(a),b-a\rangle+\frac{L-m}{2}\|b-a\|^2.\]

Since \(L-m\ge0\), cocoercivity for these two models gives

\[\|G(x)-G(y)\|^2\le(L-m)\langle G(x)-G(y),x-y\rangle.\]

Write \(d=x-y\) and \(v=\nabla f(x)-\nabla f(y)\). The shifted gradient difference is exactly

\[G(x)-G(y)=v-md.\]

Expand the cocoercivity inequality using this identity:

\[\|v\|^2-2m\langle v,d\rangle+m^2\|d\|^2\le(L-m)\bigl(\langle v,d\rangle-m\langle d,d\rangle\bigr).\]

Use \(\langle d,d\rangle=\|d\|^2\) and collect terms. This gives

\[\|v\|^2+mL\|d\|^2\le(m+L)\langle v,d\rangle.\]

When \(m=L\), the cocoercivity constant above is zero, forcing \(G(x)=G(y)\); consequently \(\nabla f(x)-\nabla f(y)=m(x-y)\). Thus no division by \(L-m\) or limiting argument is needed.

\subsection{\texorpdfstring{The gradient difference lies in a ball}{The gradient difference lies in a ball}}


Let \(H\) be a real inner-product space, \(x,y,g_x,g_y\in H\), and \(m,L\in\mathbb R\) with \(m\le L\). Suppose

\[mL\|x-y\|^2+\|g_x-g_y\|^2\le(m+L)\langle g_x-g_y,x-y\rangle.\]

Then

\[\left\|(g_x-g_y)-\frac{m+L}{2}(x-y)\right\|\le\frac{L-m}{2}\|x-y\|.\]

The proposed radius is nonnegative because \(m\le L\). Both sides of the desired norm bound are therefore nonnegative, so it is enough to prove the corresponding squared inequality.



Set \(c=(m+L)/2\), \(r=(L-m)/2\), \(d=x-y\), and \(v=g_x-g_y\). Completing the square in the interpolation inequality gives

\[\begin{aligned}\|v-cd\|^2&=\|v\|^2-(m+L)\langle v,d\rangle+c^2\|d\|^2\\&\le(c^2-mL)\|d\|^2=r^2\|d\|^2.\end{aligned}\]

Taking nonnegative square roots proves the claim. If \(m=L\), the radius is zero and \(v=md\).

\subsection{\texorpdfstring{The centered Euler-\allowbreak{}step bound}{The centered Euler-step bound}}


Let \(H\) be a real inner-product space, \(m,L,h\in\mathbb R\) with \(m\le L\) and \(h\ge0\), and let \(x,y,g_x,g_y\in H\) satisfy

\[mL\|x-y\|^2+\|g_x-g_y\|^2\le(m+L)\langle g_x-g_y,x-y\rangle.\]

Then

\[\|(x-hg_x)-(y-hg_y)\|\le\left(\left|1-h\frac{m+L}{2}\right|+h\frac{L-m}{2}\right)\|x-y\|.\]

The preceding ball bound gives \(\|v-cd\|\le r\|d\|\), where \(d=x-y\), \(v=g_x-g_y\), \(c=(m+L)/2\), and \(r=(L-m)/2\).



Decompose the Euler difference into its centered part and its deviation:

\[d-hv=(1-hc)d-h(v-cd).\]

The triangle inequality and homogeneity of the norm give

\[\|d-hv\|\le|1-hc|\,\|d\|+h\|v-cd\|,\]

where \(|-h|=h\) because \(h\ge0\).



Multiply the ball bound by \(h\ge0\) and substitute it:

\[\|d-hv\|\le\bigl(|1-hc|+hr\bigr)\|d\|.\]

For \(m=L\), this reduces to the exact scalar factor \(|1-mh|\).

\subsection{\texorpdfstring{The sharp small-\allowbreak{}step factor}{The sharp small-step factor}}


Let \(H\) be a real inner-product space, \(x,y,g_x,g_y\in H\), and \(m,L,h\in\mathbb R\). Suppose \(m\le L\), \(m+L>0\), \(0\le h\le2/(m+L)\), and

\[mL\|x-y\|^2+\|g_x-g_y\|^2\le(m+L)\langle g_x-g_y,x-y\rangle.\]

Then

\[\|(x-hg_x)-(y-hg_y)\|\le(1-mh)\|x-y\|.\]

Apply the centered bound. The step-size assumption and \(m+L>0\) imply \(h(m+L)\le2\), hence \(1-h(m+L)/2\ge0\).



Remove the absolute value with this sign and simplify its coefficient:

\[\left|1-h\frac{m+L}{2}\right|+h\frac{L-m}{2}=1-h\frac{m+L}{2}+h\frac{L-m}{2}=1-mh.\]\subsection{\texorpdfstring{The sharp large-\allowbreak{}step factor}{The sharp large-step factor}}


Let \(H\) be a real inner-product space, \(x,y,g_x,g_y\in H\), and \(m,L,h\in\mathbb R\). Suppose \(m\le L\), \(m+L>0\), \(h\ge0\), \(2/(m+L)\le h\), and

\[mL\|x-y\|^2+\|g_x-g_y\|^2\le(m+L)\langle g_x-g_y,x-y\rangle.\]

Then

\[\|(x-hg_x)-(y-hg_y)\|\le(Lh-1)\|x-y\|.\]

This is a distance bound; it is a strict contraction only when its coefficient is below one.



Apply the centered estimate. Now the step-size assumption gives \(2\le h(m+L)\), so \(1-h(m+L)/2\le0\).



Replace the absolute value by the negative of its argument. The coefficient becomes

\[\left|1-h\frac{m+L}{2}\right|+h\frac{L-m}{2}=h\frac{m+L}{2}-1+h\frac{L-m}{2}=Lh-1.\]\subsection{\texorpdfstring{Small-\allowbreak{}step contraction for a potential}{Small-step contraction for a potential}}


Let \(f:H\to\mathbb R\) be differentiable and \(m\)-strongly convex on a real Hilbert space, with an \(L\)-Lipschitz gradient, \(L\ge0\), \(m\le L\), and \(m+L>0\). If \(0\le h\le2/(m+L)\), then for every \(x,y\),

\[\|(x-h\nabla f(x))-(y-h\nabla f(y))\|\le(1-mh)\|x-y\|.\]

The quadratic models imply the strong-smooth interpolation inequality for \(g_x=\nabla f(x)\) and \(g_y=\nabla f(y)\). Apply the small-step Euler bound to these two vectors. Its hypotheses are exactly the stated curvature and step-size assumptions.

\subsection{\texorpdfstring{Large-\allowbreak{}step contraction for a potential}{Large-step contraction for a potential}}


Let \(H\) be a real Hilbert space and \(f:H\to\mathbb R\) differentiable and \(m\)-strongly convex, with an \(L\)-Lipschitz gradient. Suppose \(L\ge0\), \(m\le L\), \(m+L>0\), \(h\ge0\), and \(2/(m+L)\le h\). Then for every \(x,y\in H\),

\[\|(x-h\nabla f(x))-(y-h\nabla f(y))\|\le(Lh-1)\|x-y\|.\]

Use the potential's interpolation inequality in the large-step Euler bound. This gives the second sharp factor \(Lh-1\).

\subsection{\texorpdfstring{The endpoint-\allowbreak{}factor form}{The endpoint-factor form}}


Let \(f\) be differentiable and \(m\)-strongly convex on a real Hilbert space, with an \(L\)-Lipschitz gradient, \(L\ge0\), \(m\le L\), and \(m+L>0\). For every \(h\ge0\) and every \(x,y\),

\[\|(x-h\nabla f(x))-(y-h\nabla f(y))\|\le\max\{|1-mh|,|1-Lh|\}\,\|x-y\|.\]

If \(0<m\le L\), this factor is below one for \(0<h<2/L\). At \(m=L\), the two endpoint factors coincide.



Split according to whether \(h\le2/(m+L)\). In the first case the small-step theorem bounds the distance by \((1-mh)\|x-y\|\).



Since \(1-mh\le|1-mh|\le\max\{|1-mh|,|1-Lh|\}\), multiplying by \(\|x-y\|\ge0\) gives the desired bound in this case.



In the other case \(h>2/(m+L)\), so the large-step theorem bounds the distance by \((Lh-1)\|x-y\|\).



The elementary absolute-value inequality yields

\[Lh-1=-(1-Lh)\le|1-Lh|.\]

Finally \(|1-Lh|\le\max\{|1-mh|,|1-Lh|\}\). Multiply by the nonnegative distance to finish the second case.

\subsection{\texorpdfstring{Sharp shared-\allowbreak{}noise contraction: small steps}{Sharp shared-noise contraction: small steps}}


Let \(f:E_d\to\mathbb R\) be differentiable and \(m\)-strongly convex, with an \(L\)-Lipschitz gradient, \(L\ge0\), \(m\le L\), and \(m+L>0\). If \(0\le h\le2/(m+L)\), then for all starting states \(x,y\) and every common noise vector \(z\),

\[\|T_h^z(x)-T_h^z(y)\|\le(1-mh)\|x-y\|.\]

The common noise cancels exactly, leaving the distance between the gradient steps from \(x\) and \(y\).



Apply the potential-based small-step estimate. It gives the factor \(1-mh\) independently of the value or distribution of \(z\).

\subsection{\texorpdfstring{Sharp shared-\allowbreak{}noise contraction: large steps}{Sharp shared-noise contraction: large steps}}


Let \(f:E_d\to\mathbb R\) be differentiable and \(m\)-strongly convex, with an \(L\)-Lipschitz gradient, \(L\ge0\), \(m\le L\), and \(m+L>0\). If \(h\ge0\) and \(2/(m+L)\le h\), then for every \(x,y,z\in E_d\),

\[\|T_h^z(x)-T_h^z(y)\|\le(Lh-1)\|x-y\|.\]

Cancel the shared noise, reducing the two ULA updates to their deterministic gradient-step difference.



Apply the potential-based large-step estimate. This yields the factor \(Lh-1\) for every common noise vector.

\section{\texorpdfstring{The local Euler defect}{The local Euler defect}}
\label{proof-module-UnadjustedLangevin-LocalError-lean}
{\small\color{gray}Source module: \texttt{Unadjusted\allowbreak{}Langevin/\allowbreak{}Local\allowbreak{}Error.\allowbreak{}lean}.\par}


Let \(H\) be a real Banach space, with norm \(\|\cdot\|\). All Lipschitz constants \(L\) below are nonnegative. All time integrals are Bochner integrals. For the Langevin application we will take \(H=L^2(P;\mathbb R^d)\); the norms below then measure random variables in mean square. The scalar defect estimate permits a nonnegative real parameter \(d\); its application uses the natural-number dimension, viewed as a real number. Begin with a Lipschitz drift \(g\), a reference curve \(Y\), and its additive noise \(B\). We first derive the error of freezing the drift over one time step directly from the integral equation.



The analytic calculation is conditional on the displayed integral equation and norm bounds. It neither constructs a stationary Langevin diffusion nor proves its score moment. The Gaussian and Brownian noise moments used later are proved separately.



\texttt{local\allowbreak{}Defect} is a vector-valued Bochner interval integral. It exists as a totalized expression before integrability is proved; all norm and linearity steps below explicitly establish the continuous or integrable hypotheses needed for their ordinary mathematical meaning. The interval \(0..h\) is oriented, while the error estimates assume \(h\ge0\).

\subsection{\texorpdfstring{Definition of the local defect}{Definition of the local defect}}


For \(g:H\to H\), \(Y:\mathbb R\to H\), and \(h\in\mathbb R\), define the drift omitted when Euler freezes \(g\) at the initial point:

\[D_h(g,Y):=\int_0^h\bigl(g(Y_t)-g(Y_0)\bigr)\,dt.\]

Integrate the difference between the evolving drift and its frozen value. This defines the defect itself; the estimates below will be consequences of this definition.

\subsection{\texorpdfstring{Separate the exact and frozen drift}{Separate the exact and frozen drift}}


If \(t\mapsto g(Y_t)\) is integrable on the interval between \(0\) and \(h\), then

\[D_h(g,Y)=\int_0^h g(Y_t)\,dt-h\,g(Y_0).\]

Expand the definition. Both the actual drift and the constant function \(g(Y_0)\) are integrable, so linearity separates their difference into two integrals.



The integral of the constant vector is \((h-0)g(Y_0)=h\,g(Y_0)\), yielding the identity.

\subsection{\texorpdfstring{An exact step is Euler minus its defect}{An exact step is Euler minus its defect}}


Assume the preceding drift integrability and the endpoint integral equation

\[Y_h-Y_0=-\int_0^h g(Y_t)\,dt+B_h.\]

Then the exact reference step has the form

\[Y_h=Y_0-h\,g(Y_0)+B_h-D_h(g,Y).\]

Solve the assumed integral equation for \(Y_h\), then replace the defect by the difference between the exact drift integral and its frozen approximation.



Cancel the two frozen-drift terms. The resulting equality is precisely the Euler update with the signed remainder \(-D_h\).

\subsection{\texorpdfstring{Continuity supplies drift integrability}{Continuity supplies drift integrability}}


Let \(H\) be a normed additive group, \(g:H\to H\) be \(L\)-Lipschitz with \(L\ge0\), and \(Y:\mathbb R\to H\) be continuous. For every real \(h\), the function \(t\mapsto g(Y_t)\) is interval integrable from \(0\) to \(h\).



A Lipschitz map is continuous, so \(g\circ Y\) is continuous. A continuous function is integrable on every compact interval, proving the claim.

\subsection{\texorpdfstring{The square-\allowbreak{}root integral}{The square-root integral}}


For every \(h\ge0\),

\[\int_0^h\sqrt t\,dt=\frac23h\sqrt h.\]

Write \(\sqrt t=t^{1/2}\) and use the power-integral formula. The exponent satisfies \(1/2>-1\), so the formula applies at the zero endpoint.



The contribution at zero vanishes. Since \(h\ge0\), the other endpoint is \(h^{1/2+1}=h^{1/2}h=\sqrt h\,h\).



Divide by \(1/2+1=3/2\) and rearrange the factors to obtain \((2/3)h\sqrt h\).

\subsection{\texorpdfstring{Bound a reference increment}{Bound a reference increment}}


Let \(h\ge0\) and suppose that, for every \(0\le s\le h\),

\[\|g(Y_s)\|\le G,\qquad \|B_s\|\le N\sqrt s,\qquad Y_s-Y_0=-\int_0^s g(Y_r)\,dr+B_s.\]

Then, for each \(0\le t\le h\),

\[\|Y_t-Y_0\|\le tG+N\sqrt t.\]

First bound the drift integral on \([0,t]\) by its constant norm bound:

\[\left\|\int_0^t g(Y_s)\,ds\right\|\le G\,|t-0|=Gt,\]

where \(t\ge0\) removes the absolute value.



To justify the constant bound under that integral, note that \(s\in(0,t]\) implies \(0\le s\le t\le h\). Hence the assumed bound on \([0,h]\) applies to every such \(s\).



Insert the integral equation and apply the triangle inequality:

\[\|Y_t-Y_0\|\le\left\|\int_0^t g(Y_s)\,ds\right\|+\|B_t\|.\]

Negation leaves the norm of the drift integral unchanged.



Apply the two bounds to obtain \(Gt+N\sqrt t\), and commute the factors \(G\) and \(t\).

\subsection{\texorpdfstring{Integrate the Lipschitz increment bound}{Integrate the Lipschitz increment bound}}


Assume \(h\ge0\), \(g\) is \(L\)-Lipschitz with \(L\ge0\), and for \(0\le t\le h\) the same integral equation and bounds \(\|g(Y_t)\|\le G\), \(\|B_t\|\le N\sqrt t\) hold. Then

\[\|D_h(g,Y)\|\le L\left(\frac{h^2G}{2}+\frac23h\sqrt h\,N\right).\]

For \(0<t\le h\), Lipschitz continuity gives

\[\|g(Y_t)-g(Y_0)\|\le L\|Y_t-Y_0\|.\]

Multiply the increment estimate by \(L\ge0\) to obtain the pointwise majorant

\[\|g(Y_t)-g(Y_0)\|\le L(tG+N\sqrt t).\]

The scalar majorant \(t\mapsto L(tG+N\sqrt t)\) is continuous, hence integrable on \([0,h]\).



Expand \(D_h\). The norm estimate for a Bochner integral, together with the integrable majorant, gives

\[\|D_h\|\le\int_0^h L(tG+N\sqrt t)\,dt.\]

Pull out constants and split the two integrable summands. Evaluate

\[\int_0^h t\,dt=\frac{h^2}{2},\qquad \int_0^h\sqrt t\,dt=\frac23h\sqrt h.\]

This leaves the two drift and noise contributions in the claimed expression.



Rearrange the scalar factors to place the drift term first and the noise term second.

\subsection{\texorpdfstring{The exact rational constant}{The exact rational constant}}


The numerical factor in the corrected Dalalyan–Karagulyan estimate satisfies

\[\frac76\sqrt2\le\frac{33}{20}=1.65.\]

Use \((\sqrt2)^2=2\) and \(\sqrt2\ge0\). Comparing squares proves the bound, since \(49/18\le1089/400\).

\subsection{\texorpdfstring{The local error with stationary moment scales}{The local error with stationary moment scales}}


Let \(h,d\in\mathbb R\) with \(h,d\ge0\), let \(g\) be \(L\)-Lipschitz with \(L\ge0\), and suppose \(Lh\le2\). Assume that, for every \(0\le t\le h\),

\[Y_t-Y_0=-\int_0^t g(Y_s)\,ds+B_t,\quad \|g(Y_t)\|\le\sqrt{Ld},\quad \|B_t\|\le\sqrt{2d}\sqrt t.\]

Then

\[\boxed{\|D_h(g,Y)\|\le\frac{33}{20}Lh\sqrt{hd}.}\]

In \(L^2\), the drift bound corresponds to a score second moment at most \(Ld\); here that bound remains a hypothesis.



Compare the nonnegative quantities \(h\sqrt{Ld}\) and \(\sqrt2\sqrt{hd}\) by squaring. It suffices to show \(h^2Ld\le2hd\).



Replace each squared square root by its nonnegative argument. Multiply \(Lh\le2\) by \(hd\ge0\), obtaining \(h^2Ld\le2hd\), and thus

\[h\sqrt{Ld}\le\sqrt2\sqrt{hd}.\]

Split products under the square roots and rearrange:

\[h\sqrt h\sqrt{2d}=h\sqrt2\sqrt{hd}.\]

The nonnegativity of \(h\) and \(2\) justifies the square-root product identities.



Apply the two-term local-defect bound with \(G=\sqrt{Ld}\) and \(N=\sqrt{2d}\):

\[\|D_h\|\le L\left(\frac{h^2\sqrt{Ld}}2+\frac23h\sqrt h\sqrt{2d}\right).\]

Bound the first term using \(h\sqrt{Ld}\le\sqrt2\sqrt{hd}\), and rewrite the second with the preceding identity. Since \(1/2+2/3=7/6\),

\[\|D_h\|\le L\frac76\sqrt2\,h\sqrt{hd}.\]

Multiply \((7/6)\sqrt2\le33/20\) by the nonnegative factor \(Lh\sqrt{hd}\). This gives the claimed constant without numerical rounding.

\section{\texorpdfstring{The noise-\allowbreak{}driven iterates}{The noise-driven iterates}}
\label{proof-module-UnadjustedLangevin-Dynamics-lean}
{\small\color{gray}Source module: \texttt{Unadjusted\allowbreak{}Langevin/\allowbreak{}Dynamics.\allowbreak{}lean}.\par}


Write \(T_h(x)=x-h\nabla U(x)\), so one Euler update is \(T_h(x)+\sqrt{2h}\,z\). The noise vectors remain explicit. Common-noise algebra works for any supplied noise sequence; identifying the algorithm with Gaussian ULA also requires independent standard Gaussian noise hypotheses.



\texttt{ula\allowbreak{}Chain} recursively applies \texttt{ula\allowbreak{}Update} to a supplied sequence of vectors. Its zero and successor equations are definitional. \texttt{ula\allowbreak{}Chain\_\allowbreak{}measurable} proves measurable dependence on the sample by induction, combining the current iterate, the Lipschitz gradient's measurability, and the next supplied measurable innovation. No independence or Gaussian-law assumption is needed for this measurability assertion.

\subsection{\texorpdfstring{Definition of the chain}{Definition of the chain}}


Given \(U:\mathbb R^d\to\mathbb R\), \(h\in\mathbb R\), an initial vector \(x_0\), and noise \((\xi_k)\), define

\[X_0=x_0,\qquad X_{k+1}=X_k-h\nabla U(X_k)+\sqrt{2h}\,\xi_k.\]

Start the recursion at the supplied initial vector: \(X_0=x_0\).



Apply the Euler update to \(X_k\) using the next noise vector \(\xi_k\). This determines \(X_{k+1}\).

\subsection{\texorpdfstring{The initial iterate}{The initial iterate}}


For every potential, step, initial vector, and noise sequence, the recursively defined chain satisfies \(X_0=x_0\).



The equality is exactly the initial clause of the recursive definition.

\subsection{\texorpdfstring{The update identity}{The update identity}}


For every \(k\ge0\), the chain satisfies \(X_{k+1}=\operatorname{ulaUpdate}(U,h,X_k,\xi_k)\).



The successor clause of the recursion gives the asserted update formula.

\subsection{\texorpdfstring{Measurability of every iterate}{Measurability of every iterate}}


Let \(\Omega\) be a measurable space. If \(\nabla U\) is \(L\)-Lipschitz, the initial vector \(X_0:\Omega\to\mathbb R^d\) is measurable, and every \(\xi_k:\Omega\to\mathbb R^d\) is measurable, then \(X_n\) is measurable for every \(n\).



A Lipschitz gradient is continuous and therefore measurable.



Induct on \(n\). The initial iterate is measurable by hypothesis. For the successor step, assume \(X_n\) is measurable.



The composition \(\nabla U\circ X_n\) is measurable. Subtracting its scalar multiple from \(X_n\), then adding \(\sqrt{2h}\,\xi_n\), preserves measurability. Hence \(X_{n+1}\) is measurable.

\section{\texorpdfstring{The actual gradient on square-\allowbreak{}integrable random variables}{The actual gradient on square-integrable random variables}}
\label{proof-module-UnadjustedLangevin-LpDrift-lean}
{\small\color{gray}Source module: \texttt{Unadjusted\allowbreak{}Langevin/\allowbreak{}Lp\allowbreak{}Drift.\allowbreak{}lean}.\par}


Let \(P\) be a finite measure on \(\Omega\). The first five results need only a normed additive commutative group \(H\); all Lipschitz constants \(L\) are nonnegative. When \(H\) is a real Banach space, \(L^2(P;H)\) provides the Banach space needed for the local-defect estimate. For a Lipschitz vector field \(g:H\to H\), construct the operator \(\mathcal G[X]=[g\circ X]\), including the case \(g(0)\ne0\). Finiteness of \(P\) ensures the constant offset is square integrable.



The gradient results use a real Hilbert space \(H\). The two score-moment lemmas need only a real inner-product space, without completeness. Whenever \(g=\nabla U\), this is the potential’s actual gradient, and the resulting \(L^2\) recursion is proved to represent the original pointwise ULA chain.



\texttt{lp\allowbreak{}Drift} lifts a Lipschitz pointwise map to \(L^2\), using its linear growth and the finite probability measure to prove integrability. Its representative theorem identifies the lifted map almost everywhere with ordinary composition. All subsequent Hilbert-space estimates concern these quotient classes and are invariant under representative choices.

\subsection{\texorpdfstring{Center a Lipschitz vector field}{Center a Lipschitz vector field}}


In a normed additive group \(H\), if \(g:H\to H\) is \(L\)-Lipschitz, then so is

\[\widetilde g(x):=g(x)-g(0).\]

For arbitrary \(x,y\), the common offset cancels:

\[\|\widetilde g(x)-\widetilde g(y)\|=\|g(x)-g(y)\|\le L\|x-y\|.\]

This is precisely the required Lipschitz inequality.

\subsection{\texorpdfstring{Construct the lifted drift}{Construct the lifted drift}}


Let \(P\) be finite and \(g\) be \(L\)-Lipschitz. For \(X\in L^2(P;H)\), define

\[\mathcal G[X]:=[\widetilde g\circ X]+[g(0)].\]

The centered function satisfies \(\widetilde g(0)=0\) and \(\|\widetilde g(x)\|\le L\|x\|\), so composition produces an \(L^2\) class. Add the constant class \([g(0)]\), which belongs to \(L^2\) because \(P\) is finite.

\subsection{\texorpdfstring{The lift agrees with the original vector field}{The lift agrees with the original vector field}}


For finite \(P\), Lipschitz \(g\), and \(X\in L^2(P;H)\),

\[\mathcal G[X](\omega)=g(X(\omega))\quad\text{for }P\text{-almost every }\omega.\]

Expand \(\mathcal G\) into the centered composition class plus the constant class.



Intersect the full-measure sets on which the representative of the sum, the composition, and the constant class agree with their pointwise formulas.



On this set the value is \((g(X(\omega))-g(0))+g(0)\). Cancel the offset to recover \(g(X(\omega))\).

\subsection{\texorpdfstring{The actual pointwise drift is square integrable}{The actual pointwise drift is square integrable}}


For finite \(P\), Lipschitz \(g\), and \(X\in L^2(P;H)\),

\[g\circ X\in L^2(P;H).\]

The constructed class \(\mathcal G[X]\) has an \(L^2\) representative. Transfer its square integrability to \(g\circ X\) using their almost-everywhere equality.

\subsection{\texorpdfstring{The drift keeps its Lipschitz constant in L\ensuremath{{}^{2}}}{The drift keeps its Lipschitz constant in L²}}


For finite \(P\) and \(L\)-Lipschitz \(g\), the lifted drift is \(L\)-Lipschitz:

\[\|\mathcal G[X]-\mathcal G[Y]\|_{L^2}\le L\|X-Y\|_{L^2}.\]

Fix \(X,Y\). The two constant classes cancel in their drift difference. For the remaining centered compositions, the pointwise Lipschitz estimate passes to the \(L^2\) norm, giving the same constant \(L\).

\subsection{\texorpdfstring{Lift the small-\allowbreak{}step gradient contraction}{Lift the small-step gradient contraction}}


Let \(H\) be a real Hilbert space, \(P\) finite, and \(U:H\to\mathbb R\) differentiable and \(m\)-strongly convex on \(H\), with \(L\)-Lipschitz gradient. Assume \(m\le L\), \(m+L>0\), and \(0\le h\le2/(m+L)\). Write \(\mathcal G[X]=[\nabla U\circ X]\). Then, for all \(X,Y\in L^2(P;H)\),

\[\|(X-h\mathcal G[X])-(Y-h\mathcal G[Y])\|_{L^2}\le(1-mh)\|X-Y\|_{L^2}.\]

It suffices to prove the corresponding norm inequality almost everywhere between representatives: integrating the squared pointwise inequality gives the \(L^2\) bound.



Work on a common full-measure set where the outer difference and both Euler-step differences have their ordinary pointwise subtraction formulas.



Also require the scalar multiples, both lifted gradients, and \(X-Y\) to have their pointwise formulas. The two drift terms are then \(\nabla U(X(\omega))\) and \(\nabla U(Y(\omega))\).



Substitute these representative identities and apply the deterministic sharp contraction to \(x=X(\omega)\), \(y=Y(\omega)\):

\[\|(x-h\nabla U(x))-(y-h\nabla U(y))\|\le(1-mh)\|x-y\|.\]\subsection{\texorpdfstring{Lift the second step-\allowbreak{}size regime}{Lift the second step-size regime}}


Let \(H\) be a real Hilbert space, \(P\) finite, and \(U\) differentiable, \(m\)-strongly convex, with \(L\)-Lipschitz gradient. Assume \(m\le L\), \(m+L>0\), \(h\ge0\), and \(h\ge2/(m+L)\). Then

\[\|(X-h\mathcal G[X])-(Y-h\mathcal G[Y])\|_{L^2}\le(Lh-1)\|X-Y\|_{L^2}.\]

The inequality holds in this whole regime; it is a strict contraction when its factor is less than one.



Reduce the \(L^2\) inequality to its almost-everywhere pointwise counterpart with factor \(Lh-1\).



Choose a full-measure set on which the outer subtraction and each Euler-step subtraction agree with the pointwise vector formulas.



On the same set identify scalar multiplication, both lifted gradient representatives, and the representative of \(X-Y\).



After substitution, invoke the deterministic second-regime estimate at \(X(\omega),Y(\omega)\):

\[\|(x-h\nabla U(x))-(y-h\nabla U(y))\|\le(Lh-1)\|x-y\|.\]\subsection{\texorpdfstring{The local defect for the actual L\ensuremath{{}^{2}} gradient}{The local defect for the actual L² gradient}}


Let \(P\) be finite, \(H\) a real Hilbert space, and \(\nabla U\) be \(L\)-Lipschitz. Let \(Y,B:\mathbb R\to L^2(P;H)\), let \(h,d\in\mathbb R\), and assume \(h,d\ge0\), \(Lh\le2\). For every \(0\le t\le h\), suppose

\[\|\mathcal G[Y_t]\|_{L^2}\le\sqrt{Ld},\quad\|B_t\|_{L^2}\le\sqrt{2d}\sqrt t,\quad Y_t-Y_0=-\int_0^t\mathcal G[Y_s]\,ds+B_t.\]

Then

\[\|D_h(\mathcal G,Y)\|_{L^2}\le\frac{33}{20}Lh\sqrt{hd}.\]

Apply the Banach-space local-defect theorem to \(\mathcal G\). Its required Lipschitz property is the one just proved for the actual lifted gradient; the reference equation and moment-scale bounds are the displayed hypotheses.

\subsection{\texorpdfstring{Identify the drift norm with the score moment}{Identify the drift norm with the score moment}}


For finite \(P\), a real inner-product space \(H\), Lipschitz \(g:H\to H\), and \(X\in L^2(P;H)\),

\[\|\mathcal G[X]\|_{L^2}^2=\int_\Omega\|g(X(\omega))\|^2\,dP(\omega).\]

Write the squared \(L^2\) norm as a self-inner-product and expand the \(L^2\) inner product into the integral of the pointwise inner product.



Almost everywhere, \(\mathcal G[X](\omega)=g(X(\omega))\). Replace the representative under the integral and use \(\langle v,v\rangle=\|v\|^2\).

\subsection{\texorpdfstring{A second-\allowbreak{}moment bound gives the drift bound}{A second-moment bound gives the drift bound}}


Under the preceding finite-measure and inner-product assumptions, if

\[\int_\Omega\|g(X(\omega))\|^2\,dP(\omega)\le C,\]

then

\[\|\mathcal G[X]\|_{L^2}\le\sqrt C.\]

For \(g=\nabla U\) and \(C=Ld\), the moment inequality is an explicit input, not a consequence of this lemma.



The norm-square identity turns the assumed moment bound into \(\|\mathcal G[X]\|^2\le C\). Taking the nonnegative square root yields the conclusion.

\subsection{\texorpdfstring{Define one ULA update in L\ensuremath{{}^{2}}}{Define one ULA update in L²}}


For finite \(P\), a real Hilbert space \(H\), a Lipschitz gradient \(\nabla U\), step \(h\in\mathbb R\), and \(X,Z\in L^2(P;H)\), define

\[\mathsf T_h(X,Z):=X-h\mathcal G[X]+\sqrt{2h}\,Z.\]

Form the drift step and add the scaled noise inside \(L^2\). All terms are already equivalence classes of square-integrable variables, so the vector-space operations produce another such class.

\subsection{\texorpdfstring{The L\ensuremath{{}^{2}} update is the pointwise ULA formula}{The L² update is the pointwise ULA formula}}


For finite \(P\), a real Hilbert space \(H\), Lipschitz \(\nabla U\), and arbitrary \(X,Z\in L^2(P;H)\),

\[\mathsf T_h(X,Z)(\omega)=X(\omega)-h\nabla U(X(\omega))+\sqrt{2h}\,Z(\omega)\quad\text{a.e.}\]

This identity needs no Gaussian or independence assumption on \(Z\).



Expand the definition of \(\mathsf T_h\) into its subtraction, drift, and scaled-noise terms.



Intersect the full-measure sets on which addition, subtraction, both scalar multiples, and the lifted gradient agree with their pointwise formulas.



Substitute these identities. The representative is exactly \(X(\omega)-h\nabla U(X(\omega))+\sqrt{2h}\,Z(\omega)\).

\subsection{\texorpdfstring{Define the genuine L\ensuremath{{}^{2}} ULA chain}{Define the genuine L² ULA chain}}


For a finite measure \(P\), a Lipschitz gradient on \(E_d\), an initial class \(X_0\in L^2(P;E_d)\), and noise classes \(Z_n\in L^2(P;E_d)\), define recursively, with the supplied initial class at index zero,

\[X_{n+1}:=\mathsf T_h(X_n,Z_n).\]

The initial class is the supplied \(X_0\). At the successor index, apply the \(L^2\) ULA update to the preceding iterate and the next noise class \(Z_n\).

\subsection{\texorpdfstring{The initial iterate}{The initial iterate}}


The recursively defined chain satisfies \(X_0=\mathrm{initial}\).



This is the zero branch of the recursion, so both sides reduce to the same class.

\subsection{\texorpdfstring{The successor iterate}{The successor iterate}}


For every \(n\ge0\), the same finite-measure, Lipschitz-gradient, initial-class and noise-class data give

\[X_{n+1}=\mathsf T_h(X_n,Z_n).\]

This is the successor branch of the defining recursion; unfold that branch to obtain the equality.

\subsection{\texorpdfstring{Identify the chain with the original algorithm}{Identify the chain with the original algorithm}}


For finite \(P\), Lipschitz gradient \(\nabla U:E_d\to E_d\), square-integrable initial and noise classes, and every \(n\ge0\),

\[X_n(\omega)=\operatorname{ulaChain}\bigl(U,h,X_0(\omega),(Z_k(\omega))_{k\ge0},n\bigr)\quad\text{for }P\text{-a.e. }\omega.\]

Induct on \(n\). At \(n=0\), both recursions return the same initial representative, so equality holds for every outcome. For the successor case, assume the representatives agree at index \(n\).



On a common full-measure set, combine the one-step representative identity with the induction hypothesis. Replace \(X_n(\omega)\) by the pointwise chain’s preceding iterate and apply the two successor recursion equations. Both next iterates are the same ULA update.

\section{\texorpdfstring{Existence of the additive Langevin integral equation}{Existence of the additive Langevin integral equation}}
\label{proof-module-UnadjustedLangevin-LangevinIntegral-lean}
{\small\color{gray}Source module: \texttt{Unadjusted\allowbreak{}Langevin/\allowbreak{}Langevin\allowbreak{}Integral.\allowbreak{}lean}.\par}


In a complete real normed space, a globally Lipschitz drift and a continuous additive path \(B\), \(B(0)=0\), determine a continuous solution of \(Y_t-x=-\int_0^tg(Y_s)ds+B_t\) on every finite nonnegative horizon. The construction is deterministic and later applies with the Banach space \(L^2(P)\).

\subsection{\texorpdfstring{Solve on a short interval}{Solve on a short interval}}


Remove the additive input by writing \(Y(t)=Z(t)+B(t)\). The new unknown satisfies \(Z(t)=x+\int_0^t f(s,Z(s))ds\), with \(f(s,z)=-g(z+B(s))\). This time-dependent drift is \(L\)-Lipschitz in \(z\) and continuous in time even though \(B\) need not be differentiable.



Bound \(\|B(t)\|\le C'\) on the compact interval and set \(G=\|g(x)\|+LC'\), radius \(a=2hG\), and drift bound \(M=G+La\). The condition \(Lh\le1/2\) proves \(Mh\le a\), giving the ball-invariance condition in \texttt{Is\allowbreak{}Picard\allowbreak{}Lindelof}. The proof checks these bounds and invokes \texttt{ODE.\allowbreak{}Fun\allowbreak{}Space.\allowbreak{}exists\_\allowbreak{}is\allowbreak{}Fixed\allowbreak{}Pt\_\allowbreak{}next}.



The chosen fixed point gives a continuous \(Z\) with its integral identity. Add back \(B\), use \(B(0)=0\), and rewrite the negative drift integral to obtain \(Y(t)-x=-\int_0^tg(Y(s))ds+B(t)\). The solution is represented as a real-time function, with the required identity asserted on \([0,h]\).

\subsection{\texorpdfstring{splice langevin integral}{splice langevin integral}}


Use \(Y(t)\) up to time \(a\), and \(Z(t-a)\) afterward. The equality \(Z(0)=Y(a)\) makes the pasted path continuous at the join. Split \(\int_0^t=\int_0^a+\int_a^t\), translate the second integral, and add the two supplied equations; the terms \(B(a)\) cancel, giving the same equation on \([0,a+b]\).

\subsection{\texorpdfstring{exists langevin integral}{exists langevin integral}}


Choose a positive segment length small enough that \(L\delta\le1/2\). Construct the first Picard segment, then repeatedly solve with initial state at the previous endpoint and shifted input \(B(a+s)-B(a)\). The splicing lemma combines finitely many segments covering the prescribed horizon. The local contraction restriction determines the construction mesh and imposes no restriction on the final horizon.

\subsection{\texorpdfstring{exists lp langevin integral small}{exists lp langevin integral small}}


Apply the Banach-space Picard theorem to \(H=L^2(P;K)\), \(g(X)=\operatorname{lpDrift}(\nabla U)(X)\). The lifted drift has the same Lipschitz constant as \(\nabla U\). Completeness of \(L^2\) and continuity of the supplied \(L^2\)-noise path are the exact fixed-point hypotheses.

\subsection{\texorpdfstring{exists lp langevin integral}{exists lp langevin integral}}


Use the finite-horizon Banach-space existence theorem with the lifted gradient drift. The conclusion is continuity in the \(L^2\) norm and an equality of \(L^2\) classes at every time in the horizon. The Bochner time integral therefore lives in the same Hilbert space; no pathwise differentiability of the random noise is asserted.

\section{\texorpdfstring{Iteration of the scalar error estimate}{Iteration of the scalar error estimate}}
\label{proof-module-UnadjustedLangevin-Recurrence-lean}
{\small\color{gray}Source module: \texttt{Unadjusted\allowbreak{}Langevin/\allowbreak{}Recurrence.\allowbreak{}lean}.\par}


Once coupling gives \(e_{k+1}\le qe_k+b\), the remaining calculation is deterministic. We retain the exact geometric sum before obtaining the simpler bias-floor bound. The upper envelope approaches \(b/(1-q)\); this does not assert that the actual error converges to that value.



The scalar sequence need not have an algorithmic interpretation for these lemmas. The proofs separate the exact geometric sum from its uniform bias-floor envelope and explicitly require the contraction factor in \([0,1)\). Later the coupling construction proves that the actual ULA error satisfies this recurrence.

\subsection{\texorpdfstring{Iteration gives a geometric sum}{Iteration gives a geometric sum}}


If \(q\ge0\) and \(e_{k+1}\le qe_k+b\) for every \(k\), then

\[e_n\le q^ne_0+b\sum_{j=0}^{n-1}q^j.\]

Induct on \(n\). At zero the sum is empty and both sides equal \(e_0\). Assume the bound at \(n\).



Multiply the induction inequality by \(q\ge0\), preserving its direction.



Write \(S_n=\sum_{j=0}^{n-1}q^j\). The geometric identity gives \(qS_n+1=S_n+q^n\); multiply this equality by \(b\).



Combine the one-step bound with the multiplied induction bound. Since \(S_{n+1}=S_n+q^n\), obtain \(e_{n+1}\le q^{n+1}e_0+bS_{n+1}\).

\subsection{\texorpdfstring{The exact finite-\allowbreak{}horizon envelope}{The exact finite-horizon envelope}}


If \(0\le q<1\) and \(e_{k+1}\le qe_k+b\), then

\[e_n\le q^ne_0+\frac{b(1-q^n)}{1-q}.\]

Apply the geometric-sum bound and evaluate the sum as \((q^n-1)/(q-1)\). Its denominator is nonzero because \(q<1\).



Change the signs of numerator and denominator simultaneously and put \(b\) in the numerator to obtain \(b(1-q^n)/(1-q)\).

\subsection{\texorpdfstring{A uniform bias-\allowbreak{}floor bound}{A uniform bias-floor bound}}


If \(0\le q<1\), \(b\ge0\), and \(e_{k+1}\le qe_k+b\), then

\[e_n\le q^ne_0+\frac b{1-q}.\]

Start from the exact finite-horizon bound. Its denominator \(1-q\) is positive.



Since \(q^n\ge0\) and \(b\ge0\), the numerator satisfies \(b(1-q^n)\le b\).



Divide by the positive denominator, add \(q^ne_0\), and compose with the exact bound.

\subsection{\texorpdfstring{Limit of the upper envelope}{Limit of the upper envelope}}


For \(0\le q<1\) and arbitrary real \(b,e_0\),

\[q^ne_0+\frac b{1-q}\longrightarrow\frac b{1-q}\quad(n\to\infty).\]

The geometric powers satisfy \(q^n\to0\).



Multiply this limit by \(e_0\) and add the constant \(b/(1-q)\).

\subsection{\texorpdfstring{Every accuracy above the floor is eventually attained}{Every accuracy above the floor is eventually attained}}


Assume \(0\le q<1\), \(b\ge0\), and \(e_{k+1}\le qe_k+b\). For every \(\varepsilon>b/(1-q)\), the errors satisfy \(e_n<\varepsilon\) for every sufficiently large \(n\).



The upper envelope tends to a value strictly below \(\varepsilon\), so it is eventually below \(\varepsilon\).



At each such index, place \(e_n\) below the envelope by the recurrence bound and then below \(\varepsilon\).

\section{\texorpdfstring{Euler approximations converge to the integral solution}{Euler approximations converge to the integral solution}}
\label{proof-module-UnadjustedLangevin-EulerLimit-lean}
{\small\color{gray}Source module: \texttt{Unadjusted\allowbreak{}Langevin/\allowbreak{}Euler\allowbreak{}Limit.\allowbreak{}lean}.\par}


This module proves a deterministic convergence theorem for additive Euler discretization. With mesh \(h=T/n\), its terminal value converges to the continuous integral solution, including when the state space itself is \(L^2(P)\).

\subsection{\texorpdfstring{additive\allowbreak{}Euler\allowbreak{}Chain}{additiveEulerChain}}


Define \(X_0=x\) and \(X_{k+1}=X_k-hg(X_k)+B((k+1)h)-B(kh)\). The additive input is one fixed function evaluated at mesh endpoints; consequently the numerical chain and the integral solution can be compared with exactly matching noise increments.

\subsection{\texorpdfstring{additive\allowbreak{}Euler\allowbreak{}Chain zero}{additiveEulerChain zero}}


Evaluating the recursive definition at index zero gives \(X_0=x\). Lean proves this by definitional equality; no continuity of the input or positivity of the mesh is required for this initialization identity.

\subsection{\texorpdfstring{additive\allowbreak{}Euler\allowbreak{}Chain succ}{additiveEulerChain succ}}


Unfold the successor clause once to expose \(X_{k+1}=X_k-hg(X_k)+\Delta B_k\). This definitional identity is the equation used when subtracting the exact solution's increment and deriving the error recurrence.

\subsection{\texorpdfstring{langevin integral increment}{langevin integral increment}}


Subtract the integral equations at \(b\) and \(a\). Continuity of \(g\circ Y\) makes it interval integrable, so interval-integral additivity gives \(\int_0^b-\int_0^a=\int_a^b\). Thus \(Y(b)-Y(a)=-\int_a^b g(Y(s))ds+B(b)-B(a)\), including oriented intervals when \(b<a\).

\subsection{\texorpdfstring{finite euler recurrence bound}{finite euler recurrence bound}}


From \(e_{k+1}\le qe_k+b\), \(e_0=0\), \(q\ge1\), and \(b\ge0\), induction proves \(e_n\le nbq^n\). In the successor step bound \(q\cdot kbq^k+b\) by \((k+1)bq^{k+1}\), using \(q^{k+1}\ge1\). This coarse bound avoids dividing by \(q-1\), so it also covers \(q=1\).

\subsection{\texorpdfstring{explicit\allowbreak{}Euler norm sub le}{explicitEuler norm sub le}}


Write the difference as \((x-y)-h(g(x)-g(y))\). The triangle inequality and Lipschitzness give \(\|\cdot\|\le\|x-y\|+hL\|x-y\|=(1+Lh)\|x-y\|\). The hypothesis \(h\ge0\) removes the absolute value in the scalar norm.

\subsection{\texorpdfstring{additive\allowbreak{}Euler\allowbreak{}Chain error le of drift modulus}{additiveEulerChain error le of drift modulus}}


At one step the common noise cancels. The Euler map contributes \((1+Lh)e_k\), and the exact integral remainder contributes at most \(h\varepsilon\), because the drift changes by at most \(\varepsilon\) within one mesh interval. Apply the finite recurrence bound and \((1+Lh)^n\le e^{Lnh}\), then use \(nh=T\) to obtain \(e_n\le T\varepsilon e^{LT}\).

\subsection{\texorpdfstring{additive\allowbreak{}Euler\allowbreak{}Chain tendsto}{additiveEulerChain tendsto}}


The continuous drift path \(g\circ Y\) is uniformly continuous on \([0,T]\). Given any positive error tolerance, choose a drift modulus \(\varepsilon\) making \(T\varepsilon e^{LT}\) small, then take the mesh \(T/(n+1)\) below its continuity radius. The previous error bound gives endpoint convergence. The case \(T=0\) is separated so no division by the horizon is needed.

\subsection{\texorpdfstring{lp langevin integral euler tendsto}{lp langevin integral euler tendsto}}


Apply the deterministic matching-input Euler convergence theorem in \(L^2(P;H)\), using the Lipschitz lifted gradient. The numerical endpoint converges in the \(L^2\) norm to the constructed solution at \(T\). This is the topology later used to pass smooth bounded tests and identify laws.

\section{\texorpdfstring{Couplings and quadratic transport}{Couplings and quadratic transport}}
\label{proof-module-UnadjustedLangevin-Coupling-lean}
{\small\color{gray}Source module: \texttt{Unadjusted\allowbreak{}Langevin/\allowbreak{}Coupling.\allowbreak{}lean}.\par}


Let \(H\) be a normed additive group with its Borel measurable structure and a second countable topology. Let \(P\) be a measure on a measurable space \(\Omega\), and write \(\mathcal L_P(X)=X_\#P\). The symbol \(\|Z\|_{2,P}\) denotes the extended \(L^2\) seminorm, which may be infinite. In the probability setting the following definitions give the usual quadratic transport cost.



\texttt{Coupling} stores a joint measure with exact first and second pushforward marginals. \texttt{coupling\allowbreak{}Cost} is its extended \(L^2\) displacement norm, and \texttt{wasserstein2} is the infimum over all such couplings. Finiteness guards precede \texttt{to\allowbreak{}Real} conversions; no existence of an optimal coupling is assumed.

\subsection{\texorpdfstring{Definition of a coupling}{Definition of a coupling}}


A coupling of measures \(\mu,\nu\) on \(H\) is a joint measure \(\gamma\) on \(H\times H\) with those two marginals. Denote the family by \(\Gamma(\mu,\nu)\).



Choose the joint measure \(\gamma\) on pairs \((x,y)\).



Require \((\operatorname{pr}_1)_\#\gamma=\mu\) and \((\operatorname{pr}_2)_\#\gamma=\nu\).

\subsection{\texorpdfstring{Definition of quadratic coupling cost}{Definition of quadratic coupling cost}}


For a joint measure \(\gamma\), define

\[c(\gamma)=\left(\int\|x-y\|^2\,d\gamma(x,y)\right)^{1/2}\in[0,+\infty].\]

Take the extended \(L^2(\gamma)\) norm of the displacement map \((x,y)\mapsto x-y\).

\subsection{\texorpdfstring{Definition of the transport infimum}{Definition of the transport infimum}}


Define the extended quadratic transport distance by

\[W_2(\mu,\nu)=\inf_{\gamma\in\Gamma(\mu,\nu)}c(\gamma).\]

Take the infimum over all couplings without selecting a minimizer.

\subsection{\texorpdfstring{Real conversion of transport distance}{Real conversion of transport distance}}


Define \(W_2^{\mathbb R}(\mu,\nu)=\operatorname{toReal}(W_2(\mu,\nu))\). The comparison lemmas below establish finite upper bounds before using this conversion.



Apply real conversion to the extended distance; for a finite distance this returns its ordinary real value.

\subsection{\texorpdfstring{Every coupling supplies an upper bound}{Every coupling supplies an upper bound}}


For every \(\gamma\in\Gamma(\mu,\nu)\), \(W_2(\mu,\nu)\le c(\gamma)\).



An infimum is at most each member of the family. Apply this to the chosen coupling.

\subsection{\texorpdfstring{The joint law of two random vectors}{The joint law of two random vectors}}


If \(X,Y:\Omega\to H\) are measurable, then \((X,Y)_\#P\) is a coupling of \(X_\#P\) and \(Y_\#P\).



Set \(\gamma=(\omega\mapsto(X(\omega),Y(\omega)))_\#P\).



Compose the pushforward with the first projection. Since \(\operatorname{pr}_1(X(\omega),Y(\omega))=X(\omega)\), the first marginal is \(X_\#P\).



The same calculation with the second projection gives the second marginal \(Y_\#P\).

\subsection{\texorpdfstring{Joint laws from almost everywhere measurable vectors}{Joint laws from almost everywhere measurable vectors}}


If \(X,Y\) are measurable almost everywhere with respect to \(P\), their joint law is still a coupling of their two laws.



Again define \(\gamma=(X,Y)_\#P\).



Apply pushforward composition for almost everywhere measurable maps to identify the first marginal with \(X_\#P\).



Apply its second-projection version to identify the second marginal with \(Y_\#P\).

\subsection{\texorpdfstring{Cost of the joint-\allowbreak{}law coupling}{Cost of the joint-law coupling}}


For measurable \(X,Y:\Omega\to H\),

\[c((X,Y)_\#P)=\|X-Y\|_{2,P}.\]

Expand the two definitions. The cost is the \(L^2\) norm of \((x,y)\mapsto x-y\) under the joint law.



Change variables through the pushforward. The displacement composed with the paired random vector is \(\omega\mapsto X(\omega)-Y(\omega)\), giving the identity.

\subsection{\texorpdfstring{A coupled pair bounds transport}{A coupled pair bounds transport}}


For measurable \(X,Y\),

\[W_2(X_\#P,Y_\#P)\le\|X-Y\|_{2,P}.\]

Use the joint law as an admissible coupling, then replace its cost by the \(L^2\) distance just computed.

\subsection{\texorpdfstring{The almost everywhere measurable transport bound}{The almost everywhere measurable transport bound}}


The same transport comparison holds for almost everywhere measurable \(X,Y\):

\[W_2(X_\#P,Y_\#P)\le\|X-Y\|_{2,P}.\]

The almost everywhere measurable joint-law construction supplies an admissible coupling, whose cost bounds the infimum.



The pushforward change of variables identifies that cost with \(\|X-Y\|_{2,P}\).

\subsection{\texorpdfstring{Finite coupled cost gives a real bound}{Finite coupled cost gives a real bound}}


If \(X,Y\) are measurable and \(X-Y\in L^2(P;H)\), then

\[W_2^{\mathbb R}(X_\#P,Y_\#P)\le\operatorname{toReal}(\|X-Y\|_{2,P}).\]

The extended comparison has a finite upper bound because \(X-Y\in L^2\). Apply monotonicity of real conversion under a finite upper bound.

\subsection{\texorpdfstring{The transport bound for L\ensuremath{{}^{2}} classes}{The transport bound for L² classes}}


For \(X,Y\in L^2(P;H)\), their representatives define laws satisfying

\[W_2(X_\#P,Y_\#P)\le\operatorname{ofReal}(\|X-Y\|_{L^2}).\]

The representative of \(X-Y\) agrees almost everywhere with the difference of the representatives. Their extended \(L^2\) norms agree.



The L\ensuremath{{}^{2}} norm is the real value of the finite extended seminorm. Embedding it back into extended reals recovers that seminorm.



Apply the almost everywhere measurable transport bound to the two representatives and use the equality of displacement seminorms.

\subsection{\texorpdfstring{The real L\ensuremath{{}^{2}}-\allowbreak{}to-\allowbreak{}transport bridge}{The real L²-to-transport bridge}}


For \(X,Y\in L^2(P;H)\),

\[W_2^{\mathbb R}(X_\#P,Y_\#P)\le\|X-Y\|_{L^2}.\]

The preceding extended bound is dominated by the finite number \(\operatorname{ofReal}(\|X-Y\|_{L^2})\), so convert monotonically to real numbers.



Real conversion of the embedding of a nonnegative norm returns the norm itself.

\subsection{\texorpdfstring{Pointwise domination passes to L\ensuremath{{}^{2}}}{Pointwise domination passes to L²}}


Let \(q\ge0\), let \(A\) be almost everywhere strongly measurable, and suppose \(\|A(\omega)\|\le q\|D(\omega)\|\) almost everywhere. Then

\[\|A\|_{2,P}\le q\|D\|_{2,P}.\]

Apply monotonicity and positive homogeneity of the extended \(L^2\) seminorm to the almost everywhere domination. No independence is required.

\subsection{\texorpdfstring{One-\allowbreak{}step error from contraction and a defect}{One-step error from contraction and a defect}}


Suppose \(q\ge0\), \(A\) is almost everywhere strongly measurable, and almost everywhere

\[X'-Y'=A-V,\qquad\|A\|\le q\|X-Y\|.\]

Then

\[\|X'-Y'\|_{2,P}\le q\|X-Y\|_{2,P}+\|V\|_{2,P}.\]

Replace \(X'-Y'\) by \(A-V\) in the seminorm using the almost everywhere error identity.



Minkowski gives \(\|A-V\|_2\le\|A\|_2+\|V\|_2\). Pointwise contraction bounds the first term by \(q\|X-Y\|_2\).

\subsection{\texorpdfstring{One-\allowbreak{}step transport estimate}{One-step transport estimate}}


Under the preceding decomposition and contraction hypotheses, if \(X',Y'\) are measurable, then

\[W_2(X'_\#P,Y'_\#P)\le q\|X-Y\|_{2,P}+\|V\|_{2,P}.\]

The joint-law coupling bounds transport by \(\|X'-Y'\|_2\). Compose this with the synchronous \(L^2\) estimate.

\subsection{\texorpdfstring{The finite-\allowbreak{}moment one-\allowbreak{}step inequality}{The finite-moment one-step inequality}}


Assume \(q\ge0\), \(A\) is almost everywhere strongly measurable, \(X'-Y'=A-V\) and \(\|A\|\le q\|X-Y\|\) almost everywhere. If \(X-Y,V\in L^2(P;H)\), then

\[\operatorname{toReal}(\|X'-Y'\|_{2,P})\le q\operatorname{toReal}(\|X-Y\|_{2,P})+\operatorname{toReal}(\|V\|_{2,P}).\]

The product \(q\|X-Y\|_2\) is finite because both factors are finite.



Its sum with the finite defect norm is finite. Use this finite upper bound to convert the synchronous extended inequality monotonically to real numbers.



Distribute real conversion through the finite sum and the product, giving the stated real inequality.

\section{\texorpdfstring{Passing from a chosen initial coupling to an infimum}{Passing from a chosen initial coupling to an infimum}}
\label{proof-module-UnadjustedLangevin-InitialCoupling-lean}
{\small\color{gray}Source module: \texttt{Unadjusted\allowbreak{}Langevin/\allowbreak{}Initial\allowbreak{}Coupling.\allowbreak{}lean}.\par}


The analytic estimate first uses the cost of one initial coupled pair. The infimum calculation here replaces that cost by \(W_2(\mu,\nu)\) when all finite-cost couplings give the same output laws. \texttt{Product\allowbreak{}Reference.\allowbreak{}lean} proves this common-law property on concrete product extensions; \texttt{Canonical\allowbreak{}ULA.\allowbreak{}lean} supplies their stationary references.



The mathematical infimum here is the infimum in \texttt{wasserstein2}. The proof works with arbitrary finite-cost couplings and approximates that infimum when necessary. This gives the true initial transport term without postulating an optimal transport plan.

\subsection{\texorpdfstring{A uniform affine estimate passes to the infimum}{A uniform affine estimate passes to the infimum}}


Let \(a\ge0\) be finite, \(b\in[0,+\infty]\), and \(\Gamma(\mu,\nu)\ne\varnothing\). If the fixed output measures satisfy

\[W_2(\mu',\nu')\le ac(\gamma)+b\quad\text{for every }\gamma\in\Gamma(\mu,\nu),\]

then \(W_2(\mu',\nu')\le aW_2(\mu,\nu)+b\).



Use the nonempty coupling family and expand transport distances as infima.



Multiplication by finite \(a\ge0\) and addition of \(b\) commute with the nonempty extended-real infimum, including \(a=0\). Thus \(aW_2(\mu,\nu)+b=\inf_\gamma(ac(\gamma)+b)\).



The output distance lies below every member of this last family, hence below its infimum.

\subsection{\texorpdfstring{A finite infimum supplies a finite-\allowbreak{}cost coupling}{A finite infimum supplies a finite-cost coupling}}


If \(W_2(\mu,\nu)<+\infty\), then some \(\gamma\in\Gamma(\mu,\nu)\) has \(c(\gamma)<+\infty\). Attainment of the infimum is not asserted.



The coupling-cost infimum is strictly below \(+\infty\) by the finiteness hypothesis.



An infimum strictly below a threshold has a member strictly below that threshold. Apply this with threshold \(+\infty\).

\subsection{\texorpdfstring{The real uniform-\allowbreak{}coupling estimate}{The real uniform-coupling estimate}}


Let \(a,b\ge0\) be finite and assume \(W_2(\mu,\nu)<+\infty\). If \(W_2(\mu',\nu')\le ac(\gamma)+b\) for every initial coupling, then

\[W_2^{\mathbb R}(\mu',\nu')\le aW_2^{\mathbb R}(\mu,\nu)+b.\]

A finite initial distance supplies a coupling. Apply the extended uniform-coupling lemma to obtain the desired inequality in extended reals.



The product \(aW_2(\mu,\nu)\) is finite.



Its sum with finite \(b\) is finite, so the extended output inequality permits monotone real conversion.



Distribute real conversion through the finite sum and product to obtain the real inequality.

\subsection{\texorpdfstring{Finite-\allowbreak{}cost initial couplings suffice}{Finite-cost initial couplings suffice}}


Let \(a,b\ge0\) be finite. Assume both the initial and output transport distances are finite and that every finite-cost \(\gamma\in\Gamma(\mu,\nu)\) satisfies

\[W_2^{\mathbb R}(\mu',\nu')\le a\operatorname{toReal}(c(\gamma))+b.\]

Then \(W_2^{\mathbb R}(\mu',\nu')\le aW_2^{\mathbb R}(\mu,\nu)+b\).



For each finite-cost coupling, embed the real estimate in extended nonnegative reals. Finiteness of the output distance and coupling cost recovers their original extended values, giving \(W_2(\mu',\nu')\le ac(\gamma)+b\).



Apply the all-couplings theorem. For a finite-cost coupling, the needed extended estimate has just been proved.



For an infinite-cost coupling, split into the cases \(a=0\) and \(a\ne0\).



If \(a=0\), choose one finite-cost coupling supplied by the finite initial distance. Its estimate gives \(W_2(\mu',\nu')\le b\), which also supplies the required zero-coefficient estimate for the infinite-cost coupling.



If \(a\ne0\), then \(a>0\), so \(a\cdot(+\infty)+b=+\infty\) and the required inequality is automatic.

\section{\texorpdfstring{The shared Hilbert-\allowbreak{}space proof and both step regimes}{The shared Hilbert-space proof and both step regimes}}
\label{proof-module-UnadjustedLangevin-HilbertDalalyan-lean}
{\small\color{gray}Source module: \texttt{Unadjusted\allowbreak{}Langevin/\allowbreak{}Hilbert\allowbreak{}Dalalyan.\allowbreak{}lean}.\par}


Work on a probability space with \(\mathcal H=L^2(P;H)\), and write \(g=\widehat{\nabla U}\). This is the single analytic proof used by both Euclidean ULA and the lifted IPLA system. The reference record exposes the integral equation, stationary laws, and moment bounds used at this stage. \texttt{Canonical\allowbreak{}ULA.\allowbreak{}lean} constructs these inputs from strong convexity, a Lipschitz gradient, and independent Gaussian innovations.



\texttt{L2\allowbreak{}Reference} stores a stationary reference path in the same \(L^2\) space and with the same innovations as the numerical chain. The analytic theorem derives its recurrence and final transport estimate. The reference is a supplied argument at this layer; the canonical Brownian construction later proves its fields from the potential assumptions.

\subsection{\texorpdfstring{The common L\ensuremath{{}^{2}} algorithm}{The common L² algorithm}}


For an initial \(X_0\in L^2(P;H)\) and innovations \(\xi_k\in L^2(P;H)\), set \(X_{k+1}=X_k-h\widehat{\nabla U}(X_k)+\sqrt{2h}\,\xi_k\).



The initial iterate is the supplied random vector.



Each successor uses the actual-gradient Euler update.

\subsection{\texorpdfstring{Unfolding one step}{Unfolding one step}}


The next iterate is exactly one L\ensuremath{{}^{2}} ULA update of the current iterate.



Evaluate the successor clause of the definition.

\subsection{\texorpdfstring{The continuous reference assumptions}{The continuous reference assumptions}}


For an \(L\)-Lipschitz gradient, step \(h\), noise vectors \(\xi_k\in\mathcal H\), and a measure \(\pi\), a reference consists of segments \(Y_{k,t}\in\mathcal H\) and increments \(B_{k,t}\in\mathcal H\) satisfying the following conditions.



Supply the two families \((Y_{k,t})\) and \((B_{k,t})\) in the Hilbert space \(\mathcal H\).



Require continuity of \(t\mapsto Y_{k,t}\) in the \(L^2\) norm for every \(k\).



Require compatible endpoints, common endpoint noise, and stationary grid laws:

\[Y_{k,h}=Y_{k+1,0},\qquad B_{k,h}=\sqrt{2h}\,\xi_k,\qquad(Y_{k,0})_\#P=\pi.\]

For \(0\le t\le h\), require

\[Y_{k,t}-Y_{k,0}=-\int_0^t g(Y_{k,s})\,ds+B_{k,t}.\]

Require the score moment bound \(\|g(Y_{k,t})\|_2\le\sqrt{Ld}\) throughout each closed step interval.



Require the increment moment bound \(\|B_{k,t}\|_2\le\sqrt{2d}\sqrt t\) on the same interval.

\subsection{\texorpdfstring{Full convergence calculation from the integral reference}{Full convergence calculation from the integral reference}}


Suppose a reference with the preceding properties is supplied, \(h\ge0\), and \(Lh\le2\). Let \(0\le q<1\), and assume the lifted Euler map contracts:

\[\|(X-hg(X))-(Y-hg(Y))\|_2\le q\|X-Y\|_2.\]

For any initial vector \(X_0\in\mathcal H\), the \(L^2\) chain driven by the reference endpoint noise satisfies

\[W_2^{\mathbb R}((X_n)_\#P,\pi)\le q^n\|X_0-Y_{0,0}\|_2+\frac{(33/20)Lh\sqrt{hd}}{1-q}.\]

The statement is conditional on the supplied reference, and the initial term is the cost of the chosen coupling.



Set \(g=\widehat{\nabla U}\), name the actual iterates \(X_k\), and define

\[b=\frac{33}{20}Lh\sqrt{hd},\qquad D_k=\int_0^h\bigl(g(Y_{k,t})-g(Y_{k,0})\bigr)\,dt.\]

The bound \(b\) is nonnegative. Apply the local-defect theorem to the actual lifted gradient and the supplied equation and moment bounds, deriving \(\|D_k\|_2\le b\).



We claim the reference endpoints obey

\[Y_{k+1,0}=Y_{k,0}-hg(Y_{k,0})+\sqrt{2h}\,\xi_k-D_k.\]

The lifted gradient is Lipschitz and the reference segment is continuous, so their composition is continuous and Bochner integrable on \([0,h]\).



Apply the Euler-defect identity at \(t=h\), then substitute the segment join and endpoint noise identities. This proves the claimed reference update. Its defect has a minus sign.



Fix \(k\) and compare the actual update with the reference endpoint equation, both using the same \(\xi_k\).



Cancel their identical noise terms. Because the reference subtracts \(D_k\),

\[X_{k+1}-Y_{k+1,0}=\bigl(X_k-hg(X_k)\bigr)-\bigl(Y_{k,0}-hg(Y_{k,0})\bigr)+D_k.\]

Apply the triangle inequality, Euler contraction, and the defect bound. Setting \(e_k=\|X_k-Y_{k,0}\|_2\) gives \(e_{k+1}\le qe_k+b\).



The actual pair supplies a finite-cost coupling, so its transport distance is at most \(e_n\). The second marginal is \(\pi\) by the stationary grid-law assumption.



Apply the scalar recurrence bound to the derived inequality and compose it with the transport comparison.

\subsection{\texorpdfstring{The small-\allowbreak{}step bound}{The small-step bound}}


Suppose \(U\) is differentiable and \(m\)-strongly convex on \(H\), \(\nabla U\) is \(L\)-Lipschitz, and \(0<m\le L\). Let \(0<h\le2/(m+L)\), and supply a reference with the preceding properties. For any initial vector \(X_0\in\mathcal H\),

\[W_2^{\mathbb R}((X_n)_\#P,\pi)\le(1-mh)^n\|X_0-Y_{0,0}\|_2+\frac{33}{20}\frac Lm\sqrt{hd}.\]

This uses a chosen initial coupling and a supplied stationary reference. Replacing the initial cost by \(W_2(\mathcal L(X_0),\pi)\) needs the additional uniform family of reference couplings explained in the final section.



Since \(m+L>0\), clear the positive denominator in the step hypothesis to obtain \(h(m+L)\le2\).



Use \(m\le L\) to get \(2mh\le h(m+L)\le2\), hence \(mh\le1\).



The same threshold gives \(Lh\le2\). Thus \(q=1-mh\ge0\); since \(mh>0\), also \(q<1\).



The sharp small-step contraction of the lifted Euler map gives the contraction hypothesis with \(q=1-mh\). Apply the integral-reference theorem.



Because \(m,h>0\), simplify the bias term:

\[\frac{(33/20)Lh\sqrt{hd}}{1-(1-mh)}=\frac{(33/20)Lh\sqrt{hd}}{mh}=\frac{33}{20}\frac Lm\sqrt{hd}.\]

Substitute the bias identity into the integral-reference bound.

\subsection{\texorpdfstring{The large-\allowbreak{}step bound}{The large-step bound}}


Suppose \(U\) is differentiable and \(m\)-strongly convex on \(H\), \(\nabla U\) is \(L\)-Lipschitz, and \(0<m\le L\). Let \(h>0\), \(h\ge2/(m+L)\), and \(Lh<2\), and supply a reference with the stated properties. For any initial vector \(X_0\in\mathcal H\),

\[W_2^{\mathbb R}((X_n)_\#P,\pi)\le(Lh-1)^n\|X_0-Y_{0,0}\|_2+\frac{33}{20}\frac{Lh}{2-Lh}\sqrt{hd}.\]

The denominator is positive by strict stability. The reference remains an input, and the initial term remains the chosen coupling cost.



Since \(m+L>0\), clear the denominator in the lower step threshold to obtain \(2\le h(m+L)\).



Use \(m\le L\) and \(h\ge0\): \(2\le h(m+L)\le2Lh\). Thus \(Lh\ge1\) and \(q=Lh-1\ge0\).



Strict stability \(Lh<2\) gives \(q=Lh-1<1\).



Apply the sharp large-step Euler contraction with \(q=Lh-1\), then invoke the integral-reference theorem.



Use \(1-(Lh-1)=2-Lh\) and factor out \(\sqrt{hd}\):

\[\frac{(33/20)Lh\sqrt{hd}}{1-(Lh-1)}=\left(\frac{(33/20)Lh}{2-Lh}\right)\sqrt{hd}.\]

Substitute this identity into the integral-reference bound to obtain the large-step formula.

\section{\texorpdfstring{A noise-\allowbreak{}driven recursion has a law independent of its realization}{A noise-driven recursion has a law independent of its realization}}
\label{proof-module-UnadjustedLangevin-ChainLaw-lean}
{\small\color{gray}Source module: \texttt{Unadjusted\allowbreak{}Langevin/\allowbreak{}Chain\allowbreak{}Law.\allowbreak{}lean}.\par}


For measurable \(F(x,z)\), the recursion \(X_{n+1}=F(X_n,Z_n)\) has a law determined by the initial law and the full noise law when the initialization and noise are independent.

\subsection{\texorpdfstring{noise\allowbreak{}Driven\allowbreak{}Chain}{noiseDrivenChain}}


Define \(X_0=x\), \(X_{n+1}=F(X_n,w_n)\) for any update and noise path. No distribution assumptions are part of the recursion; they enter only when a law is assigned to its inputs.

\subsection{\texorpdfstring{noise\allowbreak{}Driven\allowbreak{}Chain zero}{noiseDrivenChain zero}}


The zero clause of the recursion evaluates to the supplied initial state. This definitional identity is the base case for all law and measurability inductions.

\subsection{\texorpdfstring{noise\allowbreak{}Driven\allowbreak{}Chain succ}{noiseDrivenChain succ}}


The successor clause unfolds to \(X_{n+1}=F(X_n,w_n)\). This exact recursion lets the generic construction later be identified with ULA by induction.

\subsection{\texorpdfstring{noise\allowbreak{}Driven\allowbreak{}Chain measurable}{noiseDrivenChain measurable}}


Induct on \(n\). The base map is the initial-state projection. At the successor, pair the measurable current state with the measurable coordinate evaluation \(w\mapsto w_n\), then compose with the measurable uncurried update. Thus each iterate is a measurable function of initialization and the full noise path.

\subsection{\texorpdfstring{chain\allowbreak{}Law}{chainLaw}}


Push \(\mu\otimes Q\) through the measurable recursion output at index \(n\). The product gives independent initialization and noise path. This canonical law depends only on the initial law, noise-path law, and update, not on a particular sample-space implementation.

\subsection{\texorpdfstring{chain\allowbreak{}Law zero}{chainLaw zero}}


At index zero the output map is first-coordinate projection. Since \(Q\) has mass one, the first marginal of \(\mu\otimes Q\) is exactly \(\mu\). This proves that the canonical law starts from the supplied initial law.

\subsection{\texorpdfstring{noise\allowbreak{}Driven\allowbreak{}Chain law of indep}{noiseDrivenChain law of indep}}


Independence gives the joint input law \(P_{(X,W)}=P_X\otimes P_W\). Push that identity through the measurable recursion output and use pushforward composition. The resulting output law is exactly \texttt{chain\allowbreak{}Law F (P.\allowbreak{}map X) (P.\allowbreak{}map W) n}.

\subsection{\texorpdfstring{noise\allowbreak{}Driven\allowbreak{}Chain law of iid}{noiseDrivenChain law of iid}}


Joint independence of the innovations and their common law \(\eta\) identify the entire noise-path law with \(\bigotimes_n\eta\). Substitute this identity into the independent-initialization law theorem. Countability of the time index supplies path-valued measurability.

\subsection{\texorpdfstring{noise\allowbreak{}Driven\allowbreak{}Chain law unique}{noiseDrivenChain law unique}}


Apply the canonical iid-law formula on both probability spaces. Equality of their initial laws and the shared innovation law makes the right sides identical, so their iterate laws agree. The theorem compares distributions without coupling the two implementations.

\subsection{\texorpdfstring{noise\allowbreak{}Driven\allowbreak{}Chain eq ula\allowbreak{}Chain}{noiseDrivenChain eq ulaChain}}


Induct on the index with \(F=\operatorname{ulaUpdate}\). The zero values agree, and the successor equations apply the same update to equal previous states and the same noise vector. This proves pointwise equality of the generic recursion and the repository's ULA chain.

\subsection{\texorpdfstring{ula\allowbreak{}Update measurable}{ulaUpdate measurable}}


Lipschitzness makes \(\nabla U\) continuous. The joint ULA update is obtained from coordinate projections and that gradient by scalar multiplication and addition, hence is measurable. This discharges the generic recursion's update-measurability premise.

\subsection{\texorpdfstring{ula\allowbreak{}Chain law of iid}{ulaChain law of iid}}


Use measurability of \texttt{ula\allowbreak{}Update} in the generic iid law theorem, then rewrite the generic recursion using \texttt{noise\allowbreak{}Driven\allowbreak{}Chain\_\allowbreak{}eq\_\allowbreak{}ula\allowbreak{}Chain}. The result identifies the actual ULA implementation with its canonical product-noise law.

\subsection{\texorpdfstring{lp\allowbreak{}Ula\allowbreak{}Chain law of iid}{lpUlaChain law of iid}}


The \(L^2\) chain representative agrees almost everywhere with the pointwise ULA recursion of the input representatives. Apply the pointwise law theorem and transfer pushforwards across that almost-everywhere equality. Quotient representative choices therefore do not change the algorithm law.

\section{\texorpdfstring{Finite quadratic transport transfers second moments}{Finite quadratic transport transfers second moments}}
\label{proof-module-UnadjustedLangevin-TransportMoments-lean}
{\small\color{gray}Source module: \texttt{Unadjusted\allowbreak{}Langevin/\allowbreak{}Transport\allowbreak{}Moments.\allowbreak{}lean}.\par}


Finite \(W_2\) distance to a finite-second-moment law supplies the other second moment automatically. This removes a redundant integrability premise when passing an initial law to the L\ensuremath{{}^{2}} formulation.

\subsection{\texorpdfstring{mem\allowbreak{}Lp id of wasserstein2 ne top}{memLp id of wasserstein2 ne top}}


Finite Wasserstein distance supplies a finite-cost coupling. On that coupling space the second coordinate is in \(L^2\) by the reference moment, and the coordinate difference is in \(L^2\) by finite coupling cost. Their sum is the first coordinate, so its law \(\mu\) has a finite second moment. No optimal coupling is selected.

\subsection{\texorpdfstring{wasserstein2 ne top of mem\allowbreak{}Lp}{wasserstein2 ne top of memLp}}


Use the explicit independent coupling \(\mu\otimes\pi\). Each coordinate is in \(L^2\), hence their difference is in \(L^2\) and has finite cost. Wasserstein distance is at most this coupling cost, proving finiteness.

\subsection{\texorpdfstring{mem\allowbreak{}Lp of law of wasserstein2 ne top}{memLp of law of wasserstein2 ne top}}


The finite-transport theorem first proves \(\mathrm{id}\in L^2(\mu)\). Replace \(\mu\) by the supplied pushforward law of \(X\), then use the change-of-law \(L^2\) criterion to pull this moment back to \(X\) on its original sample space.

\subsection{\texorpdfstring{mem\allowbreak{}Lp of wasserstein2 map ne top}{memLp of wasserstein2 map ne top}}


Specialize the preceding theorem to \(\mu=P\circ X^{-1}\), whose law identity is reflexive. Finite transport of the actual pushforward to a finite-moment reference therefore gives \(X\in L^2(P)\).

\section{\texorpdfstring{Realize every initial coupling independently of the noise}{Realize every initial coupling independently of the noise}}
\label{proof-module-UnadjustedLangevin-InitialExtension-lean}
{\small\color{gray}Source module: \texttt{Unadjusted\allowbreak{}Langevin/\allowbreak{}Initial\allowbreak{}Extension.\allowbreak{}lean}.\par}


Given any coupling \(\gamma\) of \(\mu\) and \(\pi\), extend it to \(\gamma\otimes Q\). This preserves its quadratic cost and supplies independent noise without requiring an optimal coupling.

\subsection{\texorpdfstring{Coupling.\allowbreak{}is\allowbreak{}Probability\allowbreak{}Measure}{Coupling.isProbabilityMeasure}}


The coupling's first marginal is the initial probability. Pushforward preserves total mass, so \(c(H\times H)=\mu(H)=1\). This constructs normalization of the joint measure using a marginal identity already stored in \texttt{Coupling}.

\subsection{\texorpdfstring{initial\allowbreak{}Extension}{initialExtension}}


Extend an initial coupling \(c(dx,dy)\) by an independent noise sample \(s\sim Q\), using \(c\otimes Q\) on \((H\times H)\times S\). Both algorithms retain the prescribed initial joint law while sharing fresh noise independent of the pair.

\subsection{\texorpdfstring{initial\allowbreak{}Extension is\allowbreak{}Probability\allowbreak{}Measure}{initialExtension isProbabilityMeasure}}


The coupling is a probability by its first marginal, and \(Q\) is a probability by hypothesis. Their product therefore has total mass one. This instance supports all subsequent \(L^2\) constructions on the extension.

\subsection{\texorpdfstring{initial\allowbreak{}Extension joint}{initialExtension joint}}


Projecting \(c\otimes Q\) to the first factor integrates out the probability noise coordinate. Since \(Q(S)=1\), the resulting measure is exactly \(c\), preserving the entire initial dependence structure.

\subsection{\texorpdfstring{initial\allowbreak{}Extension initial}{initialExtension initial}}


Compose the joint projection with its first state projection. The extension's joint marginal is \(c\), whose first state marginal is \(\mu\). Pushforward composition therefore yields the specified algorithm initialization law.

\subsection{\texorpdfstring{initial\allowbreak{}Extension reference}{initialExtension reference}}


Compose projection to the initial pair with its second state projection. The coupling's second marginal is \(\pi\), so the reference initialization on the extension has exactly law \(\pi\).

\subsection{\texorpdfstring{initial\allowbreak{}Extension noise}{initialExtension noise}}


The coupling measure has mass one. Projecting its product with \(Q\) to the noise factor therefore returns \(Q\). The \texttt{SFinite} hypothesis supplies the product marginal theorem's measure-theoretic condition.

\subsection{\texorpdfstring{initial\allowbreak{}Extension independent}{initialExtension independent}}


The two coordinate projections of a product probability measure are independent. Here those coordinates are the full coupled initial pair and the fresh noise sample. Independence of either individual initial state from any noise-derived random object follows by composition.

\subsection{\texorpdfstring{initial\allowbreak{}Extension initial noise}{initialExtension initial noise}}


Map \(((x,y),s)\) to \((x,s)\). Product pushforward and the first marginal identity of \(c\) give \(\mu\otimes Q\). Thus removing the reference coordinate leaves the same algorithm input law for every initial coupling.

\subsection{\texorpdfstring{initial\allowbreak{}Extension map}{initialExtension map}}


Factor the map through \(((x,y),s)\mapsto(x,s)\), whose law is \(\mu\otimes Q\). Pushforward composition proves that every measurable output depending only on algorithm initialization and noise has the coupling-independent law \(f_*(\mu\otimes Q)\).

\subsection{\texorpdfstring{initial\allowbreak{}Extension chain\allowbreak{}Law}{initialExtension chainLaw}}


Apply \texttt{initial\allowbreak{}Extension\_\allowbreak{}map} to the measurable recursion output. Its canonical pushforward is exactly \texttt{chain\allowbreak{}Law}, so the same output law is realized on every initial-coupling extension. This equality is derived rather than made a field of the coupling.

\subsection{\texorpdfstring{initial\allowbreak{}Extension iid noise\allowbreak{}Law}{initialExtension iid noiseLaw}}


First project the extension to its product-path noise law, then evaluate coordinate \(n\). The marginal of \(\bigotimes_k\eta\) at that coordinate is \(\eta\), giving the requested noise law on the extension.

\subsection{\texorpdfstring{initial\allowbreak{}Extension iid independent}{initialExtension iid independent}}


The iid path has independent coordinate evaluations. Pull this independence back along the extension's measure-preserving noise projection. The resulting innovation sequence remains jointly independent on the coupled initialization space.

\subsection{\texorpdfstring{initial\allowbreak{}Extension initial independent}{initialExtension initial independent}}


The algorithm's initial state is a projection of the initial pair. Compose the pair/\allowbreak{}noise independence theorem with this state projection. This proves independence from the entire noise sample, not just from each individual noise vector.

\subsection{\texorpdfstring{initial\allowbreak{}Extension cost}{initialExtension cost}}


The difference \(x-y\) depends only on the initial pair. Its \(L^2\) extended norm under \(c\otimes Q\) equals the same norm under \(c\), because the noise factor has mass one. The latter is exactly \texttt{coupling\allowbreak{}Cost c.\allowbreak{}joint}.

\subsection{\texorpdfstring{initial\allowbreak{}Extension mem\allowbreak{}Lp difference}{initialExtension memLp difference}}


The state-difference function is measurable. Its extended \(L^2\) norm equals the finite initial coupling cost by the previous theorem. These two facts prove \texttt{Mem\allowbreak{}Lp} for the difference on the extension.

\subsection{\texorpdfstring{initial\allowbreak{}Extension mem\allowbreak{}Lp reference}{initialExtension memLp reference}}


The reference coordinate has law \(\pi\). Pull its supplied \(L^2(\pi)\) identity moment back through this law to get square integrability on the extension. This establishes the reference initialization without choosing a separate realization.

\subsection{\texorpdfstring{initial\allowbreak{}Extension mem\allowbreak{}Lp initial}{initialExtension memLp initial}}


Write \(x=(x-y)+y\). The difference is in \(L^2\) by finite coupling cost and the reference coordinate is in \(L^2\) by its law. Closure of \(L^2\) under addition gives square integrability of the initial algorithm state; no independent moment premise on \(\mu\) is introduced.

\subsection{\texorpdfstring{initial\allowbreak{}Extension lp\allowbreak{}Initial}{initialExtension lpInitial}}


Use \texttt{initial\allowbreak{}Extension\_\allowbreak{}mem\allowbreak{}Lp\_\allowbreak{}initial} to convert the first state coordinate into an actual \(L^2\) class. The class is represented almost everywhere by \(((x,y),s)\mapsto x\), preserving its concrete relation to the coupling coordinates.

\subsection{\texorpdfstring{initial\allowbreak{}Extension lp\allowbreak{}Reference}{initialExtension lpReference}}


Convert the square-integrable second state coordinate into an \(L^2\) class. Its representative equals \(((x,y),s)\mapsto y\) almost everywhere, with law \(\pi\), and it shares the same product probability space as the algorithm initialization.

\subsection{\texorpdfstring{initial\allowbreak{}Extension lp\allowbreak{}Initial law}{initialExtension lpInitial law}}


The representative chosen by \texttt{to\allowbreak{}Lp} agrees almost everywhere with the first state coordinate. Pushforwards are unchanged by such an equality. Combine this with \texttt{initial\allowbreak{}Extension\_\allowbreak{}initial} to obtain the initial law \(\mu\).

\subsection{\texorpdfstring{initial\allowbreak{}Extension lp\allowbreak{}Reference law}{initialExtension lpReference law}}


Replace the selected \(L^2\) representative by the second state projection almost everywhere. Its already proved marginal is \(\pi\), hence the class used as reference initialization has exactly that law.

\subsection{\texorpdfstring{initial\allowbreak{}Extension lp cost}{initialExtension lp cost}}


The representative of the class difference equals \(x-y\) almost everywhere. The norm of that \(L^2\) difference is the real value of its extended norm; \texttt{initial\allowbreak{}Extension\_\allowbreak{}cost} identifies it with \((\operatorname{couplingCost}c)_{\mathbb R}\). The finite-cost premise justifies this real conversion.

\section{\texorpdfstring{A reference for every initial coupling with one output law}{A reference for every initial coupling with one output law}}
\label{proof-module-UnadjustedLangevin-ProductReference-lean}
{\small\color{gray}Source module: \texttt{Unadjusted\allowbreak{}Langevin/\allowbreak{}Product\allowbreak{}Reference.\allowbreak{}lean}.\par}


This construction turns the analytic reference estimate into an exact initial Wasserstein term by realizing each finite initial coupling on a product with one fixed noise source.

\subsection{\texorpdfstring{pointwise\allowbreak{}Update}{pointwiseUpdate}}


Define the ordinary Euler update \(F_h(x,z)=x-h\nabla U(x)+\sqrt{2h}z\) in an arbitrary Hilbert-space presentation. It is the pointwise function whose lifted composition gives the \(L^2\) numerical chain.

\subsection{\texorpdfstring{pointwise\allowbreak{}Update measurable}{pointwiseUpdate measurable}}


The Lipschitz gradient is continuous. Compose it with the first coordinate projection and combine the remaining continuous scalar and vector operations. The uncurried map \((x,z)\mapsto F_h(x,z)\) is therefore measurable.

\subsection{\texorpdfstring{product\allowbreak{}Noise mem\allowbreak{}Lp}{productNoise memLp}}


The noise projection from \(c\otimes Q\) to \(Q\) is measure preserving because the coupling is a probability. Pull back the supplied \(L^2(Q)\) noise function along this projection. Its square integrability is unchanged.

\subsection{\texorpdfstring{product\allowbreak{}Noise}{productNoise}}


Use the preceding square-integrability proof to form the class of \(((x,y),s)\mapsto W_n(s)\) on the concrete extension. Both the numerical chain and reference use this same \(L^2\) innovation.

\subsection{\texorpdfstring{product\allowbreak{}Noise coe}{productNoise coe}}


The representative selected by \texttt{to\allowbreak{}Lp} agrees almost everywhere with the pulled-back original noise \(W_n(s)\). This identity connects quotient-space noise arithmetic to the actual source random functions.

\subsection{\texorpdfstring{product\allowbreak{}Output\allowbreak{}Law}{productOutputLaw}}


Push \(\mu\otimes Q\) through the pointwise noise-driven chain starting at \(x\) with source path \(W(s)\). This definition omits the reference coupling entirely, making explicit the common algorithm output law for all coupling realizations.

\subsection{\texorpdfstring{product chain coe}{product chain coe}}


Induct through the lifted chain and use the representative identities for initial state, product noise, lifted drift, addition, and scalar multiplication. At each fixed index the result agrees almost everywhere with the ordinary recursion driven by \(x\) and \(W(s)\).

\subsection{\texorpdfstring{product chain output\allowbreak{}Law}{product chain outputLaw}}


Replace the \(L^2\) chain by its almost-everywhere pointwise recursion. That recursion depends only on the algorithm's initial coordinate and noise sample, so \texttt{initial\allowbreak{}Extension\_\allowbreak{}map} gives its coupling-independent law. The result is exactly \texttt{product\allowbreak{}Output\allowbreak{}Law}.

\subsection{\texorpdfstring{product\allowbreak{}Probability}{productProbability}}


A coupling of the initial probability with the reference measure is itself a probability, and its product with the fixed probability noise source is a probability. This local instance supplies normalization on every extension separately.

\subsection{\texorpdfstring{Product\allowbreak{}Reference}{ProductReference}}


For every finite-cost initial coupling, the structure supplies a concrete \texttt{L2\allowbreak{}Reference} on that coupling's product extension, with its initial state matching the second coordinate. Algorithm output laws and initial costs are proved outside the structure; no convergence estimate is a field.

\subsection{\texorpdfstring{Product\allowbreak{}Reference.\allowbreak{}output finite}{ProductReference.output finite}}


Finite initial Wasserstein distance supplies at least one finite-cost initial coupling. On its concrete extension, the numerical and reference variables belong to \(L^2\), giving a finite-cost output coupling. Their established laws identify its marginals with the common output law and \(\pi\), hence the output Wasserstein distance is finite.

\subsection{\texorpdfstring{dalalyan small step of product reference}{dalalyan small step of product reference}}


For each finite-cost initial coupling \(c\), apply the analytic \(L^2\) Dalalyan theorem on \(c\otimes Q\). The initial class-distance theorem identifies its initial error exactly with the real coupling cost, and the common output-law theorem identifies its numerical law.



The resulting bound holds for every finite \(c\), with coefficient \((1-mh)^n\ge0\). Infimizing their costs gives the genuine \(W_2(\mu,\pi)\), without selecting an optimal coupling. Output finiteness is proved first so the final real-valued inequality cannot conceal infinity.

\subsection{\texorpdfstring{dalalyan large step of product reference}{dalalyan large step of product reference}}


Repeat the couplingwise analytic argument for the stable large-step branch, whose contraction factor is \(Lh-1\in[0,1)\). The geometric error sum has denominator \(2-Lh>0\). The finite-coupling infimum theorem replaces the initial class distance by the actual initial Wasserstein distance and yields the displayed large-step bias term.

\section{\texorpdfstring{Coordinate Gaussian moments and L\ensuremath{{}^{2}} Brownian continuity}{Coordinate Gaussian moments and L² Brownian continuity}}
\label{proof-module-UnadjustedLangevin-GaussianCoordinates-lean}
{\small\color{gray}Source module: \texttt{Unadjusted\allowbreak{}Langevin/\allowbreak{}Gaussian\allowbreak{}Coordinates.\allowbreak{}lean}.\par}


Coordinatewise Gaussian laws suffice for the second moment \(\mathbb E\|X\|^2=dv\). This module carefully obtains Brownian \(L^2\) continuity and increment norms without assuming independence between spatial coordinates.

\subsection{\texorpdfstring{gaussian coordinates mem\allowbreak{}Lp}{gaussian coordinates memLp}}


Each centered Gaussian coordinate has a finite second moment. In finite Euclidean dimension, \(\|X\|^2=\sum_iX_i^2\), so the finite sum is integrable. Coordinatewise measurability also assembles to vector measurability, giving \(X\in L^2\). Coordinate independence is unnecessary for this moment statement.

\subsection{\texorpdfstring{gaussian coordinates second moment}{gaussian coordinates second moment}}


Use \(\|X\|^2=\sum_iX_i^2\) and finite-sum integral linearity. The coordinate law \(X_i\sim\mathcal N(0,v)\) gives \(\mathbb EX_i^2=v\), hence \(\mathbb E\|X\|^2=|\iota|v\). The dimension factor is exactly the cardinality of the finite coordinate set.

\subsection{\texorpdfstring{gaussian coordinates Lp norm sq}{gaussian coordinates Lp norm sq}}


For the given \(L^2\) class, its squared norm is the integral of the squared norm of its representative. Apply the preceding Gaussian coordinate moment formula to obtain \(\|X\|_2^2=|\iota|v\).

\subsection{\texorpdfstring{gaussian coordinates Lp norm}{gaussian coordinates Lp norm}}


Take nonnegative square roots in \(\|X\|_2^2=|\iota|v\). Norm nonnegativity selects the positive root, giving \(\|X\|_2=\sqrt{|\iota|v}\), including zero variance and zero dimension.

\subsection{\texorpdfstring{gaussian coordinates to\allowbreak{}Lp coordinate\allowbreak{}Law}{gaussian coordinates toLp coordinateLaw}}


The representative of \texttt{Mem\allowbreak{}Lp.\allowbreak{}to\allowbreak{}Lp X} agrees almost everywhere with \(X\). Coordinate evaluation preserves that equality, so each coordinate law transfers to the canonical \(L^2\) representative by almost-everywhere invariance of pushforwards.

\subsection{\texorpdfstring{gaussian coordinates to\allowbreak{}Lp norm}{gaussian coordinates toLp norm}}


Apply the exact \(L^2\) Gaussian norm theorem to \texttt{to\allowbreak{}Lp X}, using its transferred coordinate laws. This gives \(\|[X]\|_2=\sqrt{|\iota|v}\) for the class constructed from the original random vector.

\subsection{\texorpdfstring{gaussian coordinates scaled l2 norm}{gaussian coordinates scaled l2 norm}}


The scalar-norm identity gives \(\|\sqrt2[X]\|_2=\sqrt2\sqrt{|\iota|v}\). Nonnegativity permits the square-root product rules, producing \(\sqrt{2|\iota|}\sqrt v\). This is the precise Langevin noise constant used in the local error estimate.

\subsection{\texorpdfstring{brownian\allowbreak{}Coordinates}{brownianCoordinates}}


Assemble the scalar coordinate processes into \(B_t(\omega)=(W_i(t,\omega))_{i\in\iota}\). Finite Euclidean coordinates determine the vector; no new random source or independence is introduced by this definition.

\subsection{\texorpdfstring{brownian\allowbreak{}Coordinates coordinate\allowbreak{}Law}{brownianCoordinates coordinateLaw}}


At a fixed time \(t\ge0\), each scalar pre-Brownian coordinate has law \(\mathcal N(0,t)\). Unfolding the assembled vector shows that its \(i\)-th coordinate is exactly that scalar variable, giving the required coordinate law.

\subsection{\texorpdfstring{brownian\allowbreak{}Coordinates\allowbreak{}Increment}{brownianCoordinatesIncrement}}


Define the vector increment \(\Delta_aB_t=B_{a+t}-B_a\) coordinate by coordinate. This records both the starting time and elapsed nonnegative time, matching the local noise input on an arbitrary Langevin segment.

\subsection{\texorpdfstring{brownian\allowbreak{}Coordinates\allowbreak{}Increment coordinate\allowbreak{}Law}{brownianCoordinatesIncrement coordinateLaw}}


The scalar Brownian increment law depends only on elapsed time: \(W_i(a+t)-W_i(a)\sim\mathcal N(0,t)\). Simplify \((a+t)-a=t\) in the variance and unfold the vector increment to obtain its coordinate law.

\subsection{\texorpdfstring{brownian\allowbreak{}Coordinates\allowbreak{}Increment mem\allowbreak{}Lp}{brownianCoordinatesIncrement memLp}}


Apply \texttt{gaussian\_\allowbreak{}coordinates\_\allowbreak{}mem\allowbreak{}Lp} to the increment's coordinate laws of variance \(t\). The assembled Brownian vector increment is square integrable even without independence between different coordinate processes.

\subsection{\texorpdfstring{brownian\allowbreak{}Coordinates\allowbreak{}Increment scaled l2 norm}{brownianCoordinatesIncrement scaled l2 norm}}


Apply the exact scaled Gaussian-coordinate norm theorem with variance \(t\). It yields \(\|\sqrt2[\Delta_aB_t]\|_2=\sqrt{2|\iota|}\sqrt t\), establishing the noise radius from Brownian laws alone.

\subsection{\texorpdfstring{brownian\allowbreak{}Coordinates\allowbreak{}Lp}{brownianCoordinatesLp}}


Convert the assembled Brownian vector at time \(t\) to its \(L^2\) class using its Gaussian-coordinate square-integrability proof. This places the time-indexed process in the Hilbert space where the deterministic integral equation is solved.

\subsection{\texorpdfstring{brownian\allowbreak{}Coordinates\allowbreak{}Lp coe}{brownianCoordinatesLp coe}}


The \texttt{to\allowbreak{}Lp} representative theorem gives \([B_t](\omega)=B_t(\omega)\) almost everywhere for each fixed time. This is the bridge from class-valued continuity statements to coordinate laws of the original Brownian process.

\subsection{\texorpdfstring{brownian\allowbreak{}Coordinates\allowbreak{}Lp sub coordinate\allowbreak{}Law}{brownianCoordinatesLp sub coordinateLaw}}


The difference of two selected representatives agrees almost everywhere with \(B_t-B_s\). Brownian increment laws, including the sign reversal when \(t<s\), give a centered Gaussian with variance \(|t-s|\); centered Gaussian symmetry handles the reversed order.

\subsection{\texorpdfstring{brownian\allowbreak{}Coordinates\allowbreak{}Lp norm sub}{brownianCoordinatesLp norm sub}}


Apply the Gaussian \(L^2\) norm formula to the difference class, whose coordinates have variance \(|t-s|\). Thus \(\|[B_t]-[B_s]\|_2=\sqrt{|\iota|\,|t-s|}\). This exact increment norm is stronger than the continuity statement derived from it.

\subsection{\texorpdfstring{brownian\allowbreak{}Coordinates\allowbreak{}Lp continuous}{brownianCoordinatesLp continuous}}


As \(s\to t\), \(|t-s|\to0\), hence the exact norm formula gives \(\|[B_t]-[B_s]\|_2\to0\). Lean uses continuity of the square root and the metric characterization of continuity. This needs finite-dimensional Brownian laws, with no sample-path continuity premise.

\subsection{\texorpdfstring{brownian\allowbreak{}Coordinates\allowbreak{}Lp increment scaled norm}{brownianCoordinatesLp increment scaled norm}}


Set the two times in the difference-norm identity to \(a+t\) and \(a\). Their distance is \(t\), and scalar multiplication by \(\sqrt2\) gives \(\sqrt{2|\iota|}\sqrt t\). This is the class-valued version of the earlier increment estimate.

\subsection{\texorpdfstring{brownian\allowbreak{}Coordinates\allowbreak{}Lp zero}{brownianCoordinatesLp zero}}


At time zero every coordinate has the degenerate centered Gaussian law of variance zero. Equivalently its \(L^2\) norm is zero by the moment formula. Norm-zero equivalence proves that the whole class is zero, irrespective of the chosen pointwise representative.

\subsection{\texorpdfstring{brownian\allowbreak{}Coordinates\allowbreak{}Real\allowbreak{}Lp}{brownianCoordinatesRealLp}}


Extend the nonnegative-time \(L^2\) process to real times by \(t\mapsto[B_{\max(t,0)}]\), implemented with \texttt{Real.\allowbreak{}to\allowbreak{}NNReal}. It is constant zero for negative times and agrees with the original class at nonnegative times.

\subsection{\texorpdfstring{brownian\allowbreak{}Coordinates\allowbreak{}Real\allowbreak{}Lp continuous}{brownianCoordinatesRealLp continuous}}


The clipping map \(t\mapsto\max(t,0)\) into nonnegative reals is continuous. Compose it with the established \(L^2\)-continuous Brownian map to obtain continuity of the real-time extension, including at zero.

\subsection{\texorpdfstring{brownian\allowbreak{}Coordinates\allowbreak{}Real\allowbreak{}Lp zero}{brownianCoordinatesRealLp zero}}


At real time zero the clipping map returns nonnegative zero. The previously proved zero-time \(L^2\) identity therefore simplifies the real-time extension to zero.

\subsection{\texorpdfstring{brownian\allowbreak{}Coordinates\allowbreak{}Real\allowbreak{}Lp increment scaled norm}{brownianCoordinatesRealLp increment scaled norm}}


For \(a,t\ge0\), clipping leaves \(a\) and \(a+t\) unchanged. Rewrite the real-time extension at both endpoints as the nonnegative-time Brownian classes, then apply their scaled increment formula. This supplies \(\sqrt{2d}\sqrt t\) in the real interval-integral formulation.

\section{\texorpdfstring{Construct the L\ensuremath{{}^{2}} Langevin equation from Brownian-\allowbreak{}coordinate laws}{Construct the L² Langevin equation from Brownian-coordinate laws}}
\label{proof-module-UnadjustedLangevin-BrownianLangevin-lean}
{\small\color{gray}Source module: \texttt{Unadjusted\allowbreak{}Langevin/\allowbreak{}Brownian\allowbreak{}Langevin.\allowbreak{}lean}.\par}


Coordinate pre-Brownian laws produce the continuous \(L^2\) additive noise needed by the deterministic integral-equation existence theorem. The exact noise norm is \(\sqrt{2d}\sqrt t\).

\subsection{\texorpdfstring{brownian\allowbreak{}Langevin\allowbreak{}Noise}{brownianLangevinNoise}}


Define \(N_a(t)=\sqrt2([B_{a+t}]-[B_a])\) in \(L^2\), using the real-time Brownian extension. The subtraction fixes the input at zero at local time zero, and the \(\sqrt2\) factor gives the overdamped Langevin diffusion normalization.

\subsection{\texorpdfstring{brownian\allowbreak{}Langevin\allowbreak{}Noise continuous}{brownianLangevinNoise continuous}}


The real-time Brownian map is continuous in \(L^2\) from its Gaussian increment moments. Composition with \(t\mapsto a+t\), subtraction of the fixed class \([B_a]\), and scalar multiplication preserve continuity. Thus the additive input satisfies the deterministic Picard theorem's continuity requirement.

\subsection{\texorpdfstring{brownian\allowbreak{}Langevin\allowbreak{}Noise zero}{brownianLangevinNoise zero}}


At local time zero both Brownian terms are \([B_a]\), so they cancel and their scalar multiple is zero. Lean unfolds the definition and uses additive identities; no exceptional-set statement is needed because this is equality in \(L^2\).

\subsection{\texorpdfstring{brownian\allowbreak{}Langevin\allowbreak{}Noise norm}{brownianLangevinNoise norm}}


For \(a,t\ge0\), apply the exact real-time Brownian increment moment formula. It gives \(\|N_a(t)\|_2=\sqrt{2d}\sqrt t\), with \(d=|\iota|\). The Gaussian laws prove this value; it is not an additional noise-bound hypothesis.

\subsection{\texorpdfstring{exists lp langevin integral of brownian small}{exists lp langevin integral of brownian small}}


Supply the continuous, zero-starting \(N_a\) to \texttt{exists\_\allowbreak{}lp\_\allowbreak{}langevin\_\allowbreak{}integral\_\allowbreak{}small}. It constructs \(Y_t-Y_0=-\int_0^t\operatorname{lpDrift}(\nabla U)(Y_s)ds+N_a(t)\) for \(Lh\le1/2\). Attach the exact noise norm theorem to the existential package. Neither the integral equation nor its noise estimate is assumed.

\subsection{\texorpdfstring{exists lp langevin integral of brownian}{exists lp langevin integral of brownian}}


Use the unrestricted finite-horizon \(L^2\) existence theorem with the same constructed Brownian input. The resulting solution is continuous in \(L^2\), starts at the prescribed random variable, and satisfies the actual gradient integral equation. The exact \(\sqrt{2d}\sqrt t\) noise formula is included as a proved property.

\section{\texorpdfstring{Gibbs normalization and second moments}{Gibbs normalization and second moments}}
\label{proof-module-UnadjustedLangevin-Gibbs-lean}
{\small\color{gray}Source module: \texttt{Unadjusted\allowbreak{}Langevin/\allowbreak{}Gibbs.\allowbreak{}lean}.\par}


On a finite-dimensional real Hilbert space, differentiability and \(m\)-strong convexity with \(m>0\) imply that \(\pi(dx)=Z^{-1}e^{-U(x)}dx\) is a probability measure with finite second moment. This establishes existence of the target measure used throughout the coupling proof.

\subsection{\texorpdfstring{gibbs\allowbreak{}Weight}{gibbsWeight}}


The unnormalized density is \(w_U(x)=e^{-U(x)}\), embedded in \([0,\infty]\) with \texttt{ENNReal.\allowbreak{}of\allowbreak{}Real}. This choice lets the partition integral retain the value infinity; positivity of the exponential ensures that embedding does not discard any part of the density.

\subsection{\texorpdfstring{gibbs\allowbreak{}Partition}{gibbsPartition}}


The partition function is \(Z_U=\int e^{-U(x)}\,dx\), formed as a nonnegative extended integral. At this definition no finiteness assumption is made. Later normalization uses separate proofs of \(0<Z_U\) and \(Z_U<\infty\), which are needed to justify cancellation with its inverse.

\subsection{\texorpdfstring{gibbs\allowbreak{}Measure}{gibbsMeasure}}


Define \(\pi_U(dx)=Z_U^{-1}e^{-U(x)}dx\) using \texttt{volume.\allowbreak{}with\allowbreak{}Density}. The inverse is the extended nonnegative inverse. The definition therefore exists for every potential, while the later probability theorem proves that it has total mass one under the stated strong convexity assumptions.

\subsection{\texorpdfstring{integrable exp neg mul norm sq}{integrable exp neg mul norm sq}}


For \(a>0\), \(e^{-a\|x\|^2}\) is Lebesgue integrable in finite dimension. The Lean proof specializes \texttt{Gaussian\allowbreak{}Fourier.\allowbreak{}integrable\_\allowbreak{}cexp\_\allowbreak{}neg\_\allowbreak{}mul\_\allowbreak{}sq\_\allowbreak{}norm\_\allowbreak{}add} to real positive precision \(a\) and zero linear term, then takes the real part. The complex Gaussian integrand is the real embedding of the desired exponential, so the final simplification identifies the functions exactly. This also includes dimension zero.

\subsection{\texorpdfstring{strong\allowbreak{}Convex quadratic lower bound}{strongConvex quadratic lower bound}}


Apply the strong supporting-plane inequality at zero: \(U(x)\ge U(0)+\langle\nabla U(0),x\rangle+\frac m2\|x\|^2\). Cauchy–Schwarz and \(ab\le a^2/m+(m/4)b^2\), valid for \(m>0\), absorb the linear term. The result is \(U(x)\ge U(0)-\|\nabla U(0)\|^2/m+(m/4)\|x\|^2\). Lean isolates this scalar square inequality before its final arithmetic rearrangement.

\subsection{\texorpdfstring{integrable exp neg of strong\allowbreak{}Convex}{integrable exp neg of strongConvex}}


Exponentiating the quadratic lower bound gives \(e^{-U(x)}\le e^{-U(0)+\|\nabla U(0)\|^2/m}e^{-(m/4)\|x\|^2}\). The right side is an integrable Gaussian multiple. Differentiability supplies continuity and measurability of the left side, so \texttt{Integrable.\allowbreak{}mono'} transfers integrability from that explicit dominating function.

\subsection{\texorpdfstring{integrable norm sq mul exp neg mul norm sq}{integrable norm sq mul exp neg mul norm sq}}


For \(a>0\), split the Gaussian decay into two equal factors. The scalar function \(s\mapsto s e^{-(a/2)s}\) is bounded on \([0,\infty)\), so \(\|x\|^2e^{-a\|x\|^2}\le C_a e^{-(a/2)\|x\|^2}\). The remaining Gaussian is integrable; the proof checks measurability and nonnegativity before applying domination.

\subsection{\texorpdfstring{integrable norm sq mul exp neg of strong\allowbreak{}Convex}{integrable norm sq mul exp neg of strongConvex}}


Multiply the previous Gibbs Gaussian domination by \(\|x\|^2\). Its right side has the form \(C\|x\|^2e^{-(m/4)\|x\|^2}\), whose integrability was proved separately. The pointwise inequality uses \(\|x\|^2\ge0\), so multiplying it preserves order; this is the second-moment input for the normalized measure.

\subsection{\texorpdfstring{gibbs\allowbreak{}Partition pos}{gibbsPartition pos}}


The continuous function \(e^{-U}\) is strictly positive everywhere. Its extended integral cannot vanish: a zero integral of a nonnegative measurable function would force it to vanish almost everywhere, contradicting positivity on a nonzero-volume open set. This proves \(0<Z_U\) without requiring an upper bound on \(Z_U\).

\subsection{\texorpdfstring{gibbs\allowbreak{}Partition ne top}{gibbsPartition ne top}}


Strong convexity with \(m>0\) has already given ordinary integrability of \(e^{-U}\). For a nonnegative integrable function its extended integral is finite. Unfolding \texttt{gibbs\allowbreak{}Partition} identifies that integral exactly with \(Z_U\), giving the guard needed for inverse cancellation.

\subsection{\texorpdfstring{gibbs\allowbreak{}Partition eq of\allowbreak{}Real integral}{gibbsPartition eq ofReal integral}}


Nonnegativity and integrability identify the extended integral with the embedded real integral: \(Z_U=\operatorname{ofReal}(\int e^{-U}\,dx)\). The proof uses the genuine integral/\allowbreak{}lintegral comparison theorem, so it cannot identify an infinite partition with the default real value of an undefined integral.

\subsection{\texorpdfstring{integral exp neg pos}{integral exp neg pos}}


Combine \(0<Z_U\) with \(Z_U=\operatorname{ofReal}(\int e^{-U}\,dx)\). Positivity of \texttt{of\allowbreak{}Real} is equivalent to positivity of its real argument, proving that the ordinary normalizing constant is strictly positive and allowing real division by it later.

\subsection{\texorpdfstring{is\allowbreak{}Probability\allowbreak{}Measure gibbs\allowbreak{}Measure}{isProbabilityMeasure gibbsMeasure}}


The mass of a density measure is its extended density integral. Pull the constant \(Z_U^{-1}\) out to get \(\pi_U(H)=Z_U^{-1}Z_U\). The proofs that \(Z_U\ne0\) and \(Z_U\ne\infty\) justify this product being one. This theorem constructs the \texttt{Is\allowbreak{}Probability\allowbreak{}Measure} instance rather than assuming normalization.

\subsection{\texorpdfstring{integrable norm sq gibbs\allowbreak{}Measure}{integrable norm sq gibbsMeasure}}


Changing measure converts integrability under \(\pi_U\) to integrability of \(Z_U^{-1}e^{-U(x)}\|x\|^2\) under volume. The weighted second moment is already integrable, and the real normalizing factor is a finite scalar. The with-density theorem supplies the almost-everywhere measurability and finiteness guards for this conversion.

\subsection{\texorpdfstring{mem\allowbreak{}Lp id gibbs\allowbreak{}Measure}{memLp id gibbsMeasure}}


For the identity random variable, membership in \(L^2(\pi_U;H)\) is equivalent to integrability of \(\|x\|^2\). The identity is measurable, and the preceding weighted-moment theorem supplies the integral. Lean uses the \(p=2\) characterization of \texttt{Mem\allowbreak{}Lp}, retaining the actual function before passing to any quotient class.

\subsection{\texorpdfstring{mem\allowbreak{}Lp sub const gibbs\allowbreak{}Measure}{memLp sub const gibbsMeasure}}


Since \(\pi_U\) is a probability, the constant map \(x\mapsto z\) belongs to \(L^2\). Subtract it from the identity, already in \(L^2(\pi_U)\), to obtain \(x\mapsto x-z\in L^2(\pi_U)\). This is a moment assertion for every chosen center, without requiring that the center minimize the potential.

\section{\texorpdfstring{Subgaussian Gibbs score and its sharp second moment}{Subgaussian Gibbs score and its sharp second moment}}
\label{proof-module-UnadjustedLangevin-GibbsScore-lean}
{\small\color{gray}Source module: \texttt{Unadjusted\allowbreak{}Langevin/\allowbreak{}Gibbs\allowbreak{}Score.\allowbreak{}lean}.\par}


Differentiability, an \(L\)-Lipschitz gradient, and finite Gibbs normalization suffice for \(\int\|\nabla U\|^2d\pi\le Ld\). Strong convexity is used only by the final wrapper to derive normalization.

\subsection{\texorpdfstring{gibbs\allowbreak{}Weight mul exp score le}{gibbsWeight mul exp score le}}


The descent lemma gives \(U(x+v)\le U(x)+\langle\nabla U(x),v\rangle+(L/2)\|v\|^2\). Negate and exponentiate, then rearrange to \(e^{-U(x)}e^{-\langle\nabla U(x),v\rangle}\le e^{L\|v\|^2/2}e^{-U(x+v)}\). Lean first establishes the real inequality and then embeds it in the extended nonnegative numbers.

\subsection{\texorpdfstring{gibbs score lmgf le}{gibbs score lmgf le}}


Integrate the pointwise exponential-score bound. Translation invariance gives \(\int e^{-U(x+v)}dx=Z_U\); normalization cancels this factor, leaving \(\mathbb E_{\pi_U}e^{-\langle\nabla U,v\rangle}\le e^{L\|v\|^2/2}\). Finiteness of \(Z_U\) is explicit because extended inverse cancellation needs it.

\subsection{\texorpdfstring{gibbs score subgaussian}{gibbs score subgaussian}}


Apply the preceding bound with the vector \(-tv\), for every real \(t\). The exponent becomes \(t\langle\nabla U,v\rangle\) and the right side is \(e^{t^2L\|v\|^2/2}\). This is exactly \texttt{Has\allowbreak{}Subgaussian\allowbreak{}MGF} with variance proxy \(L\|v\|^2\); the proof also establishes exponential integrability required by that structure.

\subsection{\texorpdfstring{sq le exp add exp neg sub two}{sq le exp add exp neg sub two}}


The series \(\cosh a=\sum_{n\ge0}a^{2n}/(2n)!\) has nonnegative terms, so retaining its first two terms gives \(\cosh a\ge1+a^2/2\). Substitute \(2\cosh a=e^a+e^{-a}\) and rearrange to \(a^2\le e^a+e^{-a}-2\). Lean uses \texttt{Real.\allowbreak{}has\allowbreak{}Sum\_\allowbreak{}cosh} with \texttt{sum\_\allowbreak{}le\_\allowbreak{}has\allowbreak{}Sum}, explicitly proving nonnegativity of every discarded even-power term.

\subsection{\texorpdfstring{integral sq le of subgaussian}{integral sq le of subgaussian}}


Integrate \(t^2X^2\le e^{tX}+e^{-tX}-2\), giving \(t^2\mathbb EX^2\le2(e^{ct^2/2}-1)\). Subgaussianity provides the exponential integrability and hence the square integrability needed here. Divide by \(t^2\) for \(t\ne0\), and let \(t\to0\) to obtain \(\mathbb EX^2\le c\). Lean computes the scalar limit before applying closedness of order.

\subsection{\texorpdfstring{integral norm sq gradient gibbs le}{integral norm sq gradient gibbs le}}


Choose an orthonormal basis \((e_i)_{i=1}^d\) and use Parseval: \(\|\nabla U(x)\|^2=\sum_i\langle\nabla U(x),e_i\rangle^2\). Each coordinate is subgaussian with proxy \(L\), hence has second moment at most \(L\). Finite-sum integrability permits integration term by term, yielding \(\mathbb E_{\pi_U}\|\nabla U\|^2\le Ld\).

\subsection{\texorpdfstring{mem\allowbreak{}Lp gradient gibbs\allowbreak{}Measure}{memLp gradient gibbsMeasure}}


The coordinate subgaussian bounds prove integrability of every squared gradient coordinate. Summing over the finite orthonormal basis proves integrability of \(\|\nabla U\|^2\). Lipschitz continuity provides strong measurability, and the \texttt{Mem\allowbreak{}Lp} square criterion packages those facts as \(\nabla U\in L^2(\pi_U)\).

\subsection{\texorpdfstring{integral norm sq gradient gibbs le of strong\allowbreak{}Convex}{integral norm sq gradient gibbs le of strongConvex}}


Strong convexity with \(m>0\) supplies the finite-partition premise of the score theorem. Substitute that proof into the first-order smoothness result to obtain \(\mathbb E_{\pi_U}\|\nabla U\|^2\le Ld\). The strong convexity constant is used for normalization here; it does not appear in the sharp score bound.

\section{\texorpdfstring{Transfer the Gibbs score estimate to an L\ensuremath{{}^{2}} reference}{Transfer the Gibbs score estimate to an L² reference}}
\label{proof-module-UnadjustedLangevin-StationaryMoments-lean}
{\small\color{gray}Source module: \texttt{Unadjusted\allowbreak{}Langevin/\allowbreak{}Stationary\allowbreak{}Moments.\allowbreak{}lean}.\par}


A reference variable \(Y\) whose law is \(\pi\) inherits \(\|\nabla U(Y)\|_{L^2}\le\sqrt{Ld}\). The assertion that its law is Gibbs is a displayed premise here.

\subsection{\texorpdfstring{lp\allowbreak{}Drift norm le of gibbs law}{lpDrift norm le of gibbs law}}


The representative of \texttt{lp\allowbreak{}Drift} equals \(\nabla U(Y)\) almost everywhere. Change variables using \(P\circ Y^{-1}=\pi_U\), so its squared \(L^2(P)\) norm is \(\int\|\nabla U\|^2d\pi_U\le Ld\). Nonnegativity of the norm and of \(Ld\) then gives \(\|\nabla U(Y)\|_2\le\sqrt{Ld}\). The law identity, rather than independence, is the input to this transfer.

\subsection{\texorpdfstring{lp\allowbreak{}Drift norm le of strong\allowbreak{}Convex gibbs law}{lpDrift norm le of strongConvex gibbs law}}


Use strong convexity to prove \(Z_U<\infty\), and apply the previous law-transfer theorem with that proof. The conclusion is the same \(\sqrt{Ld}\) bound for the lifted drift of any \(L^2\) variable with Gibbs law. There is no separate assumed drift-moment bound hidden in this specialization.

\section{\texorpdfstring{Assemble continuous Brownian reference segments}{Assemble continuous Brownian reference segments}}
\label{proof-module-UnadjustedLangevin-BrownianReference-lean}
{\small\color{gray}Source module: \texttt{Unadjusted\allowbreak{}Langevin/\allowbreak{}Brownian\allowbreak{}Reference.\allowbreak{}lean}.\par}


Construct consecutive \(L^2\) Langevin segments from actual Brownian-coordinate noise and package them as an analytic reference. At this intermediate stage Gibbs stationarity remains an explicit premise, later discharged in CanonicalULA.

\subsection{\texorpdfstring{brownian\allowbreak{}Segment\allowbreak{}From}{brownianSegmentFrom}}


Choose the integral solution supplied by \texttt{exists\_\allowbreak{}lp\_\allowbreak{}langevin\_\allowbreak{}integral\_\allowbreak{}of\_\allowbreak{}brownian} for the segment starting at absolute time \(nh\), with the supplied initial \(L^2\) variable. Classical choice selects data from a proved existence theorem; its defining equation and continuity are recovered by the specification theorem.

\subsection{\texorpdfstring{brownian\allowbreak{}Segment\allowbreak{}From spec}{brownianSegmentFrom spec}}


Apply the choice specification to obtain continuity, the time-zero initial value, the gradient integral equation on \([0,h]\), and the exact noise norm. These four properties describe the selected segment independently of its noncomputable implementation.

\subsection{\texorpdfstring{brownian\allowbreak{}Knots}{brownianKnots}}


Set the zeroth knot to the prescribed initial variable. Define the next knot to be the endpoint at local time \(h\) of the segment started at the current knot. This recursion ensures consecutive segments meet exactly as \(L^2\) classes.

\subsection{\texorpdfstring{brownian\allowbreak{}Segments}{brownianSegments}}


At index \(n\), select \texttt{brownian\allowbreak{}Segment\allowbreak{}From} with the recursively constructed \(n\)-th knot as initial state. The absolute noise start is \(nh\), so all segments are driven by increments of the same Brownian process.

\subsection{\texorpdfstring{brownian\allowbreak{}Segments start}{brownianSegments start}}


The chosen-segment specification states its value at zero equals its supplied initial state. Here that state is the \(n\)-th knot, so \(Y^{(n)}_0=K_n\) follows directly by specialization.

\subsection{\texorpdfstring{brownian\allowbreak{}Segments initial}{brownianSegments initial}}


Apply the segment-start identity at \(n=0\), then unfold the zeroth knot. The actual constructed reference therefore begins at exactly the prescribed random variable, preserving its chosen joint realization with the numerical initialization.

\subsection{\texorpdfstring{brownian\allowbreak{}Euler\allowbreak{}Noise}{brownianEulerNoise}}


Normalize the actual segment endpoint noise by \(1/\sqrt{2h}\). This defines the \(L^2\) innovation used by the Euler chain, so it shares the exact Brownian increment driving the reference segment.

\subsection{\texorpdfstring{brownian\allowbreak{}Euler\allowbreak{}Noise endpoint}{brownianEulerNoise endpoint}}


For \(h>0\), \(\sqrt{2h}\ne0\). Multiply the normalized noise class by \(\sqrt{2h}\) and cancel the inverse scalar. The resulting equality identifies the Euler noise contribution with the same Brownian endpoint input in the integral equation.

\subsection{\texorpdfstring{brownian\allowbreak{}Reference}{brownianReference}}


Populate \texttt{L2\allowbreak{}Reference} with the constructed segments. Their specification proves continuity and the integral equation; recursive knots prove endpoint compatibility; normalized noise proves endpoint matching; Gibbs-score moment transfer proves the drift bound. The supplied all-time Gibbs-law premise remains visible and is later discharged by the stationarity construction.

\subsection{\texorpdfstring{brownian\allowbreak{}Reference of strong\allowbreak{}Convex}{brownianReference of strongConvex}}


Use strong convexity to derive finiteness of the Gibbs partition, then call \texttt{brownian\allowbreak{}Reference}. All analytic, normalization, moment, and noise fields are proved internally. The all-time law premise is still explicit in this constructor and must be provided by the stationary-law theorem.

\section{\texorpdfstring{Strong convexity forces m \ensuremath{\le} L}{Strong convexity forces m ≤ L}}
\label{proof-module-UnadjustedLangevin-PotentialAssumptions-lean}
{\small\color{gray}Source module: \texttt{Unadjusted\allowbreak{}Langevin/\allowbreak{}Potential\allowbreak{}Assumptions.\allowbreak{}lean}.\par}


On a nontrivial Hilbert space the numerical compatibility \(m\le L\) follows from strong convexity and a Lipschitz gradient of the same differentiable potential.

\subsection{\texorpdfstring{Compare monotonicity with Cauchy–Schwarz}{Compare monotonicity with Cauchy–Schwarz}}


Choose \(x\ne0\), possible because the Hilbert space is nontrivial. Strong monotonicity gives \(m\|x\|^2\le\langle\nabla U(x)-\nabla U(0),x\rangle\), while Cauchy–Schwarz and Lipschitzness bound the same quantity by \(L\|x\|^2\). Cancel the strictly positive \(\|x\|^2\) to prove \(m\le L\). No positivity premise on \(m\) is needed here.

\section{\texorpdfstring{Normalized increments are independent standard Gaussians}{Normalized increments are independent standard Gaussians}}
\label{proof-module-UnadjustedLangevin-BrownianInnovations-lean}
{\small\color{gray}Source module: \texttt{Unadjusted\allowbreak{}Langevin/\allowbreak{}Brownian\allowbreak{}Innovations.\allowbreak{}lean}.\par}


For a pre-Brownian coordinate and positive grid spacing \(h\), \(Z_n=(B_{(n+1)h}-B_{nh})/\sqrt h\) has standard normal law. Across independent coordinate paths, the entire coordinate-time family is independent.

\subsection{\texorpdfstring{normalized\allowbreak{}Brownian\allowbreak{}Increment}{normalizedBrownianIncrement}}


Define \(\xi_n=(B_{(n+1)h}-B_{nh})/\sqrt h\). The denominator is interpreted as real scalar division; positive-step law results below explicitly require \(h>0\), so that the variance rescaling is legitimate.

\subsection{\texorpdfstring{normalized\allowbreak{}Brownian\allowbreak{}Increment has\allowbreak{}Law}{normalizedBrownianIncrement hasLaw}}


A Brownian increment over an interval of length \(h\) is \(\mathcal N(0,h)\). Scaling by \(1/\sqrt h\) multiplies its variance by \(1/h\), producing \(\mathcal N(0,1)\). The proof rewrites \((n+1)h-nh=h\) and uses \((\sqrt h)^2=h\), justified by \(h>0\).

\subsection{\texorpdfstring{normalized\allowbreak{}Brownian\allowbreak{}Increment mem\allowbreak{}Lp}{normalizedBrownianIncrement memLp}}


The exact standard Gaussian law provides every finite moment, in particular the second. Pull this square integrability back through \texttt{normalized\allowbreak{}Brownian\allowbreak{}Increment\_\allowbreak{}has\allowbreak{}Law} to prove \(\xi_n\in L^2(P)\).

\subsection{\texorpdfstring{normalized\allowbreak{}Brownian\allowbreak{}Increment second moment}{normalizedBrownianIncrement second moment}}


Transfer the integral of \(x^2\) along the standard Gaussian law of \(\xi_n\). Its mean is zero and variance one, hence \(\mathbb E\xi_n^2=1\). This is an exact second-moment identity for the normalized increment.

\subsection{\texorpdfstring{normalized\allowbreak{}Brownian\allowbreak{}Increment independent}{normalizedBrownianIncrement independent}}


Brownian increments on the consecutive grid intervals are independent. Normalizing each by the same deterministic scalar is a coordinatewise measurable transformation and preserves independence. The formulation permits \(h=0\), where every increment is the constant zero and the independence assertion remains valid.

\subsection{\texorpdfstring{i\allowbreak{}Indep\allowbreak{}Fun uncurry nat of aemeasurable}{iIndepFun uncurry nat of aemeasurable}}


Assume the row-valued random maps \((X_{i,n})_n\) are independent and each row is internally independent. Replace the almost-everywhere measurable coordinates by measurable versions; countability of \(\mathbb N\) lets all replacements in a row hold on one full-measure set. Apply independence of independent families, then transfer the result back across the coordinatewise almost-everywhere equalities.

\subsection{\texorpdfstring{normalized\allowbreak{}Brownian\allowbreak{}Increments independent}{normalizedBrownianIncrements independent}}


Each grid increment depends measurably on two entries of its coordinate grid path. Independence of those paths gives independence across coordinate rows, while Brownian independent increments gives independence within each row. Apply \texttt{i\allowbreak{}Indep\allowbreak{}Fun\_\allowbreak{}uncurry\_\allowbreak{}nat\_\allowbreak{}of\_\allowbreak{}aemeasurable} to obtain joint independence of all \((i,n)\) innovations.

\subsection{\texorpdfstring{normalized\allowbreak{}Brownian\allowbreak{}Increments independent of paths}{normalizedBrownianIncrements independent of paths}}


Restrict each independent full coordinate path to the grid \(nh\), preserving independence under composition. The preceding theorem then yields joint coordinate-time independence of the normalized increments. This avoids requiring a separate measurable map into an uncountable path space.

\subsection{\texorpdfstring{The entire innovation law}{The entire innovation law}}


Every coordinate-time innovation has law \(\mathcal N(0,1)\), and the family is jointly independent. The countable product-law theorem therefore identifies the whole innovation map with \(\bigotimes_{(i,n)}\mathcal N(0,1)\). Countability supplies measurability of the path-valued map from coordinatewise almost-everywhere measurability.

\section{\texorpdfstring{Pull Brownian noise onto every coupled initialization}{Pull Brownian noise onto every coupled initialization}}
\label{proof-module-UnadjustedLangevin-ProductBrownian-lean}
{\small\color{gray}Source module: \texttt{Unadjusted\allowbreak{}Langevin/\allowbreak{}Product\allowbreak{}Brownian.\allowbreak{}lean}.\par}


Brownian-coordinate laws, innovation independence and their L\ensuremath{{}^{2}} realizations survive passage to the product \(\gamma\otimes Q\). This builds all reference fields except the explicitly displayed stationarity premise, discharged by the later canonical assembly.

\subsection{\texorpdfstring{pre\allowbreak{}Brownian comp measure\allowbreak{}Preserving}{preBrownian comp measurePreserving}}


A measure-preserving map leaves the law of every measurable random variable unchanged after pullback. Apply this to every finite-dimensional Brownian vector. Their Gaussian laws and covariance data are therefore preserved, giving a pre-Brownian process on the new sample space.

\subsection{\texorpdfstring{brownian\allowbreak{}Coordinates\allowbreak{}Lp pullback}{brownianCoordinatesLp pullback}}


Both sides represent the vector \((B_i(t,f(\omega)))_i\) almost everywhere. Pull back the source representative equality using measure preservation and combine it with the target \texttt{to\allowbreak{}Lp} equality. \texttt{Lp.\allowbreak{}ext} identifies the two quotient classes.

\subsection{\texorpdfstring{brownian\allowbreak{}Euler\allowbreak{}Noise pullback}{brownianEulerNoise pullback}}


The \(L^2\) pullback is linear, hence commutes with subtraction and scalar multiplication. Apply the preceding Brownian-vector pullback identity to the two endpoints in the normalized increment. This identifies the pulled-back Euler noise with the noise constructed after pullback.

\subsection{\texorpdfstring{source\allowbreak{}Euler\allowbreak{}Noise}{sourceEulerNoise}}


Take the actual selected representative of the previously constructed \(L^2\) Brownian Euler noise on the source space. This is the random function supplied to the generic product algorithm before choosing any initial coupling.

\subsection{\texorpdfstring{source\allowbreak{}Euler\allowbreak{}Noise mem\allowbreak{}Lp}{sourceEulerNoise memLp}}


Every representative of a given \(L^2\) class is square integrable by \texttt{Lp.\allowbreak{}mem\allowbreak{}Lp}. Apply that fact to the Brownian Euler noise class. This avoids assuming an extra noise-moment condition in the product construction.

\subsection{\texorpdfstring{source\allowbreak{}Euler\allowbreak{}Noise coe}{sourceEulerNoise coe}}


Unfold the nested class-valued noise definition and apply representative equalities for scalar multiplication, subtraction, and the Brownian vector classes. Almost everywhere the result is \((\sqrt{2h})^{-1}\sqrt2(B_{nh+h}-B_{nh})\), with real-time clipping as in the original definition.

\subsection{\texorpdfstring{source\allowbreak{}Euler\allowbreak{}Noise coordinate coe}{sourceEulerNoise coordinate coe}}


Evaluate the vector representative formula in coordinate \(i\). For \(h>0\), use \(\sqrt{2h}=\sqrt2\sqrt h\) and cancel the nonzero \(\sqrt2\), obtaining \((B_i((n+1)h)-B_i(nh))/\sqrt h\).

\subsection{\texorpdfstring{source\allowbreak{}Euler\allowbreak{}Noise coordinate coe nnreal}{sourceEulerNoise coordinate coe nnreal}}


For nonnegative mesh times, \texttt{Real.\allowbreak{}to\allowbreak{}NNReal} leaves their real values unchanged. Rewrite both clipped grid endpoints into the corresponding nonnegative-real products. The coordinate formula becomes the ordinary Brownian grid increment divided by \(\sqrt h\).

\subsection{\texorpdfstring{source\allowbreak{}Euler\allowbreak{}Noise normalized coordinate}{sourceEulerNoise normalized coordinate}}


The nonnegative-grid formula is exactly the definition of \texttt{normalized\allowbreak{}Brownian\allowbreak{}Increment}. After converting the positive real step to its nonnegative-real form, the two random functions agree almost everywhere. This gives a direct bridge to the established innovation-law theorems.

\subsection{\texorpdfstring{source\allowbreak{}Euler\allowbreak{}Noise coordinate\allowbreak{}Law}{sourceEulerNoise coordinateLaw}}


A normalized Brownian increment has law \(\mathcal N(0,1)\). Transfer this law across \texttt{source\allowbreak{}Euler\allowbreak{}Noise\_\allowbreak{}normalized\_\allowbreak{}coordinate} using almost-everywhere invariance of laws. Hence the actual source noise representative has the required Gaussian coordinate law.

\subsection{\texorpdfstring{source\allowbreak{}Euler\allowbreak{}Noise independent}{sourceEulerNoise independent}}


The Brownian grid-path theorem gives independence of all normalized coordinates indexed by \((i,n)\). Reindex by the injective coordinate swap to \((n,i)\), then transfer across the almost-everywhere coordinate identities. This matches the index order expected by the vector innovation theorem.

\subsection{\texorpdfstring{product\allowbreak{}Pre\allowbreak{}Brownian}{productPreBrownian}}


Projection from an initial-coupling product extension to the noise source is measure preserving. Pull each scalar Brownian coordinate back along it using \texttt{pre\allowbreak{}Brownian\_\allowbreak{}comp\_\allowbreak{}measure\allowbreak{}Preserving}. The actual coordinate processes on the extension retain their Brownian finite-dimensional laws.

\subsection{\texorpdfstring{product\allowbreak{}Noise eq brownian\allowbreak{}Euler\allowbreak{}Noise}{productNoise eq brownianEulerNoise}}


The lifted source noise and the extension's constructed Brownian noise are both the measure-preserving pullback of the source \(L^2\) Brownian Euler noise. Use representative extensionality to identify the former with this pullback, then apply \texttt{brownian\allowbreak{}Euler\allowbreak{}Noise\_\allowbreak{}pullback} to identify the latter.

\subsection{\texorpdfstring{product\allowbreak{}Euler\allowbreak{}Noise coordinate\allowbreak{}Law}{productEulerNoise coordinateLaw}}


Pull the source standard-Gaussian coordinate law through the extension's noise projection. The lifted product-noise representative agrees almost everywhere with the raw pulled-back source noise. Transfer the law across that equality to obtain the coordinate law for the exact classes used by the numerical chain.

\subsection{\texorpdfstring{product\allowbreak{}Euler\allowbreak{}Noise independent}{productEulerNoise independent}}


Identify the source coordinate-time path law as the product of standard Gaussians. Pull that law to the extension; its coordinate family is independent by the product-law independence criterion. Finally replace raw source pullbacks by the almost-everywhere equal product-noise representatives.

\subsection{\texorpdfstring{product\allowbreak{}Euler\allowbreak{}Noise independent\allowbreak{}Initial}{productEulerNoise independentInitial}}


The raw first initial coordinate depends only on the coupling factor, while the whole source noise sequence depends only on the noise factor. Product independence therefore applies. Countability of the sequence lets representative equalities hold simultaneously at all indices; transport independence to the selected \(L^2\) initialization and noise representatives.

\subsection{\texorpdfstring{product\allowbreak{}Euler\allowbreak{}Noise independent\allowbreak{}Reference}{productEulerNoise independentReference}}


Use the same product independence argument with the second initial coordinate. Both initialization and whole noise-path representative equalities are transferred almost everywhere, proving the actual reference class is independent of all fresh Euler noise.

\subsection{\texorpdfstring{brownian\allowbreak{}Product\allowbreak{}Probability}{brownianProductProbability}}


The coupling is a probability by its marginal identity and the Brownian source measure is a probability. Their product extension is normalized. This local instance supplies the probability-space hypotheses of the Brownian pullback and \(L^2\) reference constructors.

\subsection{\texorpdfstring{product\allowbreak{}Gibbs\allowbreak{}Mem\allowbreak{}Lp}{productGibbsMemLp}}


Name the already proved finite-second-moment theorem for the strongly convex Gibbs law. The abbreviation supplies exactly the \(L^2(\pi_U)\) identity moment needed to construct the reference initialization on every coupling extension.

\subsection{\texorpdfstring{brownian\allowbreak{}Product\allowbreak{}L2\allowbreak{}Reference}{brownianProductL2Reference}}


Construct the Brownian reference on the extension using the supplied stationarity of its selected segments. Replace its Brownian-defined Euler noise by the common lifted source noise using \texttt{product\allowbreak{}Noise\_\allowbreak{}eq\_\allowbreak{}brownian\allowbreak{}Euler\allowbreak{}Noise}. All remaining continuity, integral, score, and noise fields come from the established constructors.

\subsection{\texorpdfstring{product\allowbreak{}Reference of constructed brownian}{productReference of constructed brownian}}


For each finite initial coupling, use \texttt{brownian\allowbreak{}Product\allowbreak{}L2\allowbreak{}Reference}. Its initial reference equals the second-coordinate \(L^2\) class by the segment-initial identity. These constructions fill the \texttt{Product\allowbreak{}Reference} fields. Only the explicitly quantified segment-stationarity premise remains an input at this intermediate stage.

\subsection{\texorpdfstring{dalalyan of constructed product brownian}{dalalyan of constructed product brownian}}


Derive \(m\le L\) from strong convexity and smoothness, construct the product reference from the stated stationary segment laws, and apply the small-step product-reference theorem. This intermediate endpoint has the genuine initial \(W_2\) term; the later independent Brownian construction discharges its stationarity premise.

\section{\texorpdfstring{Identify recursive sampling laws with iterated Markov transitions}{Identify recursive sampling laws with iterated Markov transitions}}
\label{proof-module-UnadjustedLangevin-MarkovLaw-lean}
{\small\color{gray}Source module: \texttt{Unadjusted\allowbreak{}Langevin/\allowbreak{}Markov\allowbreak{}Law.\allowbreak{}lean}.\par}


The path-space law of an iid-noise recursion agrees with the repeated one-step transition \(\mu\mapsto F_\#(\mu\otimes\eta)\). The proof explicitly establishes independence of the current state and the fresh innovation.

\subsection{\texorpdfstring{finite\allowbreak{}Noise\allowbreak{}Chain}{finiteNoiseChain}}


Evaluate the same recursion using only a finite prefix \((w_0,\ldots,w_{n-1})\). At the successor, recurse on the first \(n\) coordinates and use the final coordinate as the next innovation. This makes finite dependence explicit for the fresh-noise independence proof.

\subsection{\texorpdfstring{finite\allowbreak{}Noise\allowbreak{}Chain eq}{finiteNoiseChain eq}}


Induct on prefix length. Restricting an infinite path to its first \(n\) entries gives the same previous state, and the final finite coordinate is \(w_n\). The successor updates thus agree, identifying finite-prefix and infinite-path implementations pointwise.

\subsection{\texorpdfstring{finite\allowbreak{}Noise\allowbreak{}Chain measurable}{finiteNoiseChain measurable}}


The zero map is the initial projection. At the successor, restriction to the earlier finite coordinates and evaluation of the last coordinate are measurable; combine them with the measurable update and the induction hypothesis. This proves the state depends measurably on a finite prefix.

\subsection{\texorpdfstring{iid prefix independent fresh}{iid prefix independent fresh}}


The index sets \(\{0,\ldots,n-1\}\) and \(\{n\}\) are disjoint. Independence of the coordinate family on the iid product source therefore gives independence of the whole prefix vector and the next coordinate. This is stronger than pairwise independence of individual past entries.

\subsection{\texorpdfstring{markov\allowbreak{}Step}{markovStep}}


Define \(M_\eta(\mu)=F_*(\mu\otimes\eta)\). Its product measure means the innovation is fresh and independent of the current state. The subsequent theorem derives that independence for the canonical recursion, rather than assuming it at each step.

\subsection{\texorpdfstring{markov\allowbreak{}Laws}{markovLaws}}


Set \(\mu_0=\mu\), \(\mu_{n+1}=M_\eta(\mu_n)\). These are iterated independent-noise transition laws. The recursion alone does not yet assert that a previously defined sample-path algorithm realizes them.

\subsection{\texorpdfstring{chain\allowbreak{}Law joint fresh}{chainLaw joint fresh}}


The current state is a measurable function of initialization and the first \(n\) noises. Product-source independence and prefix/\allowbreak{}fresh independence make it independent of \(w_n\). Its marginal is \texttt{chain\allowbreak{}Law}, and the innovation marginal is \(\eta\); their joint law is therefore the product of those actual marginals.

\subsection{\texorpdfstring{chain\allowbreak{}Law succ iid}{chainLaw succ iid}}


The pointwise successor is \(F(X_n,w_n)\). Push the proved product joint law of \((X_n,w_n)\) through \(F\); this yields \(\mu_{n+1}=M_\eta(\mu_n)\) for the canonical recursion laws. No later-time state/\allowbreak{}noise independence is assumed.

\subsection{\texorpdfstring{chain\allowbreak{}Law eq markov\allowbreak{}Laws}{chainLaw eq markovLaws}}


The canonical path law and the iterated Markov law have the same initial value and the same successor equation. Induction on \(n\) identifies them at every index. This closes the connection between sample-path and transition-law definitions.

\subsection{\texorpdfstring{markov\allowbreak{}Step is\allowbreak{}Probability}{markovStep isProbability}}


The independent product of two probability measures has mass one. A measurable pushforward preserves that total mass, so \texttt{markov\allowbreak{}Step} is again a probability measure.

\subsection{\texorpdfstring{markov\allowbreak{}Laws is\allowbreak{}Probability}{markovLaws isProbability}}


Induct from the initial probability and apply \texttt{markov\allowbreak{}Step\_\allowbreak{}is\allowbreak{}Probability} at each successor. This proves normalization of every transition iterate and provides the probability instances used in later law comparisons.

\subsection{\texorpdfstring{noise\allowbreak{}Driven\allowbreak{}Chain law eq markov\allowbreak{}Laws}{noiseDrivenChain law eq markovLaws}}


First identify the independent iid realization with the canonical product-path output law. Then apply \texttt{chain\allowbreak{}Law\_\allowbreak{}eq\_\allowbreak{}markov\allowbreak{}Laws}. Every realization satisfying the original initialization/\allowbreak{}noise assumptions consequently has the same proved Markov iterate law.

\section{\texorpdfstring{A finite O(h\ensuremath{{}^{3/2}}) rejected-\allowbreak{}displacement bound}{A finite O(h³ᐟ²) rejected-displacement bound}}
\label{proof-module-UnadjustedLangevin-MALAError-lean}
{\small\color{gray}Source module: \texttt{Unadjusted\allowbreak{}Langevin/\allowbreak{}MALAError.\allowbreak{}lean}.\par}


The Metropolis-adjusted Langevin proposal provides an exactly invariant approximation to the diffusion. This module proves that its mean rejected displacement under Gibbs initialization is at most \(Ch\sqrt h\), with a finite constant derived from the potential assumptions.

\subsection{\texorpdfstring{mala\allowbreak{}Proposal}{malaProposal}}


Define \(T_h(x,\xi)=x-h\nabla U(x)+\sqrt{2h}\xi\). This is the same Euler proposal used by ULA; MALA applies an accept/\allowbreak{}reject rule to this proposed state.

\subsection{\texorpdfstring{mala\allowbreak{}Log\allowbreak{}Ratio}{malaLogRatio}}


The log acceptance ratio is \(R_h(x,y)=U(x)-U(y)+(\|y-x+h\nabla U(x)\|^2-\|x-y+h\nabla U(y)\|^2)/(4h)\). The common Gaussian normalizers cancel in the reverse/\allowbreak{}forward proposal density ratio.

\subsection{\texorpdfstring{mala\allowbreak{}Log\allowbreak{}Ratio expansion}{malaLogRatio expansion}}


Expand both squared norms using the real inner-product norm identity. The \(\|y-x\|^2\) terms cancel, and division by \(4h\) gives \(R_h=U(x)-U(y)+\frac12\langle\nabla U(x)+\nabla U(y),y-x\rangle-\frac h4(\|\nabla U(y)\|^2-\|\nabla U(x)\|^2)\). The premise \(h\ne0\) justifies denominator cancellation.

\subsection{\texorpdfstring{gradient norm sq difference bound}{gradient norm sq difference bound}}


Factor \(|\|g_y\|^2-\|g_x\|^2|=|\|g_y\|-\|g_x\||(\|g_y\|+\|g_x\|)\). The reverse triangle inequality bounds the first factor by \(\|g_y-g_x\|\), then Lipschitzness bounds it by \(L\|y-x\|\). The other factor is nonnegative, so multiplication preserves the estimate.

\subsection{\texorpdfstring{potential trapezoid bound}{potential trapezoid bound}}


Use the smooth upper and lower first-order models for the potential increment at the two endpoints. Compare the increment with the average endpoint directional derivative, and use the Lipschitz bound on the derivative difference. The resulting trapezoid error is at most \(L\|y-x\|^2\); no second derivative of \(U\) is assumed.

\subsection{\texorpdfstring{mala\allowbreak{}Log\allowbreak{}Ratio abs bound}{malaLogRatio abs bound}}


Insert the exact log-ratio expansion and use the triangle inequality. Bound its potential trapezoid term by \(L\|y-x\|^2\) and its squared-gradient difference by \(L\|y-x\|(\|g_x\|+\|g_y\|)\). Positivity of \(h\) gives the displayed coefficient \(hL/4\) without an absolute value on the step.

\subsection{\texorpdfstring{metropolis rejection le abs}{metropolis rejection le abs}}


If \(r\ge0\), \(e^r\ge1\), so rejection is zero. If \(r<0\), \(1-e^r\le-r=|r|\) by \(e^r\ge1+r\). Both cases also give \(0\le1-\min(1,e^r)\), establishing a nonnegative rejection probability and its log-ratio bound.

\subsection{\texorpdfstring{rejected displacement le}{rejected displacement le}}


Multiply \(1-\min(1,e^r)\le|r|\) by the nonnegative displacement norm \(\|y-x\|\). This yields the rejection-weighted coupling error bound with no condition on the states.

\subsection{\texorpdfstring{mala\allowbreak{}Amplitude}{malaAmplitude}}


Use the step-independent amplitude \(A(x,\xi)=\|\nabla U(x)\|+\sqrt2\|\xi\|\). For \(0\le h\le1\), it bounds the proposal displacement after extracting \(\sqrt h\); its cubic integrability will give a uniform rejection-error constant.

\subsection{\texorpdfstring{mala\allowbreak{}Rejection\allowbreak{}Displacement}{malaRejectionDisplacement}}


Define \(D_h(x,\xi)=\|T_h(x,\xi)-x\|(1-\min(1,e^{R_h(x,T_h(x,\xi))}))\). This is the acceptance-variable-averaged distance lost on rejection, so no additional uniform random variable appears in its integral.

\subsection{\texorpdfstring{mala\allowbreak{}Amplitude nonneg}{malaAmplitude nonneg}}


Both score and noise norms are nonnegative, and \(\sqrt2\ge0\). Their sum defining \(A\) is therefore nonnegative. This sign permits later powers, products, and integral comparisons without absolute-value ambiguity.

\subsection{\texorpdfstring{mala\allowbreak{}Rejection\allowbreak{}Displacement nonneg}{malaRejectionDisplacement nonneg}}


The proposal displacement norm is nonnegative and the rejection factor is nonnegative by \texttt{metropolis\_\allowbreak{}rejection\_\allowbreak{}le\_\allowbreak{}abs}. Their product is nonnegative, permitting use of extended integrals and domination by an integrable envelope.

\subsection{\texorpdfstring{mala\allowbreak{}Proposal displacement le}{malaProposal displacement le}}


The triangle inequality gives \(\|T_h-x\|\le h\|g_x\|+\sqrt{2h}\|\xi\|\). For \(0\le h\le1\), \(h\le\sqrt h\) and \(\sqrt{2h}=\sqrt2\sqrt h\). Factor out \(\sqrt h\) to obtain \(\|T_h-x\|\le\sqrt h A(x,\xi)\).

\subsection{\texorpdfstring{mala\allowbreak{}Rejection\allowbreak{}Displacement pointwise}{malaRejectionDisplacement pointwise}}


Use rejection \(\le|R_h|\), the log-ratio estimate, and \(\|g_{T_h}\|\le\|g_x\|+L\|T_h-x\|\). Replace the displacement by \(\sqrt hA\), and absorb the remaining step factors with \(h\le1\). The result is \(D_h\le h\sqrt h\,L(2+L)A^3\), whose envelope has no step-size dependence.

\subsection{\texorpdfstring{integrable norm fourth mul exp neg mul norm sq}{integrable norm fourth mul exp neg mul norm sq}}


A polynomial of degree four can be absorbed by part of Gaussian decay: \(\|x\|^4e^{-a\|x\|^2}\le C e^{-(a/2)\|x\|^2}\) for \(a>0\). Gaussian integrability and measurable domination therefore prove the fourth weighted moment is finite.

\subsection{\texorpdfstring{integrable norm fourth mul exp neg of strong\allowbreak{}Convex}{integrable norm fourth mul exp neg of strongConvex}}


The quadratic strong-convexity lower bound gives \(e^{-U(x)}\le C e^{-(m/4)\|x\|^2}\). Multiply by \(\|x\|^4\) and invoke the fourth Gaussian weighted-moment theorem. This derives the unnormalized Gibbs fourth moment from the original potential assumptions.

\subsection{\texorpdfstring{integrable norm fourth gibbs\allowbreak{}Measure}{integrable norm fourth gibbsMeasure}}


Change measure to the normalized Gibbs density. Integrability of \(\|x\|^4\) under \(\pi_U\) reduces to the just-proved integrability of \(\|x\|^4e^{-U(x)}\) times a finite normalization factor. No fourth-moment premise is added.

\subsection{\texorpdfstring{mem\allowbreak{}Lp three id gibbs\allowbreak{}Measure}{memLp three id gibbsMeasure}}


The fourth-moment theorem gives the identity variable in \(L^4(\pi_U)\). On a probability space, \(L^4\subset L^3\), so the identity belongs to \(L^3\). This supplies the cubic moment needed for the rejection envelope.

\subsection{\texorpdfstring{mala\allowbreak{}Pair\allowbreak{}Measure}{malaPairMeasure}}


Take \(\pi_U\otimes\mathcal N(0,I)\), representing a Gibbs state and an independent standard Gaussian proposal innovation. This is the joint law under which the local rejection error is averaged.

\subsection{\texorpdfstring{mala\allowbreak{}Amplitude mem\allowbreak{}Lp three}{malaAmplitude memLp three}}


The Gibbs state is in \(L^3\); Lipschitz growth \(\|\nabla U(x)\|\le\|\nabla U(0)\|+L\|x\|\) puts its score in \(L^3\). Gaussian noise also belongs to \(L^3\). Lift both to their independent product and add their norm envelopes to obtain \(A\in L^3\).

\subsection{\texorpdfstring{mala\allowbreak{}Amplitude cube integrable}{malaAmplitude cube integrable}}


For nonnegative \(A\), its third power is \(\|A\|^3\). The \(L^3\) characterization gives integrability of this function. This is the finite ordinary integral used to define the rejection constant.

\subsection{\texorpdfstring{mala\allowbreak{}Rejection\allowbreak{}Constant}{malaRejectionConstant}}


Set \(C_{U,L}=L(2+L)\int A(x,\xi)^3d(\pi_U\otimes\gamma)\). The preceding declaration proves the integral is finite under the potential assumptions. All dependence on the positive step size has been removed from this constant.

\subsection{\texorpdfstring{mala\allowbreak{}Rejection\allowbreak{}Constant nonneg}{malaRejectionConstant nonneg}}


The cubic amplitude integrand is nonnegative and its integral is nonnegative. Since \(L\ge0\), the factor \(L(2+L)\) is nonnegative as well. Their product gives \(C_{U,L}\ge0\).

\subsection{\texorpdfstring{mala\allowbreak{}Rejection\allowbreak{}Displacement integrable}{malaRejectionDisplacement integrable}}


Continuity of the potential and gradient makes the rejection-displacement function measurable. Its absolute value equals itself by nonnegativity, and its pointwise bound is the integrable function \(h\sqrt h\,L(2+L)A^3\). Apply measurable domination to prove ordinary integrability.

\subsection{\texorpdfstring{mala\allowbreak{}Rejection\allowbreak{}Displacement integral bound}{malaRejectionDisplacement integral bound}}


Integrate the pointwise rejection envelope using the established integrability of both sides. Pull the step-independent factors out of the integral and recognize \(C_{U,L}\). This gives \(\mathbb ED_h\le C_{U,L}h\sqrt h\) for \(0<h\le1\).

\subsection{\texorpdfstring{exists mala\allowbreak{}Rejection\allowbreak{}Constant}{exists malaRejectionConstant}}


Choose the explicitly defined \texttt{mala\allowbreak{}Rejection\allowbreak{}Constant U L}. Its nonnegativity and the integrated envelope theorem verify the existential statement for every \(0<h\le1\) simultaneously. The same constant works throughout the small-step limit.

\section{\texorpdfstring{Detailed balance through accepted and rejected flows}{Detailed balance through accepted and rejected flows}}
\label{proof-module-UnadjustedLangevin-Metropolis-lean}
{\small\color{gray}Source module: \texttt{Unadjusted\allowbreak{}Langevin/\allowbreak{}Metropolis.\allowbreak{}lean}.\par}


The Metropolis correction preserves the initial marginal whenever the accepted proposal flow is invariant under swapping its two states. A positive proposal-pair density makes this symmetry explicit.

\subsection{\texorpdfstring{accepted\allowbreak{}Flow}{acceptedFlow}}


Weight the current/\allowbreak{}proposal joint measure \(\nu(dx,dy)\) by acceptance probability \(a(x,y)\). Thus \(A(dx,dy)=a(x,y)\nu(dx,dy)\), with \texttt{of\allowbreak{}Real} embedding its real nonnegative weight. The bounds \(0\le a\le1\) are supplied in the conservation theorems below.

\subsection{\texorpdfstring{rejected\allowbreak{}Flow}{rejectedFlow}}


The rejected pair measure is \(R(dx,dy)=(1-a(x,y))\nu(dx,dy)\). It retains the proposed state as an auxiliary coordinate; the eventual output maps it to the original state \(x\), preserving all rejection mass.

\subsection{\texorpdfstring{metropolis\allowbreak{}Output}{metropolisOutput}}


Define the output as \(\operatorname{snd}_*A+\operatorname{fst}_*R\). Accepted moves output \(y\), while rejected moves output \(x\). This is the complete transition measure and requires no renormalization of accepted moves.

\subsection{\texorpdfstring{accepted\allowbreak{}Flow add rejected\allowbreak{}Flow}{acceptedFlow add rejectedFlow}}


Since \(0\le a\le1\), the two nonnegative weights add to one after embedding in extended scalars. Additivity of with-density measures then gives \(A+R=\nu\). The proof checks both nonnegativity guards before combining \texttt{of\allowbreak{}Real} values.

\subsection{\texorpdfstring{marginals eq of swap invariant}{marginals eq of swap invariant}}


The identity \(\operatorname{fst}\circ\operatorname{swap}=\operatorname{snd}\) and pushforward composition give \(\operatorname{snd}_*\nu=\operatorname{fst}_*(\operatorname{swap}_*\nu)\). If the pair measure is swap invariant, this is \(\operatorname{fst}_*\nu\), proving equality of its marginals.

\subsection{\texorpdfstring{metropolis\allowbreak{}Output eq initial}{metropolisOutput eq initial}}


Detailed balance makes the two marginals of the accepted flow equal. Replace its second marginal by its first, then combine the accepted and rejected first marginals using \(A+R=\nu\). The output is therefore \(\operatorname{fst}_*\nu\), exactly the original state law.

\subsection{\texorpdfstring{map weighted comp}{map weighted comp}}


For a measurable \(f\) and weight \(w\), prove \(f_*((w\circ f)\mu)=w(f_*\mu)\). On each measurable set, the pushforward definition and nonnegative change-of-variables formula turn both sides into \(\int\mathbf1_A(f(x))w(f(x))d\mu(x)\). This is the weight-transport identity used by the actual MALA proposal map.

\subsection{\texorpdfstring{map weighted comp equiv}{map weighted comp equiv}}


Specialize the preceding pushforward/\allowbreak{}weight identity to a measurable equivalence. Its bundled measurability supplies the required hypothesis. The application below uses the coordinate swap equivalence to prove symmetry of an accepted density.

\subsection{\texorpdfstring{metropolis\allowbreak{}Acceptance}{metropolisAcceptance}}


For positive joint density \(w(x,y)\), define \(a(x,y)=\min(1,w(y,x)/w(x,y))\). The ratio is reverse flow divided by forward flow; positivity ensures both its meaning and nonnegative acceptance.

\subsection{\texorpdfstring{weight mul metropolis\allowbreak{}Acceptance}{weight mul metropolisAcceptance}}


Distribute the positive factor \(w(x,y)\) through the minimum. Cancellation of the nonzero denominator gives \(w(x,y)a(x,y)=\min(w(x,y),w(y,x))\). This algebra is the reason the accepted joint density is symmetric.

\subsection{\texorpdfstring{accepted\allowbreak{}Flow density}{acceptedFlow density}}


Two successive changes of density multiply their weights. Combine the original positive density \(w\) and acceptance \(a\), then apply the weight/\allowbreak{}minimum identity. The accepted flow becomes the base measure weighted by \(\min(w,w\circ\operatorname{swap})\), a visibly symmetric density.

\subsection{\texorpdfstring{metropolis\allowbreak{}Output density eq initial}{metropolisOutput density eq initial}}


The base measure is swap invariant, and the accepted minimum density is unchanged by swapping its arguments. The weight-transport theorem therefore proves accepted-flow detailed balance. The general accepted/\allowbreak{}rejected conservation theorem then identifies the output with the initial marginal.

\section{\texorpdfstring{Couple ULA to an invariant Metropolis chain}{Couple ULA to an invariant Metropolis chain}}
\label{proof-module-UnadjustedLangevin-MALACoupling-lean}
{\small\color{gray}Source module: \texttt{Unadjusted\allowbreak{}Langevin/\allowbreak{}MALACoupling.\allowbreak{}lean}.\par}


A common Gaussian proposal couples unadjusted ULA with MALA. Under a stable small step, their expected distance grows by at most the Gibbs-average rejected displacement at each step. This is an auxiliary approximation for stationarity.

\subsection{\texorpdfstring{mala\allowbreak{}Acceptance}{malaAcceptance}}


Define \(a_h(x,y)=\min(1,e^{R_h(x,y)})\), with the proved target/\allowbreak{}proposal log ratio. The minimum caps acceptance at one and the exponential ensures nonnegativity; these bounds justify the accepted/\allowbreak{}rejected measure split.

\subsection{\texorpdfstring{mala\allowbreak{}Proposal\allowbreak{}Acceptance}{malaProposalAcceptance}}


Compose \(a_h\) with \((x,\xi)\mapsto(x,T_h(x,\xi))\). The resulting \(a_h(x,T_h(x,\xi))\) is the acceptance weight on the original state/\allowbreak{}noise product space, prior to pushing to state/\allowbreak{}proposal pairs.

\subsection{\texorpdfstring{mala\allowbreak{}Acceptance nonneg}{malaAcceptance nonneg}}


Both entries of the minimum, one and the exponential, are nonnegative. Their minimum is nonnegative. This direct order proof is used when combining real weights through \texttt{ENNReal.\allowbreak{}of\allowbreak{}Real}.

\subsection{\texorpdfstring{mala\allowbreak{}Acceptance le one}{malaAcceptance le one}}


The minimum \(\min(1,e^{R_h})\) is at most its first entry. This proves \(a_h\le1\), so the rejection weight \(1-a_h\) is nonnegative.

\subsection{\texorpdfstring{mala\allowbreak{}Proposal measurable}{malaProposal measurable}}


The gradient is continuous by Lipschitzness. Projections, scalar multiplication, subtraction, and addition preserve continuity, so the joint proposal map \((x,\xi)\mapsto x-h\nabla U(x)+\sqrt{2h}\xi\) is measurable.

\subsection{\texorpdfstring{mala\allowbreak{}Proposal\allowbreak{}Acceptance measurable}{malaProposalAcceptance measurable}}


Differentiability makes \(U\) continuous, and Lipschitzness makes its gradient continuous. The log-ratio formula is assembled from continuous norms, inner products, and arithmetic; exponentiation and minimum preserve measurability. Composing with the measurable proposal gives the state/\allowbreak{}noise acceptance weight.

\subsection{\texorpdfstring{ula\allowbreak{}Law\allowbreak{}Step}{ulaLawStep}}


Define \(P_h\mu\) as the pushforward of \(\mu\otimes\gamma\) by the Euler proposal \(T_h\). The product measure encodes that the fresh Gaussian is independent of the current state. This transition is the actual independent-noise ULA step on laws.

\subsection{\texorpdfstring{mala\allowbreak{}One\allowbreak{}Step}{malaOneStep}}


Split \(\mu\otimes\gamma\) into acceptance- and rejection-weighted measures. Map the accepted part to \(T_h(x,\xi)\), and the rejected part to \(x\), then add them. This defines one full MALA transition with all rejection mass retained.

\subsection{\texorpdfstring{mala\allowbreak{}Coupled\allowbreak{}Step}{malaCoupledStep}}


For a pair \((x,y)\), draw one shared Gaussian \(\xi\). ULA always moves to \(T_h(x,\xi)\); MALA moves to \(T_h(y,\xi)\) with weight \(a_h(y,T_h(y,\xi))\), otherwise stays at \(y\). The two weighted pushforwards define the explicit coupling measure.

\subsection{\texorpdfstring{weighted accept reject add}{weighted accept reject add}}


For a measurable \(a\in[0,1]\), combine the two with-density measures using \(a+(1-a)=1\). Nonnegativity of both terms justifies the \texttt{of\allowbreak{}Real} addition identity, so their sum is exactly the original measure.

\subsection{\texorpdfstring{mala\allowbreak{}Coupled\allowbreak{}Step fst}{malaCoupledStep fst}}


The first-coordinate update is identical on acceptance and rejection. Combine those two weights to one before pushing forward. The remaining product marginal is the original first marginal times a fresh Gaussian, so the first output law is \(P_h(\operatorname{fst}_*\rho)\).

\subsection{\texorpdfstring{mala\allowbreak{}Coupled\allowbreak{}Step snd}{malaCoupledStep snd}}


Project both branches of the explicit coupling to the second state. They are exactly the accepted proposal and rejected current-state branches of MALA. Product marginal identities replace \(\rho\) by its second marginal, giving \(\operatorname{snd}_*\rho'=M_h(\operatorname{snd}_*\rho)\).

\subsection{\texorpdfstring{mala\allowbreak{}Coupled\allowbreak{}Step is\allowbreak{}Probability}{malaCoupledStep isProbability}}


The first marginal is a measurable pushforward of a probability product by \texttt{mala\allowbreak{}Coupled\allowbreak{}Step\_\allowbreak{}fst}. It has mass one, and a measure and its first-coordinate pushforward have the same total mass. Therefore the joint coupling also has total mass one.

\subsection{\texorpdfstring{mala\allowbreak{}Coupled\allowbreak{}Chain}{malaCoupledChain}}


Initialize on the diagonal \((X_0,Y_0)=(X,X)\) with \(X\sim\mu\), then recursively apply the explicit coupled step. Equal starting states give zero initial transport cost, while subsequent first and second marginals follow ULA and MALA respectively.

\subsection{\texorpdfstring{mala\allowbreak{}Coupled\allowbreak{}Chain is\allowbreak{}Probability}{malaCoupledChain isProbability}}


The diagonal pushforward of the initial probability is a probability. Induct on the number of transitions, using probability preservation of \texttt{mala\allowbreak{}Coupled\allowbreak{}Step} at each successor. This produces the instance needed for product noise and finite-cost arguments at every iteration.

\subsection{\texorpdfstring{mala\allowbreak{}Coupled\allowbreak{}Chain fst}{malaCoupledChain fst}}


At zero, diagonal initialization has first marginal \(\mu\). In the successor step apply \texttt{mala\allowbreak{}Coupled\allowbreak{}Step\_\allowbreak{}fst} and the induction hypothesis. The resulting recursion is precisely the iteration \(P_h^n\mu\).

\subsection{\texorpdfstring{mala\allowbreak{}Coupled\allowbreak{}Chain snd}{malaCoupledChain snd}}


At zero the second marginal is \(\pi\). The successor marginal theorem gives \(M_h\) applied to the current second marginal; substitute the induction hypothesis and the supplied one-step invariance \(M_h\pi=\pi\). Thus the second coordinate remains stationary through the coupling.

\subsection{\texorpdfstring{mala\allowbreak{}Proposal nonexpansive}{malaProposal nonexpansive}}


Shared noise cancels in \(T_h(x,\xi)-T_h(y,\xi)\). The sharp strong-convexity/\allowbreak{}smoothness contraction theorem for the deterministic gradient map bounds the remaining norm by \(\|x-y\|\) when \(0\le h\le2/(m+L)\). No Gaussian moment is used in this pointwise comparison.

\subsection{\texorpdfstring{mala\allowbreak{}Coupled\allowbreak{}Step pointwise}{malaCoupledStep pointwise}}


The accepted distance is at most \(\|x-y\|\) by nonexpansiveness. On rejection, insert \(T_h(y,\xi)\) and use the triangle inequality to bound the distance by \(\|x-y\|+\|T_h(y,\xi)-y\|\). Weight the two branches; their common old-distance coefficients add to one, and the remaining term is the rejection displacement.

\subsection{\texorpdfstring{mala\allowbreak{}L1\allowbreak{}Cost}{malaL1Cost}}


Define \(\mathcal C(\rho)=\int\|x-y\|d\rho\) as a nonnegative extended integral. Keeping it extended avoids assuming integrability before the cost recurrence proves it. Subsequent theorems establish finiteness and then convert it to an ordinary expected distance.

\subsection{\texorpdfstring{lintegral weighted map}{lintegral weighted map}}


Pull an extended nonnegative integral through a measurable pushforward, then through a with-density measure. Multiplication of the two nonnegative real factors becomes the embedded product \(a(x)b(f(x))\). This gives an exact integration formula for each branch of the coupling.

\subsection{\texorpdfstring{mala\allowbreak{}Coupled\allowbreak{}Step cost formula}{malaCoupledStep cost formula}}


Apply \texttt{lintegral\_\allowbreak{}weighted\_\allowbreak{}map} to accepted and rejected branches of the coupled measure. Their cost is the integral of the acceptance-weighted proposal distance plus rejection-weighted stay distance over \(\rho\otimes\gamma\). The formula exposes exactly the pointwise expression bounded by \texttt{mala\allowbreak{}Coupled\allowbreak{}Step\_\allowbreak{}pointwise}.

\subsection{\texorpdfstring{mala\allowbreak{}Coupled\allowbreak{}Step cost le}{malaCoupledStep cost le}}


Integrate the pointwise coupling bound. The old distance integrates to \(\mathcal C(\rho)\); the error depends only on the second state and innovation. Its second marginal is Gibbs by hypothesis, so that integral is bounded by \(C_{U,L}h\sqrt h\). Thus \(\mathcal C(\rho')\le\mathcal C(\rho)+C_{U,L}h\sqrt h\), valid initially as an extended inequality.

\subsection{\texorpdfstring{mala\allowbreak{}Coupled\allowbreak{}Chain cost le}{malaCoupledChain cost le}}


Induct from zero diagonal cost. The second marginal remains Gibbs by the supplied one-step MALA invariance, so the one-step cost bound applies at every iteration. Summing the same local error gives \(\mathcal C(\rho_n)\le nC_{U,L}h\sqrt h\).

\subsection{\texorpdfstring{mala\allowbreak{}Coupled\allowbreak{}Chain distance integrable}{malaCoupledChain distance integrable}}


The quantitative extended cost estimate bounds \(\int\|x-y\|d\rho_n\) by a finite embedded real number. The distance is measurable and nonnegative, so finite extended integral is equivalent to ordinary integrability. This discharges the moment condition rather than assuming it.

\subsection{\texorpdfstring{mala\allowbreak{}Coupled\allowbreak{}Chain distance integral le}{malaCoupledChain distance integral le}}


Convert the finite extended cost inequality to real scalars using \texttt{ENNReal.\allowbreak{}to\allowbreak{}Real}. The newly proved integrability identifies the real value of the extended integral with the ordinary expected distance. Consequently \(\mathbb E_{\rho_n}\|X-Y\|\le nC_{U,L}h\sqrt h\).

\section{\texorpdfstring{Identify Gaussian proposals with their radial density}{Identify Gaussian proposals with their radial density}}
\label{proof-module-UnadjustedLangevin-GaussianProposalDensity-lean}
{\small\color{gray}Source module: \texttt{Unadjusted\allowbreak{}Langevin/\allowbreak{}Gaussian\allowbreak{}Proposal\allowbreak{}Density.\allowbreak{}lean}.\par}


For \(h>0\), the affine standard Gaussian proposal \(c+\sqrt{2h}\,Z\) has density proportional to \(e^{-\|y-c\|^2/(4h)}\). The proof identifies probability measures by characteristic functions and does not require a preselected coordinate formula for the Gaussian normalizer.

\subsection{\texorpdfstring{radial\allowbreak{}Gaussian\allowbreak{}Integral}{radialGaussianIntegral}}


Set \(J(a)=\int_H e^{-a\|z\|^2}dz\) as an ordinary real integral. The positive-precision results below prove its integrability and positivity before using \(J(a)^{-1}\) as a normalizing constant.

\subsection{\texorpdfstring{radial\allowbreak{}Gaussian\allowbreak{}Normalizer}{radialGaussianNormalizer}}


Define \(C(a)=J(a)^{-1}\). This is an actual integral normalizer on the given finite-dimensional Hilbert space, avoiding a coordinate-dependent choice of a formula for Lebesgue density.

\subsection{\texorpdfstring{radial\allowbreak{}Gaussian\allowbreak{}Integral pos}{radialGaussianIntegral pos}}


For \(a>0\), Gaussian integrability was established in \texttt{integrable\_\allowbreak{}exp\_\allowbreak{}neg\_\allowbreak{}mul\_\allowbreak{}norm\_\allowbreak{}sq}. The integrand is strictly positive, so its ordinary integral is strictly positive by \texttt{integral\_\allowbreak{}exp\_\allowbreak{}pos}.

\subsection{\texorpdfstring{radial\allowbreak{}Gaussian\allowbreak{}Normalizer pos}{radialGaussianNormalizer pos}}


The inverse of a strictly positive real number is strictly positive. Apply this to \(J(a)\), whose positivity was proved immediately above. This also proves the normalizer is nonzero for later cancellation.

\subsection{\texorpdfstring{radial\allowbreak{}Gaussian\allowbreak{}Density}{radialGaussianDensity}}


Define \(q_{a,c}(z)=C(a)e^{-a\|z-c\|^2}\), embedded in extended nonnegative scalars. The center changes the translated argument, while the normalizer depends only on precision and dimension.

\subsection{\texorpdfstring{integrable radial\allowbreak{}Gaussian\allowbreak{}Density real}{integrable radialGaussianDensity real}}


Translation by \(c\) preserves Lebesgue measure, so \(e^{-a\|z-c\|^2}\) is integrable whenever \(a>0\). Multiply by the finite real constant \(C(a)\) to obtain integrability of the displayed real density.

\subsection{\texorpdfstring{integral radial\allowbreak{}Gaussian\allowbreak{}Density real}{integral radialGaussianDensity real}}


Translate \(z\mapsto z-c\) in the integral and pull out \(C(a)\). The result is \(C(a)J(a)=1\), justified by \(J(a)>0\). Thus normalization is independent of the center.

\subsection{\texorpdfstring{radial\allowbreak{}Gaussian\allowbreak{}Density is\allowbreak{}Probability}{radialGaussianDensity isProbability}}


The total mass of the with-density measure is the extended integral of its density. Nonnegativity and the proved real integrability identify it with \texttt{of\allowbreak{}Real} of the real integral, which is one. This constructs probability normalization, including finiteness.

\subsection{\texorpdfstring{integral complex radial\allowbreak{}Gaussian}{integral complex radialGaussian}}


For real \(a\), \(e^{-(a:\mathbb C)\|z\|^2}\) is the complex embedding of the real Gaussian exponential. Rewrite the integrand with \texttt{Complex.\allowbreak{}of\allowbreak{}Real\_\allowbreak{}exp}; the continuous real embedding commutes with integration, giving the complex embedding of \(J(a)\).

\subsection{\texorpdfstring{integral radial\allowbreak{}Gaussian character}{integral radialGaussian character}}


Apply the finite-dimensional complex Gaussian Fourier integral with imaginary linear coefficient. Combining the two exponentials and using \(i^2=-1\) gives \(\int e^{-a\|z\|^2}e^{i\langle z,t\rangle}dz=J(a)e^{-\|t\|^2/(4a)}\). The zero-frequency Gaussian formula identifies the normalization factor with the same real integral \(J(a)\).

\subsection{\texorpdfstring{char\allowbreak{}Fun radial\allowbreak{}Gaussian\allowbreak{}Density}{charFun radialGaussianDensity}}


Change from the density measure to a weighted Lebesgue integral. Translate \(z=u+c\); the character splits into \(e^{i\langle u,t\rangle}e^{i\langle c,t\rangle}\). The Fourier formula gives a factor \(J(a)\), canceled by \(C(a)\), leaving \(e^{-\|t\|^2/(4a)}e^{i\langle c,t\rangle}\).

\subsection{\texorpdfstring{gaussian\allowbreak{}Proposal\allowbreak{}Normalizer}{gaussianProposalNormalizer}}


Specialize the precision to \(a=1/(4h)\). The resulting constant normalizes a Gaussian with covariance \(2hI\). Its independence of the current state is essential when the same factor cancels in the forward and reverse Metropolis ratio.

\subsection{\texorpdfstring{gaussian\allowbreak{}Proposal\allowbreak{}Density}{gaussianProposalDensity}}


Define \(q_h(c,y)=C(1/(4h))e^{-\|y-c\|^2/(4h)}\). For \(h>0\), this is the density of a covariance-\(2hI\) Gaussian centered at \(c\). The density identification is proved later rather than built into the definition.

\subsection{\texorpdfstring{gaussian\allowbreak{}Proposal\allowbreak{}Normalizer pos}{gaussianProposalNormalizer pos}}


For \(h>0\), the precision \(1/(4h)\) is positive. Apply positivity of the radial Gaussian normalizer at that precision. This nonzero constant is used in detailed-balance cancellation.

\subsection{\texorpdfstring{radial\allowbreak{}Gaussian\allowbreak{}Density eq proposal\allowbreak{}Density}{radialGaussianDensity eq proposalDensity}}


Unfold the two density definitions and the proposal normalizer. Their exponents agree by \(-(1/(4h))\|y-c\|^2=-\|y-c\|^2/(4h)\). Function extensionality gives equality of the density functions, with no probabilistic argument needed.

\subsection{\texorpdfstring{gaussian proposal density}{gaussian proposal density}}


Compute characteristic functions on both sides. The affine image \(c+\sqrt{2h}Z\) has character \(e^{-h\|t\|^2}e^{i\langle c,t\rangle}\); the radial-density formula with \(a=1/(4h)\) gives the same expression. Both measures are probabilities, so \texttt{Measure.\allowbreak{}ext\_\allowbreak{}of\_\allowbreak{}char\allowbreak{}Fun} identifies them. The square-root identity uses \(h>0\).

\subsection{\texorpdfstring{gaussian proposal radial density}{gaussian proposal radial density}}


Rewrite the proposal density using \texttt{radial\allowbreak{}Gaussian\allowbreak{}Density\_\allowbreak{}eq\_\allowbreak{}proposal\allowbreak{}Density} in the established affine-Gaussian law equality. This yields the equivalent precision-parameter form used in the detailed-balance proof.

\subsection{\texorpdfstring{ula proposal radial density}{ula proposal radial density}}


Set the affine Gaussian center to \(c=x-h\nabla U(x)\). The fixed-state density theorem then gives \(x-h\nabla U(x)+\sqrt{2h}Z\) its exact radial density. No regularity of \(U\) is required because \(x\) is fixed throughout this statement.

\subsection{\texorpdfstring{ula proposal density}{ula proposal density}}


Expand the radial density at center \(x-h\nabla U(x)\). The exponent becomes \(-\|y-x+h\nabla U(x)\|^2/(4h)\), with the same positive normalizer. This is the explicit ULA transition density used in the Metropolis acceptance ratio.

\section{\texorpdfstring{The joint initial-\allowbreak{}state and proposal density}{The joint initial-state and proposal density}}
\label{proof-module-UnadjustedLangevin-GaussianProposalPairDensity-lean}
{\small\color{gray}Source module: \texttt{Unadjusted\allowbreak{}Langevin/\allowbreak{}Gaussian\allowbreak{}Proposal\allowbreak{}Pair\allowbreak{}Density.\allowbreak{}lean}.\par}


For a measurable center \(c(x)\), independent \(X\sim\mu\) and \(Z\sim\mathcal N(0,I)\) produce the joint law \(\mu(dx)q_h(c(x),y)dy\).

\subsection{\texorpdfstring{Disintegrate the Gaussian change of variables}{Disintegrate the Gaussian change of variables}}


For \(X\sim\mu\) and independent \(Z\sim\mathcal N(0,I)\), fix \(x\) and apply the proved proposal-density identity to \(c(x)+\sqrt{2h}Z\). It gives conditional integration against \(q_h(c(x),y)dy\).



To obtain equality of joint measures, evaluate a measurable set \(A\subset H\times H\). Tonelli writes both sides as integrals over \(x\) of their \(y\)-sections. The fixed-center density theorem identifies each section, hence the measures agree. Measurability of \(c\) supplies measurability of the pair update and joint density; no probability or moment condition on \(\mu\) is added.

\section{\texorpdfstring{Exact Gibbs invariance of the Metropolis-\allowbreak{}adjusted step}{Exact Gibbs invariance of the Metropolis-adjusted step}}
\label{proof-module-UnadjustedLangevin-MALAInvariance-lean}
{\small\color{gray}Source module: \texttt{Unadjusted\allowbreak{}Langevin/\allowbreak{}MALAInvariance.\allowbreak{}lean}.\par}


The Gaussian proposal density identifies the MALA pair flow. Its reversed-to-forward density ratio is precisely the exponential of the implemented log ratio, so the Metropolis balance theorem proves exact Gibbs invariance.

\subsection{\texorpdfstring{mala\allowbreak{}Proposal\allowbreak{}Joint}{malaProposalJoint}}


Push \(\mu\otimes\gamma\) through \((x,\xi)\mapsto(x,T_h(x,\xi))\). Its first coordinate retains the current state, and its second records the actual Gaussian proposal. This joint measure is the input to the accepted-flow detailed-balance theorem.

\subsection{\texorpdfstring{mala\allowbreak{}Joint\allowbreak{}Weight}{malaJointWeight}}


Set \(w_h(x,y)=\exp(-U(x)-\|y-x+h\nabla U(x)\|^2/(4h))\). The Gibbs and Gaussian normalization constants are kept in a separate symmetric base measure, so \(w_h\) is strictly positive and its reverse/\allowbreak{}forward ratio is easy to compute.

\subsection{\texorpdfstring{mala\allowbreak{}Joint\allowbreak{}Weight ratio}{malaJointWeight ratio}}


Divide the two positive exponentials and use \(e^a/e^b=e^{a-b}\). Expanding the definition of \(w_h\) shows the exponent difference is exactly \texttt{mala\allowbreak{}Log\allowbreak{}Ratio}. Thus \(w_h(y,x)/w_h(x,y)=e^{R_h(x,y)}\).

\subsection{\texorpdfstring{mala\allowbreak{}One\allowbreak{}Step eq metropolis\allowbreak{}Output}{malaOneStep eq metropolisOutput}}


Transport acceptance and rejection weights through the measurable pair map using \texttt{map\_\allowbreak{}weighted\_\allowbreak{}comp}. Projecting the accepted pair to its second coordinate gives the proposal; projecting rejected pairs to the first gives the unchanged state. These are precisely the two branches defining \texttt{mala\allowbreak{}One\allowbreak{}Step}.

\subsection{\texorpdfstring{mala\allowbreak{}Proposal\allowbreak{}Joint gibbs density}{malaProposalJoint gibbs density}}


Apply the Gaussian joint-density theorem with center \(x-h\nabla U(x)\), then multiply its density by the Gibbs density. Combine the exponential factors into \(w_h(x,y)\). Pull the two state-independent normalizers into the base measure \(Z_U^{-1}C_h(dx\,dy)\), which is invariant under coordinate swap.

\subsection{\texorpdfstring{mala\allowbreak{}One\allowbreak{}Step gibbs invariant}{malaOneStep gibbs invariant}}


The joint proposal measure has a positive density \(w_h\) against a symmetric base. Its acceptance rule is \(\min(1,w_h(y,x)/w_h(x,y))\) by \texttt{mala\allowbreak{}Joint\allowbreak{}Weight\_\allowbreak{}ratio}. Hence its accepted density is the symmetric minimum \(\min(w_h(x,y),w_h(y,x))\).



Apply \texttt{metropolis\allowbreak{}Output\_\allowbreak{}density\_\allowbreak{}eq\_\allowbreak{}initial} and the exact step/\allowbreak{}output identity. The initial marginal of the proposal pair is Gibbs, since projecting \(\pi_U\otimes\gamma\) to its state returns \(\pi_U\). Thus \(M_h\pi_U=\pi_U\) is proved with no stationarity input.

\section{\texorpdfstring{A vanishing coupling distance determines the limiting law}{A vanishing coupling distance determines the limiting law}}
\label{proof-module-UnadjustedLangevin-CouplingLimit-lean}
{\small\color{gray}Source module: \texttt{Unadjusted\allowbreak{}Langevin/\allowbreak{}Coupling\allowbreak{}Limit.\allowbreak{}lean}.\par}


If \(X_n\to X\) in \(L^2\) and each \(\operatorname{Law}(X_n)\) couples to a fixed probability \(\pi\) with mean distance tending to zero, then \(\operatorname{Law}(X)=\pi\).

\subsection{\texorpdfstring{characteristic\allowbreak{}Test norm sub le}{characteristicTest norm sub le}}


The unit-circle map is \(1\)-Lipschitz, so the difference of the two characters is at most \(|\langle x-y,t\rangle|\). Cauchy–Schwarz then bounds this by \(\|t\|\|x-y\|\). Thus characteristic tests have an explicit Lipschitz constant at each fixed frequency.

\subsection{\texorpdfstring{char\allowbreak{}Fun marginals norm sub le}{charFun marginals norm sub le}}


Pull both marginal characteristic integrals back to the coupling space and subtract. Each character is bounded by one, hence integrable under the probability coupling. The integral norm inequality and the pointwise Lipschitz estimate give \(|\widehat\mu(t)-\widehat\pi(t)|\le\|t\|\mathbb E\|X-Y\|\), with the distance's integrability explicitly supplied.

\subsection{\texorpdfstring{continuous char\allowbreak{}Fun map Lp hilbert}{continuous charFun map Lp hilbert}}


Apply the continuous linear functional \(x\mapsto\langle t,x\rangle\) to an \(H\)-valued \(L^2\) variable. Its lifted map is continuous on \(L^2\). The Hilbert-valued characteristic function at \(t\) is the scalar projected characteristic function at frequency one, so the scalar continuity theorem proves the claim.

\subsection{\texorpdfstring{measure map Lp eq of coupling cost tendsto}{measure map Lp eq of coupling cost tendsto}}


For each \(t\), the coupling inequality forces the approximating characteristic functions toward \(\widehat\pi(t)\), since the expected distance tends to zero. Independently, \(L^2\) convergence forces them toward \(\widehat{P_X}(t)\). Uniqueness of limits along a nontrivial filter identifies these values, and \texttt{Measure.\allowbreak{}ext\_\allowbreak{}of\_\allowbreak{}char\allowbreak{}Fun} yields \(P_X=\pi\).

\section{\texorpdfstring{From scalar coordinate laws to independent Gaussian vectors}{From scalar coordinate laws to independent Gaussian vectors}}
\label{proof-module-UnadjustedLangevin-GaussianInnovationVectors-lean}
{\small\color{gray}Source module: \texttt{Unadjusted\allowbreak{}Langevin/\allowbreak{}Gaussian\allowbreak{}Innovation\allowbreak{}Vectors.\allowbreak{}lean}.\par}


For finitely many spatial coordinates, independent standard normal scalar innovations assemble into independent standard Gaussian vectors. Scalar normal laws without joint independence would not suffice for this conclusion.

\subsection{\texorpdfstring{gaussian\allowbreak{}Innovations path\allowbreak{}Law}{gaussianInnovations pathLaw}}


Regroup the jointly independent scalar coordinates indexed by \((n,i)\) into vector blocks \(W_n\). A finite block with independent standard Gaussian coordinates has the standard Euclidean Gaussian law. Product-measure identification and reindexing then show that the full vector path has law \(\bigotimes_n\mathcal N(0,I)\).

\subsection{\texorpdfstring{gaussian\allowbreak{}Innovations vector\allowbreak{}Law}{gaussianInnovations vectorLaw}}


Project the established whole-path product law to its \(n\)-th vector coordinate. Product marginals give \(P\circ W_n^{-1}=\mathcal N(0,I)\). This derives the multivariate law from the scalar coordinate assumptions rather than adding it as an algorithm hypothesis.

\subsection{\texorpdfstring{gaussian\allowbreak{}Innovations vector\allowbreak{}Independent}{gaussianInnovations vectorIndependent}}


The coordinate maps of the vector product law are independent. Pull that independence back through the actual innovation-path map using \texttt{gaussian\allowbreak{}Innovations\_\allowbreak{}path\allowbreak{}Law}. Thus the vector-valued \(W_n\) are independent over time, not just uncorrelated.

\section{\texorpdfstring{Prove Gibbs stationarity by refining invariant Metropolis chains}{Prove Gibbs stationarity by refining invariant Metropolis chains}}
\label{proof-module-UnadjustedLangevin-GibbsStationarity-lean}
{\small\color{gray}Source module: \texttt{Unadjusted\allowbreak{}Langevin/\allowbreak{}Gibbs\allowbreak{}Stationarity.\allowbreak{}lean}.\par}


A Brownian Langevin integral solution started from Gibbs stays Gibbs. The proof uses convergent Euler approximations and exactly invariant MALA chains, with no weak-PDE uniqueness assumption.

\subsection{\texorpdfstring{hilbert lp\allowbreak{}Chain coe}{hilbert lpChain coe}}


Induct over the chain recursion. The representative of the lifted gradient equals the pointwise gradient almost everywhere, and the representative formulas for addition and scalar multiplication give the same Euler successor. Intersect the finite set of null-set complements needed at each index to identify the \(L^2\) chain with the ordinary recursion.

\subsection{\texorpdfstring{hilbert lp\allowbreak{}Chain law eq ula\allowbreak{}Law\allowbreak{}Step}{hilbert lpChain law eq ulaLawStep}}


Use the almost-everywhere pointwise-chain identification, the independent-iid realization law theorem, and the canonical path/\allowbreak{}Markov-law identity. With innovation law \(\mathcal N(0,I)\), the generic independent-noise step is exactly \texttt{ula\allowbreak{}Law\allowbreak{}Step}. Thus the \(L^2\) implementation has law \(P_h^n(P_{X_0})\).

\subsection{\texorpdfstring{additive\allowbreak{}Euler\allowbreak{}Chain eq hilbert lp\allowbreak{}Chain}{additiveEulerChain eq hilbert lpChain}}


Induct on the mesh index. Both chains start at the same \(L^2\) variable. Their successor equations agree because the supplied additive increment equals \(\sqrt{2h}W_k\), making the noise terms identical as classes.

\subsection{\texorpdfstring{brownian\allowbreak{}Langevin\allowbreak{}Noise grid increment}{brownianLangevinNoise grid increment}}


Apply the Brownian input cocycle at \(a=kh\), \(t=h\). The elementary identity \((k+1)h=kh+h\) matches its endpoints to the grid notation. This supplies the exact additive increment used by the normalized Euler noise.

\subsection{\texorpdfstring{brownian additive\allowbreak{}Euler\allowbreak{}Chain eq lp\allowbreak{}Chain}{brownian additiveEulerChain eq lpChain}}


The Brownian grid-increment identity and \texttt{brownian\allowbreak{}Euler\allowbreak{}Noise\_\allowbreak{}endpoint} verify the increment premise of the generic chain equality. Therefore the deterministic matching-noise Euler approximation is exactly the already defined Brownian-normalized \(L^2\) Euler chain.

\subsection{\texorpdfstring{brownian\allowbreak{}Euler\allowbreak{}Noise vector\allowbreak{}Gaussian}{brownianEulerNoise vectorGaussian}}


Each normalized coordinate is standard Gaussian. The supplied joint scalar-coordinate independence lets \texttt{gaussian\allowbreak{}Innovations\_\allowbreak{}vector\allowbreak{}Law} and \texttt{gaussian\allowbreak{}Innovations\_\allowbreak{}vector\allowbreak{}Independent} assemble those coordinates into iid standard multivariate Gaussian vectors. These are the exact law hypotheses needed by the Markov-recursion theorem.

\subsection{\texorpdfstring{Pass the invariant Gibbs marginal to the continuous limit}{Pass the invariant Gibbs marginal to the continuous limit}}


Fix \(T>0\) and choose meshes \(h_n=T/N_n\to0\) small enough for the nonexpansive MALA coupling. Matching-noise Euler approximations converge in \(L^2\) to \(Y_T\). Their laws are exactly \(P_{h_n}^{N_n}\pi_U\), by the proved iid Brownian innovation laws and initialization independence.



Couple each Euler law to \(N_n\) MALA steps started from Gibbs. Exact MALA invariance makes the second marginal \(\pi_U\). The rejection bound gives expected coupling distance at most \(N_nC h_n\sqrt{h_n}=TC\sqrt{h_n}\to0\); integrability of this distance is proved from the finite extended-cost estimate.



Apply \texttt{measure\_\allowbreak{}map\_\allowbreak{}Lp\_\allowbreak{}eq\_\allowbreak{}of\_\allowbreak{}coupling\_\allowbreak{}cost\_\allowbreak{}tendsto}: characteristic functions of the Euler endpoints converge to those of \(Y_T\) by \(L^2\) convergence, and to those of \(\pi_U\) by the vanishing coupling cost. Equality of characteristic functions proves \(P_{Y_T}=\pi_U\). The zero-horizon case follows directly from the integral equation.

\section{\texorpdfstring{Stability and uniqueness for additive-\allowbreak{}noise solutions}{Stability and uniqueness for additive-noise solutions}}
\label{proof-module-UnadjustedLangevin-LangevinUniqueness-lean}
{\small\color{gray}Source module: \texttt{Unadjusted\allowbreak{}Langevin/\allowbreak{}Langevin\allowbreak{}Uniqueness.\allowbreak{}lean}.\par}


Two solutions driven by the same continuous additive noise satisfy \(\|Y_t-Z_t\|\le e^{Lt}\|x-y\|\) on the prescribed interval. Equality of initial values gives uniqueness. Only global Lipschitzness is used here; the sharper contractive estimate uses strong convexity elsewhere.

\subsection{\texorpdfstring{langevin integral stability}{langevin integral stability}}


Subtract the two integral equations so that their common additive input cancels. The integral norm inequality and Lipschitz drift give \(u(t)\le\|x-y\|+L\int_0^t u(s)ds\), where \(u(t)=\|Y(t)-Z(t)\|\) is continuous. Grönwall's inequality yields \(u(t)\le\|x-y\|e^{Lt}\). All time integrals are justified by continuity on the compact horizon.

\subsection{\texorpdfstring{langevin integral unique}{langevin integral unique}}


Set the two initial states equal in the stability estimate. The right side becomes \(0e^{Lt}=0\), forcing \(\|Y(t)-Z(t)\|=0\) for each \(t\in[0,h]\). Norm-zero equivalence gives equality of the vectors, packaged in Lean as \texttt{Eq\allowbreak{}On} on the interval.

\subsection{\texorpdfstring{lp langevin integral stability}{lp langevin integral stability}}


Specialize Banach-space stability to the actual lifted gradient on \(L^2(P;K)\). Its Lipschitz theorem supplies the required constant, so \(\|Y(t)-Z(t)\|_2\le e^{Lt}\|Y(0)-Z(0)\|_2\). This compares solutions with a common \(L^2\)-noise input and possibly different initial random variables.

\subsection{\texorpdfstring{lp langevin integral unique}{lp langevin integral unique}}


Apply the same-input uniqueness theorem in \(L^2(P;K)\). Equality here is equality of \(L^2\) classes for each time, hence almost-everywhere equality of representatives at each fixed time. The conclusion does not by itself select one exceptional null set valid simultaneously at every real time.

\section{\texorpdfstring{Identify chosen reference segments with one global solution}{Identify chosen reference segments with one global solution}}
\label{proof-module-UnadjustedLangevin-BrownianSegmentsGlobal-lean}
{\small\color{gray}Source module: \texttt{Unadjusted\allowbreak{}Langevin/\allowbreak{}Brownian\allowbreak{}Segments\allowbreak{}Global.\allowbreak{}lean}.\par}


The segment construction is compatible with the unique global integral solution. This lets stationarity of a global Brownian Langevin solution prove stationarity of every chosen reference segment.

\subsection{\texorpdfstring{brownian\allowbreak{}Langevin\allowbreak{}Noise increment}{brownianLangevinNoise increment}}


Expand the two increments from zero and subtract them. Their common \([B_0]\) term cancels, leaving \(\sqrt2([B_{a+t}]-[B_a])=N_a(t)\). This is an algebraic cocycle identity in \(L^2\), valid for the real-time extension as defined.

\subsection{\texorpdfstring{lp brownian integral shift}{lp brownian integral shift}}


Subtract the global integral equations at \(a+t\) and \(a\). Integral additivity gives the drift integral over \([a,a+t]\), and time translation turns it into \(\int_0^t\nabla U(Y_{a+s})ds\). The noise cocycle gives \(N_a(t)\), yielding the local segment equation with initial state \(Y_a\).

\subsection{\texorpdfstring{brownian\allowbreak{}Segment\allowbreak{}From eq global}{brownianSegmentFrom eq global}}


The shifted global path and the selected local segment are continuous solutions of the same \(L^2\) integral equation with the same initial value and noise input. Apply \(L^2\) uniqueness to identify them on \([0,h]\). The horizon premise ensures the shifted interval lies inside the global solution interval.

\subsection{\texorpdfstring{brownian\allowbreak{}Knots eq global}{brownianKnots eq global}}


Induct on the knot index. At zero, the global equation and zero-starting noise give the prescribed initial state. If the current knot equals \(Y_{nh}\), segment/\allowbreak{}global uniqueness makes its endpoint equal \(Y_{(n+1)h}\), establishing the next knot identity.

\subsection{\texorpdfstring{brownian\allowbreak{}Segments eq global}{brownianSegments eq global}}


The knot/\allowbreak{}global theorem supplies the matching initial value for the \(n\)-th segment. Apply \texttt{brownian\allowbreak{}Segment\allowbreak{}From\_\allowbreak{}eq\_\allowbreak{}global} to conclude \(Y^{(n)}_t=Y_{nh+t}\) on \([0,h]\). This is a deterministic identification of the constructions, with no stationarity premise.

\section{\texorpdfstring{Dalalyan’s theorem for the actual Gaussian ULA transition}{Dalalyan’s theorem for the actual Gaussian ULA transition}}
\label{proof-module-UnadjustedLangevin-CanonicalULA-lean}
{\small\color{gray}Source module: \texttt{Unadjusted\allowbreak{}Langevin/\allowbreak{}Canonical\allowbreak{}ULA.\allowbreak{}lean}.\par}


Let \(U:\mathbb R^d\to\mathbb R\) be differentiable and \(m\)-strongly convex with \(m>0\), with \(L\)-Lipschitz gradient. The canonical Brownian source derives all stochastic reference obligations and proves the Wasserstein bound for the iterated Gaussian transition \(P_h\), with finite initial \(W_2\).

\subsection{\texorpdfstring{Identify the canonical output law}{Identify the canonical output law}}


The Brownian source provides jointly independent standard Gaussian innovations; the product extension makes initialization independent of the entire sequence. The realization/\allowbreak{}Markov-law theorem identifies its output with \(P_h^k\mu\). The product output-law identity removes the chosen finite coupling from this law equality.

\subsection{\texorpdfstring{product\allowbreak{}Reference of independent brownian}{productReference of independent brownian}}


On each finite-coupling extension, pull back the independent Brownian coordinates and use the product independence theorems for the reference initialization. Construct a global \(L^2\) integral solution over the required finite horizon.



The proved MALA/\allowbreak{}Euler stationarity theorem gives Gibbs law for that global solution at every needed absolute time. Segment/\allowbreak{}global uniqueness identifies each recursively constructed reference segment with that solution. Thus the segment-stationarity premise of \texttt{product\allowbreak{}Reference\_\allowbreak{}of\_\allowbreak{}constructed\_\allowbreak{}brownian} is discharged, filling every reference field from the stated source and potential assumptions.

\subsection{\texorpdfstring{canonical\allowbreak{}Euler\allowbreak{}Noise}{canonicalEulerNoise}}


Specialize the source Euler noise to the explicitly constructed independent coordinate Brownian source. This abbreviation fixes one concrete implementation of the standard Gaussian innovations for the unconditional law theorem.

\subsection{\texorpdfstring{canonical\allowbreak{}Product\allowbreak{}Reference}{canonicalProductReference}}


The canonical Brownian construction supplies the coordinate Brownian laws and independence of their grid paths for every positive step. Apply \texttt{product\allowbreak{}Reference\_\allowbreak{}of\_\allowbreak{}independent\_\allowbreak{}brownian}. No stochastic reference structure or stationary-law assertion remains an external premise.

\subsection{\texorpdfstring{canonical\allowbreak{}Output eq ula\allowbreak{}Law\allowbreak{}Step}{canonicalOutput eq ulaLawStep}}


Finite initial transport provides a finite-cost coupling, and Gibbs has a finite second moment. Apply the general Brownian output-law identity on that extension with the canonical independent Brownian source. It proves the canonical product implementation has exactly the ordinary ULA law \(P_h^k\mu\).

\subsection{\texorpdfstring{dalalyan output finite}{dalalyan output finite}}


The canonical product reference supplies a finite-cost output coupling for every iteration whenever initial transport is finite. Its output law is \(P_h^k\mu\). Substituting this law proves \(W_2(P_h^k\mu,\pi_U)<\infty\), legitimizing the real Wasserstein notation used in the final estimates.

\subsection{\texorpdfstring{dalalyan}{dalalyan}}


Construct \texttt{canonical\allowbreak{}Product\allowbreak{}Reference} from differentiability, \(m\)-strong convexity, and \(L\)-Lipschitz gradient. Nontriviality of the finite coordinate space proves \(m\le L\); the finite-coupling product theorem gives the small-step bound with bias \((33/20)(L/m)\sqrt{hd}\).



Finally replace the canonical product output by \(P_h^k\mu\) using \texttt{canonical\allowbreak{}Output\_\allowbreak{}eq\_\allowbreak{}ula\allowbreak{}Law\allowbreak{}Step}. The resulting inequality is about the actual Gaussian ULA laws, and output finiteness has already been proved. The independent Brownian construction and stationary reference are internal parts of this proof.

\subsection{\texorpdfstring{dalalyan large step}{dalalyan large step}}


Use the same fully constructed canonical reference in the large-step product theorem. The step assumptions give \(0\le Lh-1<1\) and \(2-Lh>0\), so the bias is \((33/20)Lh\sqrt{hd}/(2-Lh)\). Rewrite its output law as the actual iterated Gaussian ULA law.


\end{document}
